% !TEX program = xelatex % (can also use "lualatex")
\documentclass[12pt,a4paper]{report}
\usepackage[utf8]{inputenc}
\usepackage[vietnamese]{babel} 
\usepackage{graphicx}
\usepackage{amsmath}
\usepackage{amssymb}
\usepackage{hyperref}
\usepackage{geometry}
\usepackage{fancyhdr}
\usepackage{tabularx} 
\usepackage{titlesec}
\usepackage{booktabs}
\usepackage{xcolor}
\usepackage{colortbl}
\usepackage{multirow}
\usepackage{makecell}
\usepackage{needspace}
\usepackage{etoolbox}
\usepackage{tikz}
\usepackage{tocloft}
\usepackage{float}
\usepackage{enumitem}
\usepackage{longtable}
\usepackage{cite}
\usepackage{listings}
\usepackage{subcaption}
\usepackage[nottoc]{tocbibind}
\usepackage{polyglossia}
\usepackage{algorithm}
\usepackage{algorithmic}
\usepackage{caption}
\usepackage{minted}
\usemintedstyle{friendly}
\usepackage{amsthm} 
\usetikzlibrary{arrows.meta, positioning, calc}

% ===== Custom Colors =====

\definecolor{hlblue}{RGB}{0,102,204}
\definecolor{hlgreen}{RGB}{0,153,76}

\definecolor{stepblue}{RGB}{52,152,219}
\definecolor{stepgreen}{RGB}{46,204,113}

% Optional extra colors
\definecolor{softblue}{RGB}{230,240,255}
\definecolor{softgreen}{RGB}{230,255,240}
\definecolor{darkblue}{RGB}{0,51,102}
\definecolor{hlyellow}{RGB}{255,204,0}
\definecolor{steporange}{RGB}{255,140,0}
\usetikzlibrary{calc}
\geometry{
    a4paper,
    left=2.5cm,
    right=2.5cm,
    top=2.5cm,
    bottom=3cm,     
    footskip=1.5cm  
}
% Định nghĩa màu sắc cho code và bảng biểu
\definecolor{excellent}{RGB}{0,128,0}
\definecolor{good}{RGB}{34,139,34}
\definecolor{medium}{RGB}{255,140,0}
\definecolor{poor}{RGB}{220,20,60}
\definecolor{codegreen}{rgb}{0,0.6,0}
\definecolor{codegray}{rgb}{0.5,0.5,0.5}
\definecolor{codepurple}{rgb}{0.58,0,0.82}
\definecolor{backcolour}{rgb}{0.95,0.95,0.92}
\definecolor{borderblue}{RGB}{16,16,144}

% Cấu hình hiển thị code Python
\lstdefinestyle{python}{
  language=Python,
  basicstyle=\ttfamily\footnotesize,
  keywordstyle=\bfseries\color{blue},
  commentstyle=\itshape\color{gray},
  stringstyle=\color{orange!70},
  numbers=left,
  numberstyle=\tiny\color{gray},
  backgroundcolor=\color{lightgray!20},
  breaklines=true,
  frame=single,
  showspaces=false,
  showstringspaces=false
}

% Cấu hình lề trang
\geometry{left=2.5cm, right=2.5cm, top=2.5cm, bottom=2.5cm}

% ---- Cấu hình hiển thị code Julia/Python ----
\lstdefinestyle{julia}{
  language=Python, % dùng Python làm base cho màu sắc
  basicstyle=\ttfamily\footnotesize,
  keywordstyle=\bfseries\color{blue},
  commentstyle=\itshape\color{gray},
  stringstyle=\color{orange!70},
  numbers=left,
  numberstyle=\tiny\color{gray},
  backgroundcolor=\color{lightgray!20},
  breaklines=true,
  frame=single,
  showspaces=false,
  showstringspaces=false,
  morekeywords={function,end,struct,mutable,const,using,import,export,
                module,where,begin,do,for,while,if,else,elseif,return,
                true,false,nothing,in,isa,typeof,push!,Dict,Vector,
                BitVector,@inbounds,@simd}
}
 

% Cấu hình Header/Footer
\pagestyle{fancy}
\fancyhf{}
\renewcommand{\headrulewidth}{0.5pt}
\renewcommand{\footrulewidth}{0.5pt}
\fancyhead[C]{\textit{Khai thác dữ liệu và ứng dụng - Nhóm 9}}
\fancyfoot[C]{\thepage}

% Cấu hình tiêu đề chương/mục
\titleformat{\chapter}[display]
{\centering\large\bfseries}{\chaptertitlename\ \thechapter}{10pt}{\large}
\titleformat{\section}
{\normalfont\large\bfseries}{\thesection}{0.3em}{}
\titleformat{\subsection}
{\normalfont\large\bfseries}{\thesubsection}{0.3em}{}


% Cấu hình mục lục
\renewcommand{\contentsname}{\begin{center} \Large \bfseries MỤC LỤC \end{center}}
\setlength{\cftbeforetoctitleskip}{20pt}
\setlength{\cftaftertoctitleskip}{10pt} 
\setlength{\cftparskip}{5pt} 
\renewcommand{\numberline}[1]{\bfseries#1\hspace{0.5em}} 

\setcounter{secnumdepth}{3} % Cho phép đánh số đến mức subsubsection

% Định nghĩa các môi trường học thuật
\newtheorem{definition}{Định nghĩa}[chapter]
\newtheorem{proposition}{Mệnh đề}[chapter]
\newtheorem{theorem}{Định lý}[chapter]
\newtheorem{corollary}{Hệ quả}[chapter]
\newtheorem{lemma}{Bổ đề}[chapter]
\newtheorem{remark}{Nhận xét}[chapter]
 
% Đổi tên nhãn chứng minh sang tiếng Việt
\renewcommand{\proofname}{\textit{Chứng minh}}

\begin{document}

% --- TRANG BÌA ---
\thispagestyle{empty}
\begin{tikzpicture}[remember picture, overlay]
  \draw[borderblue, line width=1pt] 
    ([shift={(1cm,1cm)}]current page.south west) rectangle 
    ([shift={(-1cm,-1cm)}]current page.north east);
  \draw[borderblue, line width=1pt] 
    ([shift={(1.1cm,1.1cm)}]current page.south west) rectangle 
    ([shift={(-1.1cm,-1.1cm)}]current page.north east);
\end{tikzpicture}

\begin{center}
    \vspace*{0cm}
    {\large BỘ GIÁO DỤC VÀ ĐÀO TẠO \par}
    \vspace{0.3cm}
    {\large TRƯỜNG ĐẠI HỌC KHOA HỌC TỰ NHIÊN \par}
    \vspace{0.3cm}
    {\large KHOA CÔNG NGHỆ THÔNG TIN \par}
    \vspace{0.3cm}
    {\large\textbf{\textbf{KHAI THÁC DỮ LIỆU VÀ ỨNG DỤNG}} \par}
    \vspace{0.3cm}
    
    \begin{center}
        \includegraphics[width=6cm]{LogoTruong.png}

    \end{center}
    \vspace{0.3cm}
    
    {\large\textbf{BÁO CÁO MÔN HỌC} \par}
    \vspace{0.3cm}
    {\large\textbf{PROJECT 02 - KHAI THÁC TẬP PHỔ BIẾN} \par}
    \vspace{0.5cm}
    
    \begin{tabular}{lll}
        \textbf{Giảng viên hướng dẫn} & : &GS. Lê Hoài Bắc \\
        & & ThS. Lê Nhựt Nam\\ & & ThS. Nguyễn Ngọc Đức \\ & & ThS. Nguyễn Thị Thu Hằng \\
        & & \\
        \textbf{Nhóm thực hiện} & : & Nhóm 9 \\
        & & \\
        \textbf{Lớp} & : & 23TNT1 \\
        & & \\
        \textbf{Nhóm sinh viên thực hiện} & : & \\
        & & Bàng Mỹ Linh \hspace{1.5cm} | 23122009 \\
        & & Lại Nguyễn Hồng Thanh \hspace{0.5cm} | 23122018 \\
        & & Phan Huỳnh Châu Thịnh \hspace{0.5cm} | 23122019 \\
        & & Nguyễn Gia Bảo \hspace{1.1cm} | 23122015 \\
    \end{tabular}
    
    \vfill
    {\Large\textbf{Tháng 4 năm 2026} \par}
\end{center}
\clearpage

% --- MỤC LỤC & DANH MỤC ---
\renewcommand{\chaptername}{Chương}
\renewcommand{\appendixname}{Phụ lục}
\renewcommand{\contentsname}{Mục lục}
\renewcommand{\listfigurename}{Danh sách hình ảnh}
\renewcommand{\listtablename}{Danh sách bảng}
\renewcommand{\bibname}{Tài liệu tham khảo}

\tableofcontents
\clearpage
\listoftables
\clearpage
\listoffigures
\clearpage

% ==========================================
% THÊM DANH MỤC THUẬT NGỮ Ở ĐÂY
% ==========================================
\chapter*{DANH MỤC THUẬT NGỮ}
\addcontentsline{toc}{chapter}{Danh mục thuật ngữ}

\begin{table}[H]
    \centering
    % Không dùng gạch đứng theo yêu cầu của thầy, dùng toprule/midrule/bottomrule
    \begin{tabularx}{\textwidth}{llX}
    \toprule
    \textbf{Thuật ngữ tiếng Việt} & \textbf{Thuật ngữ tiếng Anh (Gốc)} & \textbf{Ký hiệu / Ý nghĩa trong thuật toán} \\
    \midrule
    Cơ sở dữ liệu giao dịch & Transaction database & $\mathcal{D}$ - Dữ liệu đầu vào, được rút gọn trong bước tiền xử lý.\\
    \addlinespace
    Giao dịch & Transaction & $T$ - Một tập con các hạng mục. \\
    \addlinespace
    Tập hạng mục & Item set & $X, Y$ \\
    \addlinespace
    Độ hỗ trợ tuyệt đối & Support counter / Frequencies & $sup(X)$ - Được lưu tại biến đếm của mỗi phần tử trong mảng danh sách liên kết. \\
    \addlinespace
    Ngưỡng hỗ trợ tối thiểu & Minimum support & $minsup$ - Điều kiện cốt lõi để cắt tỉa item và gọi đệ quy. \\
    \addlinespace
    Tập phổ biến & Frequent item set & $F$ - Kết quả đầu ra duy nhất của thuật toán. \\
    \addlinespace
    Tiền tố & Prefix & $P$ - Tập hợp các item đã xét trong các vòng đệ quy trước đó. \\
    \addlinespace
    Phân bổ lại & Reassign & Hành động cắt bỏ item đầu tiên và chuyển giao dịch sang danh sách khác nằm bên phải. \\
    \addlinespace
    Mảng danh sách liên kết đơn & Array of singly linked lists & Cấu trúc dữ liệu tối giản được Relim sử dụng thay cho FP-tree. \\
    \bottomrule
    \end{tabularx}
\end{table}
\clearpage
% ==========================================
% --- NỘI DUNG CHÍNH ---
% CHƯƠNG 1: NỀN TẢNG LÝ THUYẾT
% CHƯƠNG 2: VÍ DỤ MINH HỌA TAY
% CHƯƠNG 3: CÀI ĐẶT
% CHƯƠNG 4: THỰC NGHIỆM VÀ ĐÁNH GIÁ
% \input{chapter1.tex}

% \input{chapter2.tex}
% \input{chapter3.tex}
% \input{chapter4.tex}

\chapter{NỀN TẢNG LÝ THUYẾT}
 
%% Người thực hiện: Châu Thịnh (3.1.1 & 3.1.2 b) e)) + Mỹ Linh (3.1.2 a) c) d))
 
% ------------------------------------------------------------
% \section{Bài toán Frequent Itemset Mining}
% % ------------------------------------------------------------
 
% %% Trình bày đầy đủ và chính xác các định nghĩa bằng ký hiệu toán học (người làm: Châu Thịnh)
 
% \subsection{Cơ sở dữ liệu giao dịch (Transaction Database)}
% %% Định nghĩa hình thức: D = {T1, T2, ..., Tn}, mỗi Ti ⊆ I với I là tập tất cả các item
% %% Ví dụ minh họa: bảng CSDL nhỏ (giống bảng Table 1 trong paper Borgelt)
 
% \subsection{Độ hỗ trợ (Support)}
% %% Định nghĩa sup(X) cho itemset X ⊆ I
% %% Phân biệt support tuyệt đối (số giao dịch chứa X) và support tương đối (tỉ lệ %)
% %% Công thức toán học chuẩn
 
% \subsection{Tập phổ biến (Frequent Itemset)}
% %% Định nghĩa: X là frequent khi sup(X) >= minsup
% %% Ghi chú về ý nghĩa minsup (ngưỡng người dùng chỉ định)
 
% \subsection{Tập đóng và Tập tối đại}
% %% Tập đóng (Closed Itemset): X đóng nếu không có superset Y ⊋ X nào có cùng support
% %% Tập tối đại (Maximal Itemset): X tối đại nếu không có superset Y ⊋ X nào là frequent
% %% Quan hệ: frequent ⊇ closed ⊇ maximal (vẽ sơ đồ Venn nếu có thể)
 
% \subsection{Tính chất Apriori (Downward Closure)}
% %% Phát biểu: nếu X không frequent thì mọi superset của X cũng không frequent
% %% Chứng minh ngắn gọn: dùng định nghĩa support và tập con
% %% Ý nghĩa: cho phép cắt tỉa không gian tìm kiếm

% ------------------------------------------------------------
\section{Bài toán Frequent Itemset Mining}
 
Phần này trình bày nền tảng lý thuyết hình thức của bài toán khai phá tập phổ biến
(Frequent Itemset Mining -- FIM), bao gồm các định nghĩa cơ bản, phân loại tập hạng
mục và tính chất giải tích quan trọng làm cơ sở cho mọi thuật toán FIM hiện đại.
 
% ----------------------------------------------------------
\subsection{Cơ sở dữ liệu giao dịch}
% ----------------------------------------------------------
 
\begin{definition}[Cơ sở dữ liệu giao dịch]
Cho $\mathcal{I} = \{i_1, i_2, \dots, i_m\}$ là tập hữu hạn các \textit{hạng mục}
(items) phân biệt. Một \textit{giao dịch} (transaction) $T$ là một bộ
$T = (\mathrm{tid}, X_T)$, trong đó $\mathrm{tid}$ là định danh duy nhất của giao
dịch và $X_T \subseteq \mathcal{I}$ là tập các hạng mục mà giao dịch đó chứa.
\textit{Cơ sở dữ liệu giao dịch} $\mathcal{D}$ là một đa tập hữu hạn gồm $n$
giao dịch:
\begin{equation}
    \mathcal{D} = \{T_1, T_2, \dots, T_n\}, \quad X_{T_k} \subseteq \mathcal{I},\;
    k = 1,\dots,n.
    \label{eq:database}
\end{equation}
Khi ngữ cảnh rõ ràng, ta đồng nhất giao dịch $T$ với tập hạng mục $X_T$ của nó
và viết $T \subseteq \mathcal{I}$.
\end{definition}
 
Bảng~\ref{tab:example-db} trình bày một ví dụ minh họa về cơ sở dữ liệu giao dịch
nhỏ sẽ được sử dụng xuyên suốt chương này.
 
\begin{table}[H]
    \centering
    \caption{Ví dụ về cơ sở dữ liệu giao dịch với $\mathcal{I} = \{1,2,3,4,5\}$}
    \label{tab:example-db}
    \begin{tabular}{cc}
        \toprule
        \textbf{TID} & \textbf{Giao dịch} \\
        \midrule
        $T_1$ & $\{1, 3, 4\}$ \\
        $T_2$ & $\{2, 3, 5\}$ \\
        $T_3$ & $\{1, 2, 3, 5\}$ \\
        $T_4$ & $\{2, 5\}$ \\
        \bottomrule
    \end{tabular}
\end{table}
 
% ----------------------------------------------------------
\subsection{Độ hỗ trợ}
% ----------------------------------------------------------
 
\begin{definition}[Độ hỗ trợ tuyệt đối]
Độ hỗ trợ tuyệt đối của một tập hạng mục $X \subseteq \mathcal{I}$ trong cơ sở dữ
liệu $\mathcal{D}$, ký hiệu $\mathrm{sup}(X)$, là số lượng giao dịch $T$ trong $\mathcal{D}$
chứa $X$:
\begin{equation}
    \mathrm{sup}(X) = \bigl|\{T \in \mathcal{D} \mid X \subseteq T\}\bigr|.
    \label{eq:support}
\end{equation}
\end{definition}
 
\begin{definition}[Độ hỗ trợ tương đối]
Độ hỗ trợ tương đối của $X$, ký hiệu $\mathrm{rsup}(X)$, là tỉ lệ giao dịch $T$ chứa $X$:
\begin{equation}
    \mathrm{rsup}(X) = \frac{\mathrm{sup}(X)}{|\mathcal{D}|} \in [0, 1].
    \label{eq:relsupport}
\end{equation}
\end{definition}
 
Độ hỗ trợ tối thiểu $\sigma$ (min support) do người dùng chỉ định
có thể được biểu diễn theo một trong hai dạng: số nguyên tuyệt đối hoặc tỉ lệ thập
phân. Trong báo cáo này, ký hiệu $\sigma$ được dùng thống nhất theo dạng tuyệt đối.
 
% ----------------------------------------------------------
\subsection{Tập phổ biến}
% ----------------------------------------------------------
 
\begin{definition}[Tập phổ biến]
Một tập hạng mục $X \subseteq \mathcal{I}$ được gọi là \textit{tập phổ biến}
(frequent itemset) với ngưỡng $\sigma$ khi và chỉ khi:
\begin{equation}
    \mathrm{sup}(X) \geq \sigma.
    \label{eq:frequent}
\end{equation}
Tập hợp tất cả các tập phổ biến trong $\mathcal{D}$ được ký hiệu là
$\mathcal{F}(\mathcal{D}, \sigma)$, hay $\mathcal{F}$ khi $\mathcal{D}$ và $\sigma$
đã được xác định rõ từ ngữ cảnh.

\begin{equation}
    \mathcal{F} = \{X \mid X \subseteq I \text{ và } \mathrm{sup}(X) \geq \mathrm{minsup}\}.
    \label{eq:fequent_itemsets}
\end{equation}
\end{definition}

Bài toán khai phá tập phổ biến đặt ra yêu cầu tìm ra toàn bộ $\mathcal{F}$, tức là
liệt kê đầy đủ mọi tập hạng mục thỏa mãn ngưỡng $\sigma$ cùng với giá trị độ hỗ trợ
tương ứng~\cite{borgelt2005}.
 
% ----------------------------------------------------------
\subsection{Tập đóng và tập tối đại}
% ----------------------------------------------------------
 
Trong thực tế, tập $\mathcal{F}$ có thể rất lớn; để thu hẹp biểu diễn, người ta định
nghĩa hai lớp con quan trọng sau.
 
\begin{definition}[Tập phổ biến đóng]

Một tập phổ biến $X \in \mathcal{F}$ được gọi là \textit{tập phổ biến đóng}
(\textit{closed frequent itemset}) nếu không tồn tại tập nào bao nó
$Y \supsetneq X$ có cùng độ hỗ trợ:
\begin{equation}
    X \text{ là tập đóng}
    \iff
    \nexists\, Y \supsetneq X :
    \mathrm{sup}(Y) = \mathrm{sup}(X).
    \label{eq:closed}
\end{equation}

Tập tất cả các tập phổ biến đóng được ký hiệu:
\begin{equation}
\mathcal{C}
=
\{X \mid X \in \mathcal{F}
\text{ và }
\nexists\, Y \supsetneq X
\text{ sao cho }
\mathrm{sup}(X)=\mathrm{sup}(Y)\}.
\label{eq:closed_itemset}
\end{equation}
\end{definition}
 
\begin{definition}[Tập tối đại]
Một tập phổ biến $X \in F$ được gọi là \textit{tập phổ biến tối đại}
(\textit{maximal frequent itemset}) nếu không tồn tại tập cha nghiêm ngặt
$Y \supsetneq X$ nào cũng là tập phổ biến:
\begin{equation}
    X \text{ là tập tối đại}
    \iff
    \nexists\, Y \supsetneq X :
    Y \in F.
    \label{eq:maximal}
\end{equation}

Gọi $M$ là tập tất cả các tập phổ biến tối đại:
\begin{equation}
M =
\left\{
X \;\middle|\;
X \in F
\text{ và }
\nexists\, Y \supsetneq X
\text{ mà }
Y \in F
\right\}.
\label{eq:maximal_itemset}
\end{equation}
\end{definition}
 
\begin{proposition}[Quan hệ phân cấp]
Ký hiệu $\mathcal{F}_{\mathrm{max}}$ và $\mathcal{F}_{\mathrm{cl}}$ lần lượt là tập
hợp các tập tối đại và tập đóng trong $\mathcal{D}$. Khi đó:
\begin{equation}
    \mathcal{F}_{\mathrm{max}} \subseteq \mathcal{F}_{\mathrm{cl}} \subseteq
    \mathcal{F}.
    \label{eq:hierarchy}
\end{equation}
\end{proposition}
 
Tập đóng và tập tối đại mang lại biểu diễn gọn hơn cho $\mathcal{F}$: từ
$\mathcal{F}_{\mathrm{cl}}$ có thể phục hồi đầy đủ độ hỗ trợ của mọi tập phổ biến,
còn $\mathcal{F}_{\mathrm{max}}$ chỉ cho biết các tập phổ biến lớn nhất (theo quan
hệ bao hàm) mà không lưu lại thông tin về độ hỗ trợ của các tập con.

\begin{figure}[H]
\centering
\begin{tikzpicture}[scale=1]
  \draw[fill=softblue, opacity=0.5, draw=borderblue, thick] (0,0) ellipse (4.5cm and 2.5cm);
  \node at (0, 1.8) {\textbf{Tập phổ biến (Frequent)}};
  
  \draw[fill=softgreen, opacity=0.6, draw=stepgreen, thick] (0,-0.3) ellipse (3cm and 1.7cm);
  \node at (0, 0.8) {\textbf{Tập đóng (Closed)}};
  
  \draw[fill=hlyellow, opacity=0.7, draw=orange, thick] (0,-0.8) ellipse (1.6cm and 0.9cm);
  \node at (0, -0.8) {\textbf{Tối đại}};
\end{tikzpicture}
\caption{Sơ đồ Venn thể hiện quan hệ bao hàm giữa Tập phổ biến, Tập đóng và Tập tối đại.}
\label{fig:venn-fcm}
\end{figure}
 
% ----------------------------------------------------------
\subsection{Tính chất Apriori}
% ----------------------------------------------------------
 
Tính chất sau đây là nền tảng lý thuyết cho phép cắt tỉa không gian tìm kiếm trong
hầu hết các thuật toán FIM.
 
\begin{proposition}[Tính đơn điệu giảm của độ hỗ trợ -- Anti-monotonicity]
\label{prop:antimonotone}
Với mọi $Y \subseteq X \subseteq \mathcal{I}$, ta có:
\begin{equation}
    \mathrm{sup}(X) \leq \mathrm{sup}(Y).
    \label{eq:apriori}
\end{equation}
Từ đó suy ra hai hệ quả quan trọng:
\begin{itemize}
    \item Nếu tập $X$ là phổ biến thì mọi tập con $Y \subseteq X$ cũng phổ biến.
    \item Nếu tập $X$ không phổ biến thì mọi tập cha $Z \supseteq X$ cũng không
    phổ biến.
\end{itemize}
\end{proposition}
 
\begin{proof}
Xét bất kỳ giao dịch $T \in \mathcal{D}$. Vì $Y \subseteq X$, điều kiện $X \subseteq T$
kéo theo $Y \subseteq T$. Do đó:
\begin{equation}
    \{T \in \mathcal{D} \mid X \subseteq T\} \subseteq \{T \in \mathcal{D} \mid
    Y \subseteq T\}.
    \label{eq:subset_proof}
\end{equation}
Số phần tử vế trái không vượt quá số phần tử vế phải; áp dụng
Định nghĩa~\eqref{eq:support}, ta thu được
$\mathrm{sup}(X) \leq \mathrm{sup}(Y)$, điều phải chứng minh.
\end{proof}
 
Tính chất này cho phép bất kỳ thuật toán FIM nào loại bỏ toàn bộ các tập cha của
một tập không phổ biến mà không cần liệt kê, từ đó giảm đáng kể số lượng
ứng viên cần xét.
 
% ------------------------------------------------------------
\section{Phân tích thuật toán RELIM}
% ------------------------------------------------------------
 
%% Tài liệu gốc: Borgelt, C. "Keeping Things Simple: Finding Frequent Item Sets
%% by Recursive Elimination." OSDM 2005.
%% Link paper: xem relim.pdf đã được cung cấp
 
%% (Người làm: Mỹ Linh)
%% Mô tả bằng lời trực quan:
%%   - RELIM lấy cảm hứng từ FP-Growth và H-Mine nhưng KHÔNG dùng prefix tree hay H-struct
%%   - Cách tiếp cận: đệ quy loại trừ từng item (recursive elimination)
%%   - Preprocessing: loại item không phổ biến, sắp xếp item theo tần suất tăng dần
%%   - Ý tưởng chính: ở mỗi bước, lấy danh sách giao dịch chứa item ít phổ biến nhất,
%%     đệ quy tìm itemset phổ biến có item đó làm tiền tố, sau đó xóa nó khỏi toàn bộ DB
%%   - Điểm khác biệt so với Apriori: không sinh candidate, duyệt theo chiều sâu
%%   - Điểm khác biệt so với FP-Growth: không cần xây dựng FP-tree phức tạp

\subsection{Ý tưởng cốt lõi}
\label{subsec:relim-idea}

Thuật toán Relim (\textit{Recursive Elimination}) được Borgelt đề xuất năm
2005~\cite{borgelt2005} với triết lý giữ mọi thứ đơn giản (\textit{Keeping
Things Simple}). Relim kế thừa cơ chế đệ quy của FP-Growth và
H-Mine, nhưng loại bỏ hoàn toàn các cấu trúc dữ liệu phức tạp như
cây tiền tố (FP-tree) hay cấu trúc siêu liên kết (H-struct): toàn bộ
thuật toán có thể được cài đặt trong một hàm đệ quy duy nhất.

Ý tưởng cốt lõi của Relim có thể hiểu trực giác như sau.
Hình dung không gian tất cả các tập phổ biến cần tìm như một cây con
liệt kê tập con. Thay vì xây dựng cây này theo chiều rộng và sinh hàng loạt
ứng viên như Apriori, Relim đặt câu hỏi đơn giản hơn:
\begin{quote}
\textit{Với mỗi hạng mục $i$, có những tập phổ biến nào chứa $i$? Nếu trả lời
được câu hỏi này cho từng $i$, ta đã liệt kê được toàn bộ tập phổ biến.}
\end{quote}
Để tránh đếm trùng (một tập có thể chứa nhiều hạng mục), Relim áp dụng nguyên
lý \textit{chia theo hạng mục đầu tiên}: cố định một thứ tự duyệt trên
các hạng mục (theo tần suất \textit{tăng dần}); với mỗi hạng mục $i$, chỉ tìm
các tập phổ biến mà $i$ là hạng mục \textit{đầu tiên} (theo thứ tự đã cố định)
xuất hiện trong tập đó. Khi đã xử lý xong tất cả tập có $i$ ở đầu, hạng mục $i$
được \textit{loại trừ} hoàn toàn khỏi cơ sở dữ liệu; thuật toán chuyển sang
hạng mục kế tiếp. Đây chính là cơ chế ``\textit{recursive elimination}'' đặt
tên cho thuật toán.

Tính đúng đắn của Relim dựa trên hai quan sát đơn giản:
\begin{itemize}
    \item[(i)] Mọi tập phổ biến $X \neq \emptyset$ đều có một hạng mục đầu tiên
    duy nhất theo thứ tự đã chọn. Do đó, nếu lần lượt liệt kê các tập theo từng
    hạng mục đầu tiên, ta liệt kê được mỗi tập \textit{đúng một lần}, không trùng
    không thiếu.
    \item[(ii)] Mọi tập phổ biến chứa $i$ làm hạng mục đầu tiên đều có dạng
    $\{i\} \cup Y$ với $Y$ là một tập phổ biến trong \textit{cơ sở dữ liệu điều
    kiện} của $i$ - tức tập các giao dịch chứa $i$, sau khi đã loại $i$ và các
    hạng mục đứng trước $i$ theo thứ tự duyệt. Việc tìm các $Y$ này lại chính là
    bài toán FIM ban đầu trên một CSDL nhỏ hơn, nên có thể giải đệ quy bằng cùng
    một thủ tục.
\end{itemize}
Cộng với tính chất Apriori (Mệnh đề~\ref{prop:antimonotone}) để cắt tỉa, ta thu
được toàn bộ tập phổ biến mà không cần sinh ứng viên hay quét lại cơ sở dữ liệu
nhiều lần.

Thứ tự duyệt là một lựa chọn thiết kế tự do, nhưng có ảnh hưởng quyết định đến
hiệu năng. Borgelt phân tích~\cite{borgelt2005}: nếu duyệt theo tần suất
\textit{tăng dần}, hạng mục được xử lý sớm là hạng mục \textit{ít phổ biến nhất},
nên CSDL điều kiện của nó nhỏ; các CSDL điều kiện về sau cũng được thu nhỏ tích
luỹ vì hạng mục đã xử lý bị loại bỏ. Thực nghiệm cho thấy thứ tự tăng dần nhanh
hơn rõ rệt so với thứ tự ngẫu nhiên, và đặc biệt nhanh hơn thứ tự
\textit{giảm dần}~\cite[Section~2.1]{borgelt2005}.

So với Apriori~\cite{apriori1993}, thuật toán duyệt không gian tìm kiếm theo \textit{chiều rộng}
(level-wise), ở vòng lặp thứ $k$ sinh tập ứng viên kích thước $k$ từ các tập
phổ biến kích thước $k-1$, sau đó quét lại toàn bộ CSDL để đếm support. Hai
nhược điểm cố hữu: (i)~số lần quét CSDL bằng kích thước tập phổ biến dài nhất;
(ii)~số ứng viên có thể bùng nổ theo hàm mũ trước khi cắt tỉa. Relim
hoàn toàn loại bỏ pha sinh ứng viên: nó duyệt theo \textit{chiều sâu}
(depth-first), chỉ cần quét CSDL gốc \textit{một lần} để đếm tần suất ban đầu;
các bước tiếp theo thao tác trực tiếp trên các CSDL điều kiện.

FP-Growth~\cite{fpgrowth2000} cũng theo lối đệ quy không sinh ứng viên, nhưng
nén CSDL thành cây tiền tố FP-tree để chia sẻ phần đầu chung giữa các giao
dịch; với mỗi hạng mục, nó dựng lại một \textit{cây điều kiện} (conditional
FP-tree). H-Mine~\cite{hmine2001} thay cây bằng cấu trúc H-struct gồm các siêu
liên kết. Cả hai cấu trúc đều \textit{phải tái cấu trúc} mỗi khi đi sâu một
tầng đệ quy - đó là chi phí đáng kể trên các tầng đệ quy sâu.

\noindent
Relim giữ nguyên ý tưởng đệ quy này nhưng dùng \textit{cấu trúc dữ liệu tối
giản}: chỉ một \textbf{mảng các danh sách liên kết đơn} (xem
Mục~\ref{subsec:ds}), trong đó các giao dịch được \textit{phân loại theo
hạng mục đầu tiên}. Khi đi sâu, thuật toán chỉ \textit{sao chép con trỏ} sang
mảng mới chứ không tái dựng cấu trúc phức tạp. Cái giá phải trả là Relim không
hưởng được lợi ích chia sẻ tiền tố như FP-tree (nên chậm hơn FP-Growth trên dữ
liệu dày), nhưng đổi lại nhận được sự đơn giản về cài đặt, ít chi phí khởi tạo
ở mỗi tầng đệ quy, và tiêu thụ bộ nhớ thấp - đặc biệt phù hợp với dữ liệu
thưa~\cite{borgelt2005}.

\begin{figure}[H]
\centering
\begin{tikzpicture}[scale=0.8, every node/.style={scale=0.85}]
% FP-Tree part
\begin{scope}[shift={(0,0)}]
    \node[font=\bfseries, text=blue!80!black] at (1.5, 1) {FP-Tree (FP-Growth)};
    \node[draw, circle, inner sep=2pt] (root) at (1.5, 0) {Root};
    \node[draw, circle, inner sep=2pt] (a) at (0, -1.5) {a:4};
    \node[draw, circle, inner sep=2pt] (b) at (3, -1.5) {b:2};
    \node[draw, circle, inner sep=2pt] (c1) at (-1, -3) {c:2};
    \node[draw, circle, inner sep=2pt] (d1) at (1, -3) {d:2};
    \node[draw, circle, inner sep=2pt] (c2) at (3, -3) {c:2};
    
    \draw[->, thick] (root) -- (a);
    \draw[->, thick] (root) -- (b);
    \draw[->, thick] (a) -- (c1);
    \draw[->, thick] (a) -- (d1);
    \draw[->, thick] (b) -- (c2);
    
    % Header table
    \draw[thick, fill=gray!10] (-2.5, 0) rectangle (-1.5, -3);
    \node at (-2, -0.5) {a}; \draw[->, dashed, blue] (-1.5, -0.5) -- (a);
    \node at (-2, -1.5) {b}; \draw[->, dashed, blue] (-1.5, -1.5) -- (b);
    \node at (-2, -2.5) {c}; \draw[->, dashed, blue] (-1.5, -2.5) -- (c1);
    \draw[->, dashed, blue] (c1) to[bend right=20] (c2);
    \draw (-2.5,-1) -- (-1.5,-1); \draw (-2.5,-2) -- (-1.5,-2);
\end{scope}

% Divider
\draw[thick, dotted, draw=black!50] (4.5, 1) -- (4.5, -4);

% Relim part
\begin{scope}[shift={(6,0)}]
    \node[font=\bfseries, text=red!80!black] at (2.5, 1) {Mảng liên kết (RELIM)};
    % Array (wider to avoid text clipping)
    \draw[thick, fill=gray!10] (-0.4, 0) rectangle (2.4, -3);
    \node at (1, -0.5) {A[a].cnt=4}; \draw (-0.4,-1) -- (2.4,-1);
    \node at (1, -1.5) {A[b].cnt=2}; \draw (-0.4,-2) -- (2.4,-2);
    \node at (1, -2.5) {A[c].cnt=4};
    
    % Nodes
    \node[draw, rounded corners, fill=white] (n1) at (4, -0.5) {(c,d)};
    \node[draw, rounded corners, fill=white] (n2) at (6.5, -0.5) {(d)};
    \node[draw, rounded corners, fill=white] (n3) at (4, -1.5) {(c)};
    
    \draw[->, thick] (2.4, -0.5) -- (n1);
    \draw[->, thick] (n1) -- (n2);
    \draw[->, thick] (2.4, -1.5) -- (n3);
    \draw[->, thick] (2.4, -2.5) -- (3.4, -2.5) node[right] {$\varnothing$};
\end{scope}
\end{tikzpicture}
\caption{So sánh trực quan: FP-Growth dùng cây tiền tố với mạng lưới con trỏ phức tạp (trái), trong khi RELIM chỉ dùng mảng danh sách liên kết đơn tối giản (phải).}
\label{fig:fp-vs-relim}
\end{figure}
 
\subsection{Cấu trúc dữ liệu chính}
\label{subsec:ds}
 
Khác với các thuật toán đương thời như FP-Growth (sử dụng cây tiền tố FP-tree) hay
H-Mine (sử dụng H-struct), thuật toán Relim thực hiện toàn bộ quá trình khai phá
trực tiếp trên cơ sở dữ liệu giao dịch nhờ một cấu trúc tối giản:
\textbf{mảng các danh sách liên kết đơn} (array of singly linked lists)~\cite{borgelt2005}.
 
Mỗi giao dịch được lưu trữ dưới dạng một mảng số nguyên, trong đó mỗi phần tử là
định danh số nguyên của một hạng mục. Sau bước tiền xử lý, các hạng mục trong mỗi
giao dịch được sắp xếp theo thứ tự tần suất tăng dần (hạng mục ít phổ biến hơn đứng
trước).
 
Thuật toán duy trì một mảng một chiều $A$ có $k$ phần tử, với $k$ là số hạng mục
phổ biến. Mỗi phần tử $A[i]$ đại diện cho hạng mục thứ $i$ và lưu trữ hai trường:
(i)~một \textit{biến đếm} $A[i].\mathit{cnt}$ ghi nhận tổng số giao dịch chứa hạng
mục đó trong cơ sở dữ liệu hiện tại; và (ii)~một \textit{con trỏ đầu danh sách}
$A[i].\mathit{head}$ trỏ đến phần tử đầu tiên của danh sách liên kết tương ứng.
 
Mỗi mắt xích trong danh sách liên kết chứa hai trường: một con trỏ đến mắt xích kế
tiếp (\textit{successor pointer}) và một con trỏ vào phần nội dung còn lại của giao
dịch (\textit{transaction pointer}).
 
Khi một giao dịch được đưa vào danh sách $A[i]$ (dựa trên hạng mục đứng đầu là $i$),
hạng mục $i$ được ``ẩn đi'' bằng cách tăng con trỏ giao dịch lên một vị trí
(\textit{pointer increment}). Thông tin về $i$ không bị mất vì nó đã được ngầm định
bởi vị trí $i$ của danh sách trong mảng.
 
Ngoài ra, thuật toán áp dụng tối ưu hóa Counter-Only: nếu sau khi ẩn hạng mục
đứng đầu, giao dịch còn lại chỉ gồm đúng một hạng mục $j$, thuật toán không cấp
phát mắt xích danh sách mà chỉ tăng $A[j].\mathit{cnt}$ lên một đơn vị. Tối ưu
hóa này loại bỏ hoàn toàn việc cấp phát bộ nhớ động cho các giao dịch ngắn, giảm
đáng kể chi phí quản lý bộ nhớ khi đi sâu vào các tầng đệ quy.
 
Hình~\ref{fig:relim-structure} minh họa trạng thái khởi tạo của mảng danh sách cho
một tập dữ liệu ví dụ, trong đó các mũi tên liền biểu diễn con trỏ danh sách, còn
mũi tên đứt biểu diễn con trỏ đến nội dung giao dịch sau khi đã ẩn hạng mục đứng đầu.
 
\begin{figure}[H]
\centering
\begin{tikzpicture}[scale=0.95]

% ======================================================
% ĐỊNH NGHĨA PHONG CÁCH (STYLES)
% ======================================================
\tikzset{
    arrayhead/.style={
        draw=black, thick, fill=gray!20,
        minimum width=2.2cm, minimum height=0.65cm,
        font=\small\bfseries
    },
    arraybox/.style={
        draw=black, thick, fill=white,
        minimum width=2.2cm, minimum height=0.65cm,
        font=\small
    },
    mynode/.style={
        draw=black, thick, rounded corners=3pt, fill=gray!5,
        minimum width=1.5cm, minimum height=0.65cm,
        font=\small
    },
    transbox/.style={
        draw=black!70, thick, fill=white,
        minimum width=2.8cm, minimum height=0.65cm,
        font=\small\ttfamily
    },
    arr/.style={
        ->, >=Stealth, thick, draw=black
    },
    darr/.style={
        ->, >=Stealth, dashed, thick, draw=black!60
    }
}

% ======================================================
% TỌA ĐỘ CÁC CỘT MẢNG
% ======================================================
\def\xA{0}
\def\xB{2.8}
\def\xC{5.6}

% ======================================================
% KHỐI MẢNG A (ARRAY HEADERS & BOXES)
% ======================================================
\node[arrayhead] (hA) at (\xA,0.7) {Item $e$};
\node[arrayhead] (hB) at (\xB,0.7) {Item $a$};
\node[arrayhead] (hC) at (\xC,0.7) {Item $c$};

\node[arraybox] (cA) at (\xA,0) {cnt $=5$};
\node[arraybox] (cB) at (\xB,0) {cnt $=4$};
\node[arraybox] (cC) at (\xC,0) {cnt $=2$};

\node[arraybox] (pA) at (\xA,-0.7) {head};
\node[arraybox] (pB) at (\xB,-0.7) {head};
\node[arraybox] (pC) at (\xC,-0.7) {head};

% Khung viền bao quanh Mảng A (Vừa khít hơn)
\draw[thick] (-1.3, 1.15) rectangle (6.9, -1.15);
\node[font=\small\bfseries, anchor=east] at (-1.4, 0) {Mảng $A$};

% ======================================================
% DANH SÁCH LIÊN KẾT CỦA ITEM E
% ======================================================
\node[mynode] (e1) at (\xA, -2.8) {Node$_1$};
\node[mynode] (e2) at (\xA, -4.2) {Node$_2$};

% Đã TĂNG KHOẢNG CÁCH sang phải (từ +3.2 lên +4.0) để mũi tên dài ra
\node[transbox] (t1) at (\xA + 4.0, -2.8) {$(a,c,b,d)$};
\node[transbox] (t2) at (\xA + 4.0, -4.2) {$(c,d)$};

% Mũi tên chính
\draw[arr] (pA.south) -- (e1.north);
\draw[arr] (e1.south) -- (e2.north);

% Mũi tên đứt nét kèm text tự động căn giữa
\draw[darr] (e1.east) -- node[above, font=\footnotesize\itshape, text=black!80] {bỏ qua $e$} (t1.west);
\draw[darr] (e2.east) -- node[above, font=\footnotesize\itshape, text=black!80] {bỏ qua $e$} (t2.west);

% Ký hiệu rỗng bám thẳng vào chân e2
\draw[thick, draw=black!60] (e2.south) -- ++(0,-0.4) node[below, inner sep=1pt, text=black] {$\varnothing$};

% ======================================================
% DANH SÁCH LIÊN KẾT CỦA ITEM C
% ======================================================
\node[font=\small\itshape, text=black!50, anchor=north] (emptyC) at (\xC, -1.3) {(danh sách rỗng)};

% ======================================================
% DANH SÁCH LIÊN KẾT CỦA ITEM A
% ======================================================
% Đặt node của Item A ngang hàng với e1 để hình cân đối
\node[mynode] (a1) at (8.8, -2.8) {Node$_1$};
% Đã TĂNG KHOẢNG CÁCH sang phải (từ 11.9 lên 12.8) để đồng nhất độ dài mũi tên
\node[transbox] (aT) at (12.8, -2.8) {$(b,d)$};

% Đường nối vuông góc đi luồn xuống dưới mảng rồi rẽ vào a1
\draw[arr, rounded corners=4pt] (pB.south) -- (\xB, -1.8) -- (8.8, -1.8) -- (a1.north);

\draw[darr] (a1.east) -- node[above, font=\footnotesize\itshape, text=black!80] {bỏ qua $a$} (aT.west);

% Ký hiệu rỗng bám thẳng vào chân a1
\draw[thick, draw=black!60] (a1.south) -- ++(0,-0.4) node[below, inner sep=1pt, text=black] {$\varnothing$};

% ======================================================
% KHUNG CHÚ GIẢI TỰ ĐỘNG KÍCH THƯỚC (Auto-sizing Legend)
% ======================================================
% Đã dịch khung chú giải sang 6.3 để nằm giữa tổng thể hình mới đã rộng hơn
\node[draw=black!40, fill=gray!2, rounded corners=4pt, inner sep=8pt, font=\small, text=black] at (6.3, -5.8) {
    \textbf{Chú giải:} 
    \texttt{cnt} -- biến đếm hỗ trợ; \quad 
    \texttt{head} -- con trỏ đầu danh sách; \quad 
    \texttt{Node} -- mắt xích liên kết.
};

\end{tikzpicture}
\caption{Cấu trúc mảng danh sách liên kết đơn của thuật toán Relim (Đã tinh chỉnh khoảng cách).}
\label{fig:relim-structure}
\end{figure}
Nhờ không phải xây dựng lại cấu trúc dữ liệu mới (như tách nhánh trong FP-tree) sau
mỗi bước đệ quy, Relim tiết kiệm đáng kể tài nguyên bộ nhớ và chi phí khởi tạo,
đặc biệt ở các tầng đệ quy sâu khi kích thước cơ sở dữ liệu điều kiện đã thu hẹp
đáng kể~\cite{borgelt2005}.
 
 
% Mỹ Linh 
\subsection{Mã giả thuật toán}
\label{subsec:relim-pseudocode}

Phần này trình bày mã giả đầy đủ của Relim, bám sát mô tả trong Section~2 của
bài báo gốc~\cite{borgelt2005}. Để tiện trình bày, thuật toán được chia thành
ba thủ tục:
\begin{itemize}
    \item \textsc{Relim}: thủ tục chính, thực hiện tiền xử lý và khởi động đệ
    quy (tương ứng Mục~2.1 và mở đầu Mục~2.3 của bài báo gốc).
    \item \textsc{Recursive-Relim}: hàm đệ quy lõi -- duyệt mảng danh sách,
    xuất tập phổ biến, gọi đệ quy và phân bổ lại (Mục~2.3).
    \item \textsc{Build-Conditional-DB}: thủ tục xây dựng cơ sở dữ liệu điều
    kiện cho một hạng mục (đoạn ``second array of lists'' trong Mục~2.3).
\end{itemize}
Các ký hiệu chính dùng trong mã giả được liệt kê ở Bảng~\ref{tab:pseudocode-notation}
và nhất quán với Danh mục thuật ngữ ở đầu báo cáo.

\begin{table}[H]
  \centering
  \caption{Quy ước ký hiệu trong mã giả}
  \label{tab:pseudocode-notation}
  \begin{tabularx}{\textwidth}{lX}
    \toprule
    \textbf{Ký hiệu} & \textbf{Ý nghĩa} \\
    \midrule
    $A[i]$ & Phần tử thứ $i$ của mảng các danh sách liên kết \\
    $A[i].\mathit{cnt}$ & Biến đếm hỗ trợ (số giao dịch chứa hạng mục $i$ làm hạng mục dẫn đầu) \\
    $A[i].\mathit{head}$ & Con trỏ đầu danh sách liên kết \\
    $e.\mathit{next}$ & Con trỏ đến mắt xích kế tiếp trong danh sách \\
    $e.\mathit{trans}$ & Con trỏ vào phần còn lại của giao dịch (sau khi ẩn các hạng mục đứng trước) \\
    $P$ & Tiền tố -- tập hạng mục đã xét trong các vòng đệ quy trước đó \\
    $B$ & Mảng danh sách của cơ sở dữ liệu điều kiện (conditional database) \\
    $\sigma$ & Ngưỡng hỗ trợ tối thiểu (minsup) \\
    \bottomrule
  \end{tabularx}
\end{table}


\subsubsection*{Thủ tục chính \textsc{Relim}}


\begin{algorithm}[H]
\caption{\textsc{Relim}: thủ tục chính của thuật toán Relim}
\label{alg:relim-main}
\begin{algorithmic}[1]
\REQUIRE Cơ sở dữ liệu giao dịch $\mathcal{D}$; ngưỡng tối thiểu $\sigma$.
\ENSURE Tập $\mathcal{F}$ gồm mọi tập phổ biến của $\mathcal{D}$ kèm độ hỗ trợ.
\STATE $\mathcal{F} \leftarrow \emptyset$ \COMMENT{tập kết quả}
\STATE $\mathit{freq} \leftarrow \textsc{CountItemFrequencies}(\mathcal{D})$
       \COMMENT{một lần quét, $O(\sum_{T \in \mathcal{D}}|T|)$}
\STATE $\mathcal{I}^* \leftarrow \{i \in \mathcal{I} : \mathit{freq}[i] \geq \sigma\}$
       \COMMENT{loại hạng mục không phổ biến}
\STATE Sắp xếp $\mathcal{I}^*$ tăng dần theo $\mathit{freq}$, gán chỉ số $1,\dots,k$
\COMMENT{$k = |\mathcal{I}^*|$, định nghĩa thứ tự $\preceq$}
\FORALL{$T \in \mathcal{D}$}
    \STATE Loại các hạng mục $\notin \mathcal{I}^*$ khỏi $T$
    \STATE Sắp xếp các hạng mục còn lại của $T$ theo thứ tự $\preceq$
\ENDFOR
\STATE Khởi tạo mảng $A$ với $k$ phần tử rỗng, $A[i].\mathit{cnt} \leftarrow 0$,
       $A[i].\mathit{head} \leftarrow \textsc{Nil}$
\FORALL{$T \in \mathcal{D}$ với $T \neq \emptyset$}
    \STATE $i \leftarrow$ chỉ số hạng mục đầu tiên của $T$
    \STATE $A[i].\mathit{cnt} \leftarrow A[i].\mathit{cnt} + 1$
    \IF{$|T| > 1$}
        \STATE Cấp phát mắt xích mới $e$; $e.\mathit{trans} \leftarrow$ con trỏ vào
               vị trí thứ hai của $T$ \COMMENT{ẩn hạng mục đầu bằng dịch con trỏ}
        \STATE $e.\mathit{next} \leftarrow A[i].\mathit{head}$;
               $A[i].\mathit{head} \leftarrow e$
    \ENDIF
    \COMMENT{nếu $|T|=1$: chỉ tăng counter, không cấp phát mắt xích (tối ưu Counter-Only)}
\ENDFOR
\STATE $\textsc{Recursive-Relim}(A, k, \emptyset, \sigma, \mathcal{F})$
\RETURN $\mathcal{F}$
\end{algorithmic}
\end{algorithm}

\noindent
\textbf{Diễn giải.} Các dòng 2-4 thực hiện đúng một lần quét CSDL gốc để đếm
tần suất và xác định tập hạng mục phổ biến đơn lẻ; thứ tự tăng dần được cố định
ngay tại đây. Vòng lặp ở dòng 5-8 \textit{thanh lọc} từng giao dịch: loại hạng
mục không phổ biến và sắp xếp theo $\preceq$. Vòng lặp ở dòng 10--17 thực hiện
\textit{phân loại theo hạng mục đầu}: hạng mục đầu của mỗi giao dịch đã sắp xếp
quyết định danh sách $A[i]$ mà mắt xích được chèn vào; bản thân hạng mục đầu
được \textit{ẩn ngầm} qua vị trí trong mảng (do đó \texttt{e.trans} trỏ tới
hạng mục thứ hai). Trường hợp giao dịch chỉ còn một hạng mục sau khi ẩn (thân
giao dịch rỗng) được xử lý theo tối ưu Counter-Only của Borgelt: chỉ tăng biến
đếm, không cấp phát mắt xích, vì không có ``đuôi'' nào để đẩy sâu hơn.


\subsubsection*{Hàm đệ quy \textsc{Recursive-Relim}}


\begin{algorithm}[H]
\caption{\textsc{Recursive-Relim}: hàm đệ quy lõi}
\label{alg:relim-recursive}
\begin{algorithmic}[1]
\REQUIRE Mảng danh sách $A$ kích thước $k$; tiền tố $P$; ngưỡng $\sigma$;
         tập kết quả $\mathcal{F}$ (\textit{tham chiếu}).
\FOR{$i \leftarrow 1$ \TO $k$}
    \IF{$A[i].\mathit{cnt} \geq \sigma$}
        \STATE $\mathcal{F} \leftarrow \mathcal{F} \cup
               \{(P \cup \{i\}, \, A[i].\mathit{cnt})\}$
               \label{line:output}
               \COMMENT{xuất tập phổ biến $P \cup \{i\}$ với support}
        \STATE $B \leftarrow \textsc{Build-Conditional-DB}(A, i)$
               \label{line:buildcond}
               \COMMENT{CSDL điều kiện cho $i$}
        \STATE $\textsc{Recursive-Relim}(B, i-1, P \cup \{i\}, \sigma, \mathcal{F})$
               \label{line:recurse}
               \COMMENT{đệ quy với tiền tố mở rộng}
    \ENDIF
    \STATE $e \leftarrow A[i].\mathit{head}$
    \WHILE{$e \neq \textsc{Nil}$}
        \STATE $e_{\text{next}} \leftarrow e.\mathit{next}$
        \IF{$e.\mathit{trans}$ vẫn còn ít nhất một hạng mục}
            \STATE $j \leftarrow$ chỉ số hạng mục đầu hiện tại của $e.\mathit{trans}$
            \STATE Dịch $e.\mathit{trans}$ sang hạng mục kế tiếp \COMMENT{ẩn $j$}
            \IF{$e.\mathit{trans}$ chưa rỗng}
                \STATE $e.\mathit{next} \leftarrow A[j].\mathit{head}$;
                       $A[j].\mathit{head} \leftarrow e$
                       \COMMENT{phân bổ lại sang $A[j]$}
            \ENDIF
            \STATE $A[j].\mathit{cnt} \leftarrow A[j].\mathit{cnt} + 1$
        \ENDIF
        \STATE $e \leftarrow e_{\text{next}}$
    \ENDWHILE
    \STATE $A[i].\mathit{head} \leftarrow \textsc{Nil}$;\;
           $A[i].\mathit{cnt} \leftarrow 0$ \COMMENT{loại trừ $i$ khỏi $A$}
\ENDFOR
\end{algorithmic}
\end{algorithm}

\noindent
\textbf{Diễn giải.} Vòng \texttt{for} ngoài duyệt mảng $A$ từ trái sang phải;
biến chạy $i$ là chỉ số của hạng mục đang đóng vai trò ``hạng mục đầu của các
tập phổ biến cần tìm''. Khối điều kiện ở dòng 2-6 hiện thực hai quan sát then
chốt đã nêu ở Mục~\ref{subsec:relim-idea}:
\begin{itemize}
    \item \textit{Dòng 3 - xuất tập phổ biến.} Nếu $A[i].\mathit{cnt} \geq
    \sigma$ thì $\{i\}$ phổ biến trong CSDL hiện tại, kết hợp với tiền tố $P$
    cho tập phổ biến $P \cup \{i\}$ ở cấp đệ quy này (Mệnh đề~\ref{prop:antimonotone}
    đảm bảo nếu $A[i].\mathit{cnt} < \sigma$ thì mọi mở rộng đều không phổ
    biến --- nhánh đệ quy được cắt tỉa).
    \item \textit{Dòng 4-5 - gọi đệ quy.} CSDL điều kiện $B$ tập hợp các
    giao dịch chứa $i$, đã loại $i$ và các hạng mục đứng trước $i$. Mọi tập phổ
    biến của $B$ mở rộng tiền tố $P \cup \{i\}$ cho ra tập phổ biến mới.
\end{itemize}
Vòng \texttt{while} ở dòng 8-19 thực hiện \textit{phân bổ lại} (reassign): với
mỗi mắt xích $e$ trong $A[i]$, lấy hạng mục đầu tiên $j$ hiện tại của
$e.\mathit{trans}$ ($j > i$ vì danh sách đã sắp xếp tăng dần), ẩn $j$ đi bằng
một bước dịch con trỏ, rồi chèn $e$ vào danh sách $A[j]$. Lưu ý cập nhật
$A[j].\mathit{cnt}$ ngay cả khi mắt xích không được chèn (trường hợp
$e.\mathit{trans}$ trở thành rỗng sau khi ẩn $j$) - đây chính là tối ưu
Counter-Only ở cấp đệ quy này. Sau khi tất cả mắt xích đã rời $A[i]$, danh sách
này được làm rỗng (dòng 20): hạng mục $i$ được \textit{loại trừ vĩnh viễn} khỏi
CSDL hiện tại; mọi tập phổ biến tiếp theo sẽ có hạng mục đầu lớn hơn $i$.

\paragraph{Lưu ý kỹ thuật.}
Việc đệ quy ở dòng~\ref{line:recurse} chỉ truyền kích thước $i-1$ (thay vì
$k$): đây là tối ưu mà Borgelt~\cite[Section~2.4]{borgelt2005} đề xuất - trong
CSDL điều kiện $B$, các vị trí $i, i+1, \dots, k$ luôn rỗng (vì mọi giao dịch
trong $B$ ban đầu chứa $i$, và sau khi ẩn $i$, hạng mục đầu mới luôn nhỏ hơn
$i$ theo thứ tự $\preceq$). Việc này tránh các vòng lặp trống ở các cấp đệ quy
sâu.


\subsubsection*{Thủ tục \textsc{Build-Conditional-DB}}


\begin{algorithm}[H]
\caption{\textsc{Build-Conditional-DB}: dựng CSDL điều kiện cho hạng mục $i$}
\label{alg:relim-build-cond}
\begin{algorithmic}[1]
\REQUIRE Mảng danh sách $A$; chỉ số $i$ của hạng mục đang xét.
\ENSURE Mảng danh sách $B$ là CSDL điều kiện ứng với tiền tố $\{i\}$.
\STATE Khởi tạo mảng $B$ với $i-1$ phần tử,
       $B[j].\mathit{cnt} \leftarrow 0$,
       $B[j].\mathit{head} \leftarrow \textsc{Nil}$ với mọi $j < i$
\STATE $e \leftarrow A[i].\mathit{head}$
\WHILE{$e \neq \textsc{Nil}$}
    \IF{$e.\mathit{trans}$ vẫn còn ít nhất một hạng mục}
        \STATE $j \leftarrow$ chỉ số hạng mục đầu hiện tại của $e.\mathit{trans}$
        \STATE Tạo bản sao mắt xích $e' \leftarrow$ \textsc{Copy}$(e)$
        \STATE Dịch $e'.\mathit{trans}$ sang hạng mục kế tiếp \COMMENT{ẩn $j$}
        \IF{$e'.\mathit{trans}$ chưa rỗng}
            \STATE $e'.\mathit{next} \leftarrow B[j].\mathit{head}$;
                   $B[j].\mathit{head} \leftarrow e'$
        \ENDIF
        \STATE $B[j].\mathit{cnt} \leftarrow B[j].\mathit{cnt} + 1$
    \ENDIF
    \STATE $e \leftarrow e.\mathit{next}$
\ENDWHILE
\RETURN $B$
\end{algorithmic}
\end{algorithm}

\noindent
\textbf{Diễn giải.} Thủ tục này song hành với vòng phân bổ lại trong
\textsc{Recursive-Relim}: cùng một mắt xích vừa được \textit{phân bổ lại} (cập
nhật vào $A$) vừa được \textit{nhân bản} (chèn bản sao vào $B$). Điểm tinh tế
là:
\begin{itemize}
    \item Mắt xích trong $B$ là \textit{bản sao} của mắt xích trong $A$, không
    phải cùng một mắt xích. Cả hai bản sao đều trỏ vào \textit{cùng một mảng
    giao dịch} (qua trường \texttt{trans}); việc dịch con trỏ chỉ thay đổi vị
    trí đọc trong mảng đó, không sao chép nội dung giao dịch. Đây là chìa khoá
    giúp Relim tiết kiệm bộ nhớ.
    \item Trường hợp $e.\mathit{trans}$ chỉ còn đúng một hạng mục: sau khi ẩn
    $j$, phần đuôi rỗng, nên không cấp phát mắt xích mới trong $B[j]$; chỉ tăng
    $B[j].\mathit{cnt}$ (tối ưu Counter-Only).
    \item Tối ưu cài đặt được Borgelt nêu trong Mục~2.4: nếu
    $A[i].\mathit{cnt} < \sigma$, có thể bỏ qua hoàn toàn bước
    \textsc{Build-Conditional-DB} vì không thể tìm thấy tập phổ biến nào trong
    nhánh đệ quy đó (điều này đã được kiểm tra ở dòng~2 của
    Thuật toán~\ref{alg:relim-recursive} trước khi gọi tới đây).
\end{itemize}

\paragraph{Mối liên hệ với cấu trúc dữ liệu.}
Trên hai cấu trúc song song $A$ và $B$, mỗi vòng lặp ngoài của
\textsc{Recursive-Relim} đồng thời (i)~\textit{thu hẹp} $A$ bằng cách loại
trừ hạng mục $i$ và (ii)~\textit{nhân bản} một phần $A[i]$ sang $B$ để đệ quy
sâu hơn. Hai thao tác này tương ứng chính xác với hai cột trái và phải
trong Hình~1 của bài báo Borgelt: cột trái là sự ``tan rã'' dần của mảng $A$
qua bốn bước cho đến khi rỗng; cột phải là các mảng $B$ được dựng tạm với tiền
tố tương ứng (prefix~$e$, prefix~$a$, prefix~$c$, prefix~$b$).

\subsection{Phân tích độ phức tạp}
\label{subsec:relim-complexity}

Bài báo gốc~\cite{borgelt2005} chủ yếu đánh giá Relim qua thực nghiệm và không
trình bày phân tích độ phức tạp hình thức. Trong phần này, chúng tôi suy ra
các cận tiệm cận trực tiếp từ giả mã ở Mục~\ref{subsec:relim-pseudocode}. Quy
ước các tham số sau, dùng thống nhất cho cả phân tích thời gian và không gian:
\begin{itemize}
    \item $n = |\mathcal{D}|$: số giao dịch của CSDL gốc;
    \item $w = \frac{1}{n}\sum_{T \in \mathcal{D}}|T|$: chiều rộng giao dịch
    trung bình;
    \item $N = \sum_{T \in \mathcal{D}}|T| = n \cdot w$: tổng kích thước (số
    cặp \textit{giao dịch-hạng mục}) - đại lượng then chốt trong các phân
    tích sau;
    \item $m = |\mathcal{I}|$: số hạng mục phân biệt; $k = |\mathcal{I}^*|$: số
    hạng mục phổ biến đơn lẻ (sau cắt tỉa, $k \leq m$);
    \item $|\mathcal{F}|$: số tập phổ biến đầu ra;
    \item $\sigma$: ngưỡng support tối thiểu.
\end{itemize}


\subsubsection*{Độ phức tạp thời gian}

\paragraph{Chi phí tiền xử lý (cố định).}
Các bước trong Thuật toán~\ref{alg:relim-main}:
\begin{itemize}
    \item Đếm tần suất (dòng~2): một lần quét toàn bộ CSDL với $O(N)$ thao tác.
    \item Cắt tỉa và sắp xếp $\mathcal{I}^*$ (dòng~3-4): $O(m + m\log m)$.
    \item Sắp xếp các hạng mục trong mỗi giao dịch (dòng~5-8): mỗi giao dịch
    tốn $O(|T|\log|T|)$ nếu dùng so sánh; với $|T| \leq w_{\max}$, tổng chi phí
    tối đa $O(N \log w_{\max})$. Trên CSDL điển hình, đây không phải là điểm
    nghẽn.
    \item Khởi tạo mảng $A$ và phân loại giao dịch theo hạng mục đầu (dòng
    9-17): $O(n)$.
\end{itemize}
\noindent
Tổng chi phí tiền xử lý: $\boxed{O(N\log w_{\max} + m\log m)}$. Khi $w_{\max}$
hằng số, đơn giản còn $O(N + m\log m)$.

\paragraph{Chi phí một cấp đệ quy.}
Xét một lời gọi \textsc{Recursive-Relim}$(A, k', P, \sigma, \cdot)$ ở một cấp
đệ quy bất kỳ, với mảng $A$ kích thước $k'$ và tổng số mắt xích là $L$. Toàn bộ
công việc ở cấp này gồm hai phần:
\begin{enumerate}
    \item[(R1)] Vòng phân bổ lại (dòng~8-19): mỗi mắt xích được \textit{dịch
    con trỏ một lần} và chèn vào danh sách khác $O(1)$. Tổng chi phí cho cả
    mảng: $O(L + k')$.
    \item[(R2)] Các lời gọi \textsc{Build-Conditional-DB} cho các nhánh phổ
    biến (dòng~4): mỗi mắt xích $e \in A[i]$ tạo \textit{tối đa} một bản sao
    với chi phí $O(1)$. Tổng chi phí phần này, cộng dồn trên mọi $i$ có
    $A[i].\mathit{cnt} \geq \sigma$, không vượt quá $O(L)$.
\end{enumerate}
\noindent
Vậy chi phí cục bộ (không tính đệ quy con) tại mỗi cấp là $O(L + k')$.

\paragraph{Tổng chi phí qua mọi cấp đệ quy.}
Quan sát then chốt~\cite[Section~2]{borgelt2005}: tổng số mắt xích sinh ra trên
\textit{toàn bộ} cây đệ quy bằng tổng kích thước các CSDL điều kiện. Dùng định
nghĩa độ hỗ trợ, có thể viết:
\begin{equation}
    \mathcal{W}_{\text{lists}} \;=\; \sum_{X \in \mathcal{F}} \mathrm{sup}(X)
    \;\;\leq\;\; |\mathcal{F}| \cdot n.
    \label{eq:total-list-work}
\end{equation}
Trực giác: với mỗi tập phổ biến $X = \{i_1 \prec i_2 \prec \dots \prec i_\ell\}$
được xuất ra, mỗi giao dịch $T \supseteq X$ đóng góp đúng một mắt xích vào CSDL
điều kiện ở cấp đệ quy ứng với tiền tố $\{i_1,\dots,i_{\ell-1}\}$, danh sách
$A[i_\ell]$. Cộng dồn cho tất cả tập phổ biến cho cận
$\sum_{X} \mathrm{sup}(X)$. Mỗi mắt xích lại được xử lý trong $O(1)$ ở cấp đệ
quy của nó, nên tổng công việc xử lý mắt xích là
$O\!\left(\sum_X \mathrm{sup}(X)\right)$.

\noindent
Kèm theo chi phí duyệt mảng (mỗi cấp đệ quy có thêm $O(k')$ phép duyệt trống),
ta có ước lượng tổng quát:
\begin{equation}
    T_{\textsc{Relim}}(\mathcal{D}, \sigma) \;=\; O\!\Bigl(N \;+\; m\log m
    \;+\; \sum_{X \in \mathcal{F}} \mathrm{sup}(X)\Bigr).
    \label{eq:relim-overall-time}
\end{equation}

\paragraph{Trường hợp tốt nhất.}
Khi $\sigma$ rất cao, hầu hết hạng mục không phổ biến: $k \to 0$ và $|\mathcal{F}|$
nhỏ. Trong cực đoan $k=0$ (không có hạng mục nào phổ biến), thuật toán dừng
ngay sau khi đếm tần suất:
\begin{equation}
    T_{\text{best}} = O(N + m\log m).
\end{equation}
Ngay cả khi $k > 0$ nhưng vẫn nhỏ, vòng đệ quy thưa thớt và thuật toán xấp xỉ
tuyến tính theo $N$.

\paragraph{Trường hợp xấu nhất.}
Khi $\sigma = 1$ (mọi tập con khác rỗng đều phổ biến), số tập phổ biến đạt
$|\mathcal{F}| = 2^k - 1$ và $\sum_X \mathrm{sup}(X)$ có thể đạt
$O(n \cdot 2^k)$. Thay vào~\eqref{eq:relim-overall-time}:
\begin{equation}
    T_{\text{worst}} = O\!\bigl(n \cdot 2^k\bigr).
\end{equation}
Đây là cận xấu nhất chung cho \textit{mọi} thuật toán FIM khai phá đầy đủ tập
phổ biến: số đầu ra đã có thể là $2^k$. Sự bùng nổ này không phải do thiết kế
của Relim mà do kích thước bài toán.

\paragraph{Trường hợp trung bình.}
Khó đưa cận khép kín do phụ thuộc phân phối dữ liệu. Borgelt báo cáo thực
nghiệm trên năm bộ dữ liệu chuẩn (BMS-Webview-1, T10I4D100K, census, chess,
mushroom) cho thấy~\cite[Section~3]{borgelt2005}:
\begin{itemize}
    \item Đường cong $\log_{10}(\text{thời gian})$ của Relim luôn nhanh hơn
    Apriori, đôi khi nhanh hơn Eclat;
    \item FP-Growth thường nhanh hơn Relim trên dữ liệu dày
    (\texttt{chess}, \texttt{mushroom}, \texttt{census}), do hưởng lợi từ chia
    sẻ tiền tố trong FP-tree mà Relim không có;
    \item Khoảng cách giữa Relim và Apriori giãn rộng khi $\sigma$ giảm.
\end{itemize}
Có thể tóm gọn: trên dữ liệu thưa, $|\mathcal{F}|$ vừa phải và
$\sum_X \mathrm{sup}(X)$ tăng tuyến tính theo $1/\sigma$; thực nghiệm phù hợp
với cận lý thuyết~\eqref{eq:relim-overall-time}.

\paragraph{Các yếu tố ảnh hưởng chính đến hiệu năng.}
Từ phân tích trên có thể nêu rõ:
\begin{itemize}
    \item \textbf{Ngưỡng $\sigma$}: ảnh hưởng \textit{lớn nhất} thông qua kích
    thước $|\mathcal{F}|$. Quan hệ $|\mathcal{F}|$ với $\sigma$ thường là siêu
    tuyến tính (tăng đột biến khi $\sigma$ giảm dưới một ngưỡng nào đó).
    \item \textbf{Kích thước CSDL $n$}: ảnh hưởng \textit{tuyến tính} đến cả
    pha tiền xử lý ($N = nw$) và pha đệ quy ($\sum_X \mathrm{sup}(X)$ tỉ lệ
    với $n$ khi cố định mật độ).
    \item \textbf{Chiều rộng giao dịch trung bình $w$}: tăng $w$ làm $N$ tăng,
    đồng thời tăng kích thước các CSDL điều kiện (giao dịch dài hơn $\Rightarrow$
    nhiều mắt xích sống sót qua nhiều tầng đệ quy). Đây là yếu tố nhạy với
    Relim hơn so với FP-Growth (do FP-tree chia sẻ tiền tố).
    \item \textbf{Số hạng mục phổ biến $k$}: quyết định độ sâu đệ quy tối đa và
    cận trên xấu nhất $2^k$.
    \item \textbf{Mật độ dữ liệu (dense vs sparse)}: dữ liệu dày sinh nhiều
    tập phổ biến hơn cùng $\sigma$; thực nghiệm xác nhận Relim chậm hơn
    FP-Growth trên dữ liệu dày như \texttt{chess} hay
    \texttt{mushroom}~\cite[Figures 5-6]{borgelt2005}.
\end{itemize}


\subsubsection*{Độ phức tạp không gian}

\paragraph{Bộ nhớ tĩnh.}
CSDL gốc lưu dưới dạng $n$ mảng số nguyên có tổng kích thước $N$ phần tử, cộng
mảng $A$ kích thước $k$. Bộ nhớ cố định: $O(N + k)$. Mỗi giao dịch chỉ được lưu
\textit{một lần}; các bước đệ quy sau không sao chép nội dung giao dịch (chỉ
dịch con trỏ vào mảng đã có), khác hẳn FP-Growth cần dựng cây mới ở mỗi cấp.

\paragraph{Bộ nhớ động cho mắt xích danh sách.}
Đây là điểm thú vị nhất trong phân tích bộ nhớ. Gọi $d$ là độ sâu đệ quy tối
đa ($d \leq k$). Khi gọi đệ quy ở dòng~\ref{line:recurse}, mảng $B$ ở cấp con
\textit{cùng tồn tại} với mảng $A$ ở cấp cha trong thời gian chạy nhánh con.
Tuy nhiên, do thiết kế phân loại theo hạng mục đầu:
\begin{equation}
    \sum_{i=1}^{k} \#\{\text{mắt xích trong } A[i]\} \;\leq\; n,
    \label{eq:lists-leq-n}
\end{equation}
nên \textit{tổng số mắt xích sống đồng thời ở mỗi cấp đệ quy không vượt quá
$n$}. Sau cấp đệ quy, một số giao dịch bị loại bỏ (do thân trở thành rỗng), nên
ở các cấp đệ quy con, số mắt xích thường nhỏ hơn $n$ đáng kể.

\noindent
Cận trên không gian (chưa tối ưu):
\begin{equation}
    S_{\text{naive}} = O(N + d \cdot n).
\end{equation}
Borgelt~\cite[Section~2.4]{borgelt2005} đề xuất tối ưu \textit{cấp phát theo
khối} (\textit{block allocation}): cấp phát $n$ mắt xích cho mỗi cấp đệ quy
trong một mảng liền kề, và lưu các mảng này ở một vùng nhớ toàn cục để
\textit{tái sử dụng} khi đệ quy quay lui. Cận trên vẫn là $O(d \cdot n)$,
nhưng (i) tránh phân mảnh, (ii) không phải gọi
\texttt{malloc}/\texttt{new} ở mỗi cấp đệ quy - một thắng lợi lớn về hằng số
trong thực tiễn.

\paragraph{Trường hợp tốt nhất, trung bình, xấu nhất.}
\begin{itemize}
    \item \textit{Tốt nhất:} $d = 0$ (không có hạng mục phổ biến) hoặc
    $d = O(\log n)$ trên dữ liệu phân tán đều: $S = O(N)$.
    \item \textit{Trung bình:} $d$ tỉ lệ với chiều dài tập phổ biến dài nhất.
    Trên dữ liệu thực, $d$ thường nhỏ hơn $k$ đáng kể: $S = O(N + d\cdot n)$
    với $d \ll k$.
    \item \textit{Xấu nhất:} $d = k$ khi $\sigma$ rất thấp:
    $S = O(N + k \cdot n) = O(N + kn)$.
\end{itemize}

\paragraph{So sánh không gian với các thuật toán đương thời.}
Apriori cần lưu các tập ứng viên ở mỗi cấp $k$, có thể vượt xa $|\mathcal{F}|$.
FP-Growth lưu FP-tree gốc (chia sẻ tiền tố nên tốt trên dày, kém trên thưa)
cộng các cây điều kiện trên đường đệ quy - chi phí khởi tạo và tái cấu trúc
cây có hằng số lớn. Relim đặt giữa: không có lợi thế chia sẻ tiền tố như
FP-Growth, nhưng cấu trúc danh sách liên kết đơn cực kỳ nhẹ (mỗi mắt xích chỉ
2~con trỏ, tức 8~byte trên máy 32-bit \cite[Section~2.4]{borgelt2005}); kết
hợp với tối ưu Counter-Only và block allocation, Relim thường tiết kiệm bộ nhớ
hơn FP-Growth trên dữ liệu thưa.
\subsection{Vị trí của Relim trong dòng phát triển FIM}
\label{subsec:history}
 
Để đánh giá đúng đóng góp của Relim, cần đặt nó trong bức tranh tổng thể của lĩnh
vực khai phá tập phổ biến từ những năm 1990 đến đầu những năm 2000.

\begin{figure}[H]
\centering
\begin{tikzpicture}[
    scale=0.9,
    every node/.style={font=\small},
    timeline/.style={draw=black!70, line width=2pt, ->, >=Stealth},
    milestone/.style={circle, fill=stepblue, minimum size=10pt, inner sep=0pt},
    year/.style={above=8pt, font=\bfseries\color{black!70}},
    algo/.style={below=8pt, font=\bfseries, align=center}
]
    % Draw timeline arrow
    \draw[timeline] (0,0) -- (14,0);
    
    % Milestones
    \node[milestone] (m1) at (1,0) {};
    \node[year] at (m1) {1993};
    \node[algo] at (m1) {Apriori};
    
    \node[milestone] (m2) at (3.5,0) {};
    \node[year] at (m2) {1997};
    \node[algo] at (m2) {Eclat};
    
    \node[milestone] (m3) at (6,0) {};
    \node[year] at (m3) {2000};
    \node[algo] at (m3) {FP-Growth};
    
    \node[milestone] (m4) at (8.5,0) {};
    \node[year] at (m4) {2001};
    \node[algo] at (m4) {H-Mine};
    
    \node[milestone, fill=red!80] (m5) at (11,0) {}; % Highlight Relim
    \node[year] at (m5) {2005};
    \node[algo, text=red!80!black] at (m5) {RELIM};
    
    \node[milestone] (m6) at (13.5,0) {};
    \node[year] at (m6) {2008+};
    \node[algo] at (m6) {PFP};

\end{tikzpicture}
\caption{Dòng thời gian phát triển của các thuật toán Khai phá tập phổ biến tiêu biểu.}
\label{fig:fim-timeline}
\end{figure}
 
\subsubsection{Thế hệ thứ nhất: Sinh tập ứng viên theo chiều rộng}

\textbf{Apriori} (Agrawal và Swami, 1993~\cite{apriori1993}) đặt nền móng cho FIM
bằng phương pháp sinh tập ứng viên theo tầng (level-wise). Thuật toán này cần $k$
lần quét cơ sở dữ liệu và số ứng viên tăng theo hàm mũ khi $\sigma$ giảm. 

\noindent
\textbf{Eclat} (Zaki, 1997~\cite{eclat1997}) chuyển sang biểu diễn dọc (vertical
format): mỗi hạng mục $i$ được gắn tập định danh $\mathrm{TID}(i)$; support tính
bằng phép giao: $\mathrm{sup}(X) = |\bigcap_{i \in X} \mathrm{TID}(i)|$. Khi dữ
liệu dày, TID-set lớn khiến phép giao tốn kém.
 
\subsubsection{Thế hệ thứ hai: Thay thế việc sinh ứng viên bằng đệ quy}
 
\textbf{FP-Growth} (Han, Pei và Yin, 2000~\cite{fpgrowth2000}) nén $\mathcal{D}$
vào FP-tree (hai lần quét), khai phá đệ quy các cây điều kiện mà không sinh ứng
viên. Rất mạnh trên dữ liệu dày nhờ chia sẻ tiền tố, nhưng việc xây dựng và tái
cấu trúc FP-tree đòi hỏi cài đặt phức tạp.

\noindent
\textbf{H-Mine} (Pei et al., 2001~\cite{hmine2001}) dùng siêu liên kết (hyperlinks)
thay cây tường minh, đơn giản hơn FP-Growth về cài đặt nhưng vẫn cần quản lý mạng
con trỏ phức tạp qua nhiều tầng.
 
\subsubsection{Đóng góp của Relim (Borgelt, 2005)}
 
Borgelt~\cite{borgelt2005} đề xuất Relim với triết lý \textit{tối giản hóa}: kế
thừa nguyên lý đệ quy loại trừ nhưng loại bỏ hoàn toàn cây và siêu liên kết. Toàn
bộ thuật toán hoạt động trên một mảng danh sách liên kết đơn và một hàm đệ quy
duy nhất với số dòng mã rất nhỏ.
 
\subsubsection{Các hướng tiếp cận cùng thời}
 
Đánh đổi cho sự đơn giản là Relim không khai thác được cơ chế chia sẻ tiền tố của
FP-tree. Trên dữ liệu dày ($\mathtt{chess}$, $\mathtt{mushroom}$), FP-Growth vượt
xa Relim vì có thể biểu diễn nhiều giao dịch bằng một nhánh dùng chung.

\noindent
Các hướng tiếp cận khác cùng giai đoạn gồm:
(i)~\textbf{dEclat} (Zaki và Gouda, 2003~\cite{declat2003}) thay TID-set bằng \textit{diffset} để
tiết kiệm bộ nhớ trên dữ liệu dày;
(ii)~\textbf{LCM} (Uno et al., 2004~\cite{lcm2004}) khai phá qua các tập đóng, đạt hiệu năng xuất
sắc trong các cuộc thi chuẩn FIM.
 
\subsubsection{Các hướng phát triển sau Relim}

Sau Relim, nghiên cứu FIM tiếp tục tập trung vào ba hướng chính:
(i)~tối ưu biểu diễn dữ liệu,
(ii)~khai phá tập đóng/tối đại,
và
(iii)~khai phá song song trên dữ liệu lớn.

\noindent
\textbf{AFOPT} (Liu et al., 2006~\cite{afopt2006}) cải tiến FP-Growth bằng cấu trúc
\textit{Ascending Frequency Ordered Prefix Tree}, giúp giảm chi phí xây dựng
conditional tree và tăng hiệu quả trên dữ liệu dày.

\noindent
Các thuật toán khai phá tập đóng như \textbf{Charm} (Zaki và Hsiao,
2002~\cite{charm2002}) và các biến thể của \textbf{LCM}~\cite{lcm2004} tiếp tục
được tối ưu nhằm giảm số lượng mẫu cần sinh, đặc biệt hiệu quả khi số tập phổ
biến rất lớn.

\noindent
Trong giai đoạn dữ liệu lớn (\textit{big data}), các thuật toán như
\textbf{Parallel FP-Growth (PFP)} (Li et al., 2008~\cite{pfp2008}) trên
Hadoop/Spark và các phương pháp GPU-based FIM được phát triển nhằm mở rộng khả
năng xử lý lên hàng tỷ giao dịch.

\noindent
Gần đây, nghiên cứu FIM còn mở rộng sang:
(i)~khai phá luồng dữ liệu (\textit{stream mining}),
(ii)~khai phá tập hữu ích cao (\textit{high-utility itemset mining}),
và
(iii)~khai phá trên dữ liệu không chắc chắn (\textit{uncertain data mining}).
 
Bảng~\ref{tab:algo-comparison} tóm tắt vị trí của Relim trong dòng phát triển này.
 
\begin{table}[H]
  \centering
  \caption{So sánh các thuật toán FIM qua các thế hệ}
  \label{tab:algo-comparison}
  \begin{tabularx}{\textwidth}{lllXX}
    \toprule
    \textbf{Thuật toán} & \textbf{Năm} & \textbf{Cấu trúc DL} & \textbf{Ưu điểm} & \textbf{Hạn chế} \\
    \midrule
    Apriori     & 1993 & Horizontal  & Đơn giản, nền tảng lý thuyết & Nhiều lần quét, sinh nhiều ứng viên \\
    \addlinespace
    Eclat       & 1997 & TID-set dọc & Một lần đọc; tốt trên thưa   & TID-set lớn khi dense \\
    \addlinespace
    FP-Growth   & 2000 & FP-tree     & Rất nhanh trên dense         & Cài đặt phức tạp, tốn RAM \\
    \addlinespace
    H-Mine      & 2001 & H-struct    & Nhanh hơn Apriori            & Con trỏ phức tạp nhiều tầng \\
    \addlinespace
    dEclat      & 2003 & Diffset dọc & Tiết kiệm RAM trên dense     & Phức tạp hơn Eclat \\
    \addlinespace
    LCM         & 2004 & Closed sets & Hiệu năng tốt nhất thực tế   & Phức tạp về lý thuyết \\
    \addlinespace
    \textbf{Relim} & \textbf{2005} & \textbf{Mảng linked list} & \textbf{Cực kỳ đơn giản; ít RAM; tốt trên thưa} & \textbf{Chậm hơn FP-Growth trên dense} \\
    \addlinespace
    AFOPT       & 2006 & Prefix tree tối ưu & Tăng tốc FP-Growth & Cấu trúc vẫn phức tạp \\
    \addlinespace
    PFP         & 2008 & Distributed FP-tree & Xử lý dữ liệu lớn song song & Overhead phân tán \\
    \bottomrule
  \end{tabularx}
\end{table}

% ============================================================
% CHƯƠNG 2: VÍ DỤ MINH HỌA TAY
% ============================================================
\chapter{VÍ DỤ MINH HỌA TAY}
 
Chương này trình bày hai ví dụ thực hiện tay hoàn toàn để làm rõ cơ chế hoạt động
của Relim. \textbf{Ví dụ~\ref{sec:ex1}} minh họa luồng xử lý tổng quát.
\textbf{Ví dụ~\ref{sec:ex2}} tập trung vào tình huống đặc biệt \textit{Counter-Only},
một tối ưu hóa đặc trưng của Relim.
 
% ============================================================
\section{Ví dụ 1: Minh họa tổng quát}
\label{sec:ex1}
% ============================================================
 
\subsection{Thiết lập}
 
\begin{center}
\begin{tikzpicture}
\node[ rounded corners=5pt,
      inner xsep=12pt, inner ysep=8pt, font=\small] {
  $\mathcal{I} = \{a,b,c,d,e,f\}$ \qquad
  $|\mathcal{D}| = 6$ \qquad
  $\sigma = 3$
};
\end{tikzpicture}
\end{center}
 
\begin{table}[H]
  \centering
  \caption{Cơ sở dữ liệu gốc (Ví dụ 1)}
  \label{tab:ex1-original}
  \begin{tabular}{cc}
    \toprule
    \textbf{TID} & \textbf{Giao dịch} \\
    \midrule
    $T_1$ & $\{a, b, c, d\}$ \\
    $T_2$ & $\{b, c, e, f\}$    \\
    $T_3$ & $\{a, b, d\}$    \\
    $T_4$ & $\{b, c, d, e\}$ \\
    $T_5$ & $\{a, c, d\}$    \\
    $T_6$ & $\{b, d, e\}$    \\
    \bottomrule
  \end{tabular}
\end{table}
 
\subsection{Bước tiền xử lý}
 
\paragraph{Đếm tần suất.}
 
\begin{table}[H]
  \centering
  \caption{Tần suất các hạng mục và thứ tự sắp xếp}
  \label{tab:ex1-freq}
  \begin{tabular}{ccccccc}
    \toprule
    Hạng mục & $f$ & $a$ & $e$ & $c$ & $b$ & $d$ \\
    \midrule
    Tần suất & 1 & 3 & 3 & 4 & 5 & 5 \\
    Thứ tự & -- & 1 & 2 & 3 & 4 & 5 \\
    Phổ biến? & $\times$ & {\checkmark} & {\checkmark} & {\checkmark} & {\checkmark} & {\checkmark} \\
    \bottomrule
  \end{tabular}
\end{table}
 
Hạng mục $f$ có tần suất $1 < \sigma = 3$ nên bị \textbf{loại bỏ} khỏi cơ sở dữ liệu. Các hạng mục khác được giữ nguyên.
Thứ tự sắp xếp tần suất tăng dần (tie-break theo bảng chữ cái):
\[
  \underbrace{a}_{3} \;<\; \underbrace{e}_{3} \;<\; \underbrace{c}_{4} \;<\; \underbrace{b}_{5} \;\approx\; \underbrace{d}_{5}.
\]
 
Cơ sở dữ liệu sau tiền xử lý được thể hiện trong Bảng~\ref{tab:ex1-sorted}.
 
\begin{table}[H]
  \centering
  \caption{Cơ sở dữ liệu sau sắp xếp ($a < e < c < b < d$)}
  \label{tab:ex1-sorted}
  \begin{tabular}{ccc}
    \toprule
    \textbf{TID} & \textbf{Giao dịch (đã sắp xếp)} & \textbf{Hạng mục đứng đầu} \\
    \midrule
    $T_1$ & $a,\; c,\; b,\; d$ & $a$ \\
    $T_2$ & $e,\; c,\; b$       & $e$ \\
    $T_3$ & $a,\; b,\; d$       & $a$ \\
    $T_4$ & $e,\; c,\; b,\; d$  & $e$ \\
    $T_5$ & $a,\; c,\; d$       & $a$ \\
    $T_6$ & $e,\; b,\; d$       & $e$ \\
    \bottomrule
  \end{tabular}
\end{table}
 
Trạng thái mảng $A$ sau khởi tạo được thể hiện trong Bảng~\ref{tab:ex1-init}.
 
\begin{table}[H]
  \centering
  \caption{Mảng $A$ sau khởi tạo - Ví dụ 1}
  \label{tab:ex1-init}
  \begin{tabular}{cclp{5.5cm}}
    \toprule
    \textbf{Hạng mục} & $\mathbf{cnt}$ & \textbf{Danh sách (hậu tố)} & \textbf{Nguồn} \\
    \midrule
    $a$ & 3 & $(c,b,d) \to (b,d) \to (c,d)$ & $T_1, T_3, T_5$ \\
    $e$ & 3 & $(c,b) \to (c,b,d) \to (b,d)$ & $T_2, T_4, T_6$ \\
    $c$ & 0 & \textit{rỗng} & - \\
    $b$ & 0 & \textit{rỗng} & - \\
    $d$ & 0 & \textit{rỗng} & - \\
    \bottomrule
  \end{tabular}
\end{table}
 
\subsection{Thực hiện thuật toán từng bước}
 
\medskip
% -- BƯỚC 1 --
\begin{tikzpicture}
\node[ rounded corners=4pt,
      inner xsep=10pt, inner ysep=6pt, text width=\linewidth-28pt] {
\textbf{{Bước 1 - Xử lý hạng mục $a$ \hfill $\mathrm{sup}(a) = 3 \geq \sigma$}}
 
\smallskip
\textbf{Ghi nhận:} $\{a\} \in \mathcal{F}$ với support = 3.
 
\smallskip
\textbf{Xây dựng $\mathit{DB}_a$} từ hậu tố các giao dịch trong $A[a]$:
 
\medskip
\begin{center}
\begin{tabular}{ccc}
  \toprule
  Giao dịch gốc & Hậu tố (sau khi ẩn $a$) & Đưa vào $\mathit{DB}_a$ \\
  \midrule
  $T_1: (a,c,b,d)$ & $(c,b,d)$ & $j{=}c$; còn $(b,d)$ $\to$ mắt xích $\to B_a[c]$; $B_a[c].cnt\mathrel{+}=1$ \\
  $T_3: (a,b,d)$   & $(b,d)$   & $j{=}b$; còn $(d)$ $\to$ mắt xích $\to B_a[b]$; $B_a[b].cnt\mathrel{+}=1$ \\
  $T_5: (a,c,d)$   & $(c,d)$   & $j{=}c$; còn $(d)$ $\to$ mắt xích $\to B_a[c]$; $B_a[c].cnt\mathrel{+}=1$ \\
  \bottomrule
\end{tabular}
\end{center}
 
\medskip
\textbf{Trạng thái $\mathit{DB}_a$ sau xây dựng:}
$B_a[c].cnt = 2$ (danh sách: $(b,d)$, $(d)$),\quad
$B_a[b].cnt = 1$ (danh sách: $(d)$),\quad
$B_a[d].cnt = 0$ (rỗng).
 
\medskip
Đệ quy trên $\mathit{DB}_a$:
\begin{itemize}[topsep=2pt, itemsep=1pt]
  \item Xử lý $c$ trong $\mathit{DB}_a$: $B_a[c].cnt = 2 < \sigma$ $\Rightarrow$ bỏ qua.
  Phân bổ lại từ $B_a[c]$: $(b,d) \to B_a[b]$, $B_a[b].cnt\mathrel{+}=1$; $(d) \to$ Counter-Only: $B_a[d].cnt\mathrel{+}=1$.
  \item Xử lý $b$ trong $\mathit{DB}_a$: $B_a[b].cnt = 2 < \sigma$ $\Rightarrow$ bỏ qua.
  Phân bổ lại từ $B_a[b]$: $(d) \to$ Counter-Only $B_a[d].cnt\mathrel{+}=1$; $(d) \to$ Counter-Only $B_a[d].cnt\mathrel{+}=1$. Tổng $B_a[d].cnt = 3$.
  \item Xử lý $d$ trong $\mathit{DB}_a$: $B_a[d].cnt = 3 \geq \sigma$ $\Rightarrow$ \textbf{ghi nhận} $\{a,d\} \in \mathcal{F}$ (support = 3).
\end{itemize}
 
\medskip
\textbf{Phân bổ lại (reassign) từ $A[a]$ về mảng $A$:}
Mỗi mắt xích trong $A[a]$ được chuyển sang danh sách theo hạng mục dẫn đầu của nó:
$(c,b,d) \to A[c]$, $A[c].cnt\mathrel{+}=1$;\quad
$(b,d) \to A[b]$, $A[b].cnt\mathrel{+}=1$;\quad
$(c,d) \to A[c]$, $A[c].cnt\mathrel{+}=1$.
\\Cập nhật tổng: $A[c].cnt \mathrel{+}= 2$,\; $A[b].cnt \mathrel{+}= 1$.
};
\end{tikzpicture}
 
\medskip
 
\begin{table}[H]
  \centering
  \caption{Trạng thái mảng $A$ sau Bước 1 (sau reassign $a$)}
  \label{tab:ex1-after1}
  \begin{tabular}{ccl}
    \toprule
    \textbf{Hạng mục} & $\mathbf{cnt}$ & \textbf{Danh sách (hậu tố)} \\
    \midrule
    $a$ & 0 \textit{(đã xử lý)} & \textit{rỗng} \\
    $e$ & 3 & $(c,b) \to (c,b,d) \to (b,d)$ \\
    $c$ & 2 & $(b,d)$ [từ $T_1$] $\to$ $(d)$ [từ $T_5$] \\
    $b$ & 1 & $(d)$ [từ $T_3$] \\
    $d$ & 0 & \textit{rỗng} \\
    \bottomrule
  \end{tabular}
\end{table}
 
\medskip
% -- BƯỚC 2 --
\begin{tikzpicture}
\node[ rounded corners=4pt,
      inner xsep=10pt, inner ysep=6pt, text width=\linewidth-28pt] {
\textbf{{Bước 2 - Xử lý hạng mục $e$ \hfill $\mathrm{sup}(e) = 3 \geq \sigma$}}
 
\smallskip
\textbf{Ghi nhận:} $\{e\} \in \mathcal{F}$ với support = 3.
 
\smallskip
\textbf{Xây dựng $\mathit{DB}_e$} từ hậu tố trong $A[e]$:
 
\medskip
\begin{center}
\begin{tabular}{ccc}
  \toprule
  Giao dịch gốc & Hậu tố (ẩn $e$) & Đưa vào $\mathit{DB}_e$ \\
  \midrule
  $T_2: (e,c,b)$    & $(c,b)$   & Mắt xích $\to B_e[c]$ \\
  $T_4: (e,c,b,d)$  & $(c,b,d)$ & Mắt xích $\to B_e[c]$ \\
  $T_6: (e,b,d)$    & $(b,d)$   & Mắt xích $\to B_e[b]$ \\
  \bottomrule
\end{tabular}
\end{center}
 
\medskip
\textbf{Đếm trong $\mathit{DB}_e$:}
$\mathrm{sup}_{DB_e}(c) = 2$,\quad $\mathrm{sup}_{DB_e}(b) = 3$,\quad $\mathrm{sup}_{DB_e}(d) = 2$.
 
\medskip
Chỉ $b$ đạt ngưỡng. Đệ quy trên $\mathit{DB}_e$:
\begin{itemize}[topsep=2pt, itemsep=1pt]
  \item $c$: $\mathrm{sup} = 2 < \sigma$ $\Rightarrow$ bỏ qua.
  \item $b$: $\mathrm{sup} = 3 \geq \sigma$ $\Rightarrow$ \textbf{ghi nhận} $\{e,b\} \in \mathcal{F}$ (support = 3).
  \quad Đệ quy sâu hơn: $\mathrm{sup}_{DB_{eb}}(d) = 2 < \sigma$ $\Rightarrow$ dừng.
\end{itemize}
 
\medskip
\textbf{Phân bổ lại từ $A[e]$ về mảng $A$:}
$(c,b) \to A[c]$, $A[c].cnt\mathrel{+}=1$;\quad
$(c,b,d) \to A[c]$, $A[c].cnt\mathrel{+}=1$;\quad
$(b,d) \to A[b]$, $A[b].cnt\mathrel{+}=1$.
\\Cập nhật tổng: $A[c].cnt \mathrel{+}= 2$,\; $A[b].cnt \mathrel{+}= 1$.
};
\end{tikzpicture}
 
\medskip
 
\begin{table}[H]
  \centering
  \caption{Trạng thái mảng $A$ sau Bước 2 (sau reassign $e$)}
  \label{tab:ex1-after2}
  \begin{tabular}{ccl}
    \toprule
    \textbf{Hạng mục} & $\mathbf{cnt}$ & \textbf{Danh sách (hậu tố)} \\
    \midrule
    $a$ & 0 \textit{(đã xử lý)} & \textit{rỗng} \\
    $e$ & 0 \textit{(đã xử lý)} & \textit{rỗng} \\
    $c$ & 4 & $(b,d)\to(d)$ [từ $a$] + $(b)\to(b,d)$ [từ $e$] \\
    $b$ & 2 & $(d)$ [từ $a$] + $(d)$ [từ $e$] \\
    $d$ & 0 & \textit{rỗng} \\
    \bottomrule
  \end{tabular}
\end{table}
 
\medskip
% -- BƯỚC 3 --
\begin{tikzpicture}
\node[ rounded corners=4pt,
      inner xsep=10pt, inner ysep=6pt, text width=\linewidth-28pt] {
\textbf{{Bước 3 - Xử lý hạng mục $c$ \hfill $\mathrm{sup}(c) = 4 \geq \sigma$}}
 
\smallskip
\textbf{Ghi nhận:} $\{c\} \in \mathcal{F}$ với support = 4.
 
\smallskip
\textbf{Xây dựng $\mathit{DB}_c$}:
$A[c]$ sau bước reassign của $a$ và $e$ chứa hậu tố từ $T_1, T_2, T_4, T_5$.
 
\medskip
\begin{center}
\begin{tabular}{ccc}
  \toprule
  Nguồn & Hậu tố trong $A[c]$ & Khi ẩn $c$ $\to$ đưa vào $\mathit{DB}_c$ \\
  \midrule
  $T_1$ (từ $a$) & $(c,b,d)$ & $(b,d)$ $\to$ mắt xích $\to B_c[b]$ \\
  $T_5$ (từ $a$) & $(c,d)$   & $(d)$ $\to$ \textbf{Counter-Only}: $B_c[d].\mathit{cnt}\mathrel{+}=1$ \\
  $T_2$ (từ $e$) & $(c,b)$   & $(b)$ $\to$ \textbf{Counter-Only}: $B_c[b].\mathit{cnt}\mathrel{+}=1$ \\
  $T_4$ (từ $e$) & $(c,b,d)$ & $(b,d)$ $\to$ mắt xích $\to B_c[b]$ \\
  \bottomrule
\end{tabular}
\end{center}
 
\medskip
\textbf{Đếm trong $\mathit{DB}_c$:}
$\mathrm{sup}_{DB_c}(b) = 3$ ($T_1, T_2, T_4$),\quad $\mathrm{sup}_{DB_c}(d) = 3$ ($T_1, T_4, T_5$).
 
\medskip
Cả $b$ và $d$ đều $\geq \sigma$. Đệ quy trên $\mathit{DB}_c$:
\begin{itemize}[topsep=2pt, itemsep=1pt]
  \item $b$: $\mathrm{sup} = 3 \geq \sigma$ $\Rightarrow$ \textbf{ghi nhận} $\{c,b\} \in \mathcal{F}$ (support = 3).
    \\ Đệ quy $\mathit{DB}_{cb}$: hậu tố $(d)$ và $(d)$ $\to$ $\mathrm{sup}_{DB_{cb}}(d) = 2 < \sigma$ $\Rightarrow$ dừng.
  \item $d$: $\mathrm{sup} = 3 \geq \sigma$ $\Rightarrow$ \textbf{ghi nhận} $\{c,d\} \in \mathcal{F}$ (support = 3).
\end{itemize}
 
\medskip
\textbf{Phân bổ lại từ $A[c]$ về mảng $A$:}
Mỗi mắt xích trong $A[c]$ được chuyển theo hạng mục dẫn đầu:
$(b,d) \to A[b]$, $A[b].cnt\mathrel{+}=1$;\quad
$(d) \to$ Counter-Only: $A[d].cnt\mathrel{+}=1$;\quad
$(b) \to$ Counter-Only: $A[b].cnt\mathrel{+}=1$;\quad
$(b,d) \to A[b]$, $A[b].cnt\mathrel{+}=1$.
\\Cập nhật tổng: $A[b].cnt \mathrel{+}= 3$,\; $A[d].cnt \mathrel{+}= 1$.
};
\end{tikzpicture}
 
\medskip
 
\begin{table}[H]
  \centering
  \caption{Trạng thái mảng $A$ sau Bước 3 (sau reassign $c$)}
  \label{tab:ex1-after3}
  \begin{tabular}{ccl}
    \toprule
    \textbf{Hạng mục} & $\mathbf{cnt}$ & \textbf{Danh sách (hậu tố)} \\
    \midrule
    $a$ & 0 \textit{(đã xử lý)} & \textit{rỗng} \\
    $e$ & 0 \textit{(đã xử lý)} & \textit{rỗng} \\
    $c$ & 0 \textit{(đã xử lý)} & \textit{rỗng} \\
    $b$ & 5 & $(d)$[từ $a$] + $(d)$[từ $e$] + $(d)$[từ $T_1$ qua $c$] + $(d)$[từ $T_4$ qua $c$] [+ 1 counter-only $T_2$] \\
    $d$ & 1 & \textit{rỗng (counter-only từ $T_5$ qua $c$)} \\
    \bottomrule
  \end{tabular}
\end{table}
 
\medskip
 
% -- BƯỚC 4 --
\begin{tikzpicture}
\node[rounded corners=4pt,
      inner xsep=10pt, inner ysep=6pt, text width=\linewidth-28pt] {
\textbf{{Bước 4 - Xử lý hạng mục $b$ \hfill $\mathrm{sup}(b) = 5 \geq \sigma$}}
 
\smallskip
\textbf{Ghi nhận:} $\{b\} \in \mathcal{F}$ với support = 5.
 
\smallskip
\textbf{Nguồn gốc $A[b].\mathit{cnt} = 5$:}
\begin{itemize}[topsep=2pt, itemsep=1pt]
  \item Bước 1 (reassign $a$): $T_3$ có hậu tố $(b,d)$ $\to$ $A[b].\mathit{cnt} \mathrel{+}= 1$, mắt xích $(d)$.
  \item Bước 2 (reassign $e$): $T_6$ có hậu tố $(b,d)$ $\to$ $A[b].\mathit{cnt} \mathrel{+}= 1$, mắt xích $(d)$.
  \item Bước 3 (reassign $c$): $T_1$ hậu tố $(b,d)$ $\to$ mắt xích $(d)$; $T_4$ hậu tố $(b,d)$ $\to$ mắt xích $(d)$; $T_2$ hậu tố $(b)$ $\to$ \textbf{Counter-Only}. Tổng: $A[b].\mathit{cnt} \mathrel{+}= 3$.
\end{itemize}
 
\smallskip
\textbf{Xây dựng $\mathit{DB}_b$}:
 
\medskip
\begin{center}
\begin{tabular}{ccc}
  \toprule
  Nguồn & Hậu tố trong $A[b]$ & Khi ẩn $b$ $\to$ đưa vào $\mathit{DB}_b$ \\
  \midrule
  $T_3$ (từ $a$) & $(d)$ & \textbf{Counter-Only}: $B_b[d].\mathit{cnt} \mathrel{+}= 1$ \\
  $T_6$ (từ $e$) & $(d)$ & \textbf{Counter-Only}: $B_b[d].\mathit{cnt} \mathrel{+}= 1$ \\
  $T_1$ (từ $c$) & $(d)$ & \textbf{Counter-Only}: $B_b[d].\mathit{cnt} \mathrel{+}= 1$ \\
  $T_4$ (từ $c$) & $(d)$ & \textbf{Counter-Only}: $B_b[d].\mathit{cnt} \mathrel{+}= 1$ \\
  \bottomrule
\end{tabular}
\end{center}
 
\medskip
\textbf{Đếm trong $\mathit{DB}_b$:}
$\mathrm{sup}_{DB_b}(d) = 4$ (từ $T_1, T_3, T_4, T_6$; \textit{tất cả đều Counter-Only vì hậu tố chỉ còn 1 item}).
 
\medskip
$d$: $\mathrm{sup} = 4 \geq \sigma$ $\Rightarrow$ \textbf{ghi nhận} $\{b,d\} \in \mathcal{F}$ (support = 4).
\\ Đệ quy $\mathit{DB}_{bd}$: rỗng (mọi hậu tố đều Counter-Only, không còn mắt xích nào) $\Rightarrow$ dừng.
 
\medskip
\textbf{Phân bổ lại:} 4 mắt xích $(d)$ chuyển sang $A[d]$.
Cập nhật: $A[d].\mathit{cnt} \mathrel{+}= 4$.
};
\end{tikzpicture}
 
\medskip
 
\begin{table}[H]
  \centering
  \caption{Trạng thái mảng $A$ sau Bước 4 (sau reassign $b$)}
  \label{tab:ex1-after4}
  \begin{tabular}{ccl}
    \toprule
    \textbf{Hạng mục} & $\mathbf{cnt}$ & \textbf{Danh sách} \\
    \midrule
    $a$ & 0 \textit{(đã xử lý)} & \textit{rỗng} \\
    $e$ & 0 \textit{(đã xử lý)} & \textit{rỗng} \\
    $c$ & 0 \textit{(đã xử lý)} & \textit{rỗng} \\
    $b$ & 0 \textit{(đã xử lý)} & \textit{rỗng} \\
    $d$ & 5 & \textit{rỗng (counter-only tích lũy)} \\
    \bottomrule
  \end{tabular}
\end{table}
 
\medskip
 
% -- BƯỚC 5 --
\begin{tikzpicture}
\node[ rounded corners=4pt,
      inner xsep=10pt, inner ysep=6pt, text width=\linewidth-28pt] {
\textbf{Bước 5 - Xử lý hạng mục $d$ \hfill $\mathrm{sup}(d) = 5 \geq \sigma$}
 
\smallskip
\textbf{Ghi nhận:} $\{d\} \in \mathcal{F}$ với support = 5.
 
\smallskip
\textbf{Nguồn gốc $A[d].\mathit{cnt} = 5$} (tích lũy qua Counter-Only và reassign dây chuyền):
\begin{itemize}[topsep=2pt, itemsep=1pt]
  \item Bước 3 (reassign $c$): $(d)$ từ $T_5$ (qua $c$) $\to$ Counter-Only $A[d].cnt\mathrel{+}=1$.
  \item Bước 4 (reassign $b$): 4 mắt xích $(d)$ từ $T_1, T_3$ (qua $a{\to}c{\to}b$), $T_4, T_6$ (qua $e{\to}c/b$) $\to$ Counter-Only $A[d].cnt\mathrel{+}=4$.
  \item Tổng tích lũy = $1 + 4 = 5$.
\end{itemize}
 
\smallskip
$d$ là hạng mục cuối cùng trong thứ tự, không có hạng mục nào đứng sau $d$.
Danh sách $A[d]$ không có mắt xích (toàn bộ đều Counter-Only) nên không xây dựng
$\mathit{DB}_d$ và không có đệ quy. Thuật toán \textbf{kết thúc}.
};
\end{tikzpicture}
 
\subsection{Tập kết quả}
 
\begin{center}
\begin{tikzpicture}
\node[ rounded corners=5pt,
      inner xsep=12pt, inner ysep=10pt, font=\small] {
$\mathcal{F} = \bigl\{
  \{a\}{:}3,\;\; \{e\}{:}3,\;\; \{c\}{:}4,\;\; \{b\}{:}5,\;\; \{d\}{:}5,\;\;
  \{a,d\}{:}3,\;\; \{e,b\}{:}3,\;\; \{c,b\}{:}3,\;\; \{c,d\}{:}3,\;\; \{b,d\}{:}4
\bigr\}$
};
\end{tikzpicture}
\end{center}

\subsection{Cây đệ quy (Recursion Tree)}

\begin{figure}[H]
\centering
\begin{tikzpicture}[
    level 1/.style={sibling distance=3cm, level distance=1.5cm},
    level 2/.style={sibling distance=2cm, level distance=1.5cm},
    every node/.style={draw=borderblue, thick, rectangle, rounded corners, fill=stepblue!10, font=\small, minimum width=1.4cm, align=center},
    edge from parent/.style={draw, thick, ->, >=Stealth}
]

  \node[fill=gray!20, draw=black!70] {Root \\ $\mathcal{D}$}
    child { node {$a:3$} 
      child { node {$\{a, d\}:3$} }
    }
    child { node {$e:3$} 
      child { node {$\{e, b\}:3$} }
    }
    child { node {$c:4$} 
      child { node {$\{c, b\}:3$} }
      child { node {$\{c, d\}:3$} }
    }
    child { node {$b:5$} 
      child { node {$\{b, d\}:4$} }
    }
    child { node {$d:5$} };
    
\end{tikzpicture}
\caption{Cây đệ quy của thuật toán RELIM trên Ví dụ 1. Các nhánh con có support $< \sigma$ bị cắt tỉa nên không hiển thị.}
\label{fig:ex1-recursion-tree}
\end{figure}
 
\subsection{Kiểm tra chéo}
 
\begin{table}[H]
  \centering
  \caption{Kiểm tra chéo toàn bộ tập hạng mục (Ví dụ 1, $\sigma = 3$)}
  \label{tab:ex1-crosscheck}
  \begin{tabular}{lcc@{\hskip 1.5em}lcc}
    \toprule
    \textbf{Tập} & $\mathbf{sup}$ & \textbf{Kết luận}
      & \textbf{Tập} & $\mathbf{sup}$ & \textbf{Kết luận} \\
    \midrule
    \multicolumn{6}{l}{\textit{\textbf{1-tập (6 tập, gồm cả $f$ đã loại ở tiền xử lý):}}} \\
    $\{a\}$      & 3 & {\textbf{Phổ biến}} & $\{b\}$     & 5 & {\textbf{Phổ biến}} \\
    $\{e\}$      & 3 & {\textbf{Phổ biến}} & $\{d\}$     & 5 & {\textbf{Phổ biến}} \\
    $\{c\}$      & 4 & {\textbf{Phổ biến}} & $\{f\}$     & 1 & {Loại (tiền xử lý)} \\
    \midrule
    \multicolumn{6}{l}{\textit{\textbf{2-tập ($\binom{5}{2} = 10$ tập):}}} \\
    $\{a,e\}$    & 0 & {Loại}             & $\{e,b\}$   & 3 & {\textbf{Phổ biến}} \\
    $\{a,c\}$    & 2 & {Loại}             & $\{e,d\}$   & 2 & {Loại} \\
    $\{a,b\}$    & 2 & {Loại}             & $\{c,b\}$   & 3 & {\textbf{Phổ biến}} \\
    $\{a,d\}$    & 3 & {\textbf{Phổ biến}} & $\{c,d\}$   & 3 & {\textbf{Phổ biến}} \\
    $\{e,c\}$    & 2 & {Loại}             & $\{b,d\}$   & 4 & {\textbf{Phổ biến}} \\
    \midrule
    \multicolumn{6}{l}{\textit{\textbf{3-tập ($\binom{5}{3} = 10$ tập):}}} \\
    $\{a,e,c\}$  & 0 & {Loại}   & $\{a,c,d\}$  & 2 & {Loại} \\
    $\{a,e,b\}$  & 0 & {Loại}   & $\{a,b,d\}$  & 2 & {Loại} \\
    $\{a,e,d\}$  & 0 & {Loại}   & $\{e,c,b\}$  & 2 & {Loại} \\
    $\{a,c,b\}$  & 1 & {Loại}   & $\{e,c,d\}$  & 1 & {Loại} \\
    $\{e,b,d\}$  & 2 & {Loại}   & $\{c,b,d\}$  & 2 & {Loại} \\
    \midrule
    \multicolumn{6}{l}{\textit{\textbf{4-tập ($\binom{5}{4} = 5$ tập):}}} \\
    $\{a,e,c,b\}$  & 0 & {Loại}   & $\{a,c,b,d\}$  & 1 & {Loại} \\
    $\{a,e,c,d\}$  & 0 & {Loại}   & $\{e,c,b,d\}$  & 1 & {Loại} \\
    $\{a,e,b,d\}$  & 0 & {Loại}   &               &   &  \\
    \midrule
    \multicolumn{6}{l}{\textit{\textbf{5-tập ($\binom{5}{5} = 1$ tập):}}} \\
    $\{a,e,c,b,d\}$ & 0 & {Loại} &               &   &  \\
    \bottomrule
  \end{tabular}
\end{table}
 
Kết quả kiểm tra chéo xác nhận $\mathcal{F}$ do Relim tìm được là đầy đủ và chính xác.
 
% ============================================================
\section{Ví dụ 2: Tình huống đặc biệt - Counter-Only Transactions}
\label{sec:ex2}
% ============================================================
 
\subsection{Mô tả tình huống đặc biệt}
 
\begin{tikzpicture}
\node[ rounded corners=4pt,
      inner xsep=10pt, inner ysep=8pt, text width=\linewidth-24pt, font=\small] {
\textbf{Tình huống Counter-Only} xảy ra khi trong quá trình xây dựng cơ sở dữ liệu
điều kiện (\textsc{Build-Conditional-DB}) hoặc phân bổ lại (\textsc{Reassign}), một
giao dịch sau khi ẩn hạng mục dẫn đầu chỉ còn lại \emph{đúng một} hạng mục duy nhất.
Thay vì cấp phát mắt xích mới (tốn bộ nhớ và thời gian), Relim chỉ tăng biến đếm
$A[j].\mathit{cnt}$. Ví dụ này được thiết kế với nhiều giao dịch ngắn (2 item) để
Counter-Only xuất hiện \textit{tự nhiên và liên tục} trong các bước đệ quy, đặc biệt
tại Bước~3 nơi \textbf{toàn bộ} thông tin được truyền qua biến đếm mà không cấp phát
bất kỳ mắt xích nào.
};
\end{tikzpicture}
 
\subsection{Thiết lập}
 
\begin{center}
\begin{tikzpicture}
\node[ rounded corners=5pt,
      inner xsep=12pt, inner ysep=8pt, font=\small] {
  $\mathcal{I} = \{a, b, c, d\}$ \qquad
  $|\mathcal{D}| = 5$ \qquad
  $\sigma = 3$
};
\end{tikzpicture}
\end{center}
 
\begin{table}[H]
  \centering
  \caption{Cơ sở dữ liệu gốc và sau tiền xử lý (Ví dụ 2)}
  \label{tab:ex2-db}
  \begin{tabular}{cll}
    \toprule
    \textbf{TID} & \textbf{Giao dịch gốc} & \textbf{Sau tiền xử lý ($a < c < d < b$)} \\
    \midrule
    $T_1$ & $\{a, b, c, d\}$ & $a, c, d, b$ \\
    $T_2$ & $\{a, b\}$       & $a, b$        \\
    $T_3$ & $\{a, c\}$       & $a, c$        \\
    $T_4$ & $\{b, c, d\}$    & $c, d, b$     \\
    $T_5$ & $\{b, d\}$       & $d, b$        \\
    \bottomrule
  \end{tabular}
\end{table}
 
\begin{table}[H]
  \centering
  \caption{Tần suất và thứ tự sắp xếp (Ví dụ 2)}
  \label{tab:ex2-freq}
  \begin{tabular}{ccccc}
    \toprule
    Hạng mục & $a$ & $c$ & $d$ & $b$ \\
    \midrule
    Tần suất & 3 & 3 & 3 & 4 \\
    Thứ tự   & 1 & 2 & 3 & 4 \\
    \bottomrule
  \end{tabular}
\end{table}
 
\subsection{Xây dựng mảng ban đầu}
 
Theo Thuật toán~\ref{alg:relim-main}, mỗi giao dịch có $|T| > 1$ luôn tạo mắt xích
trong danh sách $A[i]$ (với $i$ là hạng mục đứng đầu). Counter-Only chỉ được áp dụng
bên trong \textsc{Build-Conditional-DB} và \textsc{Reassign}, không áp dụng tại khởi tạo.
 
\begin{table}[H]
  \centering
  \caption{Phân loại giao dịch vào mảng $A$ theo Thuật toán~\ref{alg:relim-main}}
  \label{tab:ex2-init}
  \begin{tabular}{cclcl}
    \toprule
    \textbf{TID} & \textbf{H.mục đầu} & \textbf{Hậu tố còn lại} & \textbf{$|T|$} & \textbf{Hành động} \\
    \midrule
    $T_1$ & $a$ & $(c, d, b)$ & 4 & $A[a].cnt\mathrel{+}=1$; tạo mắt xích $\to A[a]$ \\
    $T_2$ & $a$ & $(b)$       & 2 & $A[a].cnt\mathrel{+}=1$; tạo mắt xích $\to A[a]$ \\
    $T_3$ & $a$ & $(c)$       & 2 & $A[a].cnt\mathrel{+}=1$; tạo mắt xích $\to A[a]$ \\
    $T_4$ & $c$ & $(d, b)$    & 3 & $A[c].cnt\mathrel{+}=1$; tạo mắt xích $\to A[c]$ \\
    $T_5$ & $d$ & $(b)$       & 2 & $A[d].cnt\mathrel{+}=1$; tạo mắt xích $\to A[d]$ \\
    \bottomrule
  \end{tabular}
\end{table}
 
\begin{table}[H]
  \centering
  \caption{Trạng thái mảng $A$ sau khởi tạo (Ví dụ 2)}
  \label{tab:ex2-array}
  \begin{tabular}{ccll}
    \toprule
    \textbf{H.mục} & $\mathbf{cnt}$ & \textbf{Danh sách liên kết} & \textbf{Ghi chú} \\
    \midrule
    $a$ & 3 & $(c,d,b) \to (b) \to (c)$ & 3 mắt xích từ $T_1, T_2, T_3$ \\
    $c$ & 1 & $(d,b)$   & 1 mắt xích từ $T_4$ \\
    $d$ & 1 & $(b)$     & 1 mắt xích từ $T_5$ \\
    $b$ & 0 & \textit{rỗng} & Không có giao dịch nào có $b$ đứng đầu \\
    \bottomrule
  \end{tabular}
\end{table}
 
\subsection{Thực hiện thuật toán từng bước}
 
\medskip
% -- EX2 BƯỚC 1 --
\begin{tikzpicture}
\node[ rounded corners=4pt,
      inner xsep=10pt, inner ysep=6pt, text width=\linewidth-28pt] {
\textbf{{Bước 1 - Xử lý hạng mục $a$ \hfill $\mathrm{sup}(a) = 3 \geq \sigma$}}
 
\smallskip
\textbf{Ghi nhận:} $\{a\} \in \mathcal{F}$ với support = 3.
 
\smallskip
\textbf{Xây dựng $\mathit{DB}_a$} từ 3 mắt xích trong $A[a]$:
 
\medskip
\begin{center}
\begin{tabular}{clll}
  \toprule
  Nguồn & Hậu tố & Ẩn h.mục đầu & Đưa vào $\mathit{DB}_a$ \\
  \midrule
  $T_1$ & $(c,d,b)$ & $j{=}c$; còn $(d,b)$ & Mắt xích $\to B_a[c]$; $B_a[c].cnt\mathrel{+}=1$ \\
  $T_2$ & $(b)$     & $j{=}b$; còn rỗng    & \textbf{Counter-Only}: $B_a[b].cnt\mathrel{+}=1$ \\
  $T_3$ & $(c)$     & $j{=}c$; còn rỗng    & \textbf{Counter-Only}: $B_a[c].cnt\mathrel{+}=1$ \\
  \bottomrule
\end{tabular}
\end{center}
 
\medskip
\textbf{Đếm trong $\mathit{DB}_a$:}
$\mathrm{sup}_{DB_a}(c) = 2$ ($T_1$ mắt xích $+$ $T_3$ Counter-Only),\quad
$\mathrm{sup}_{DB_a}(b) = 1$ ($T_2$ Counter-Only).
 
\medskip
Không có hạng mục nào đạt $\sigma = 3$ $\Rightarrow$ \textbf{không} sinh thêm tập với tiền tố $\{a\}$.
 
\medskip
\textit{Lưu ý: Counter-Only xuất hiện tự nhiên tại đây -- $T_2$ và $T_3$ có hậu tố
chỉ còn 1 item sau khi ẩn hạng mục đầu, nên không cần cấp phát mắt xích trong $\mathit{DB}_a$.}
 
\medskip
\textbf{Phân bổ lại} 3 mắt xích từ $A[a]$:
\begin{itemize}[topsep=2pt, itemsep=1pt]
  \item $T_1$: $(c,d,b)$, $j{=}c$, còn $(d,b)$ $\to$ chèn vào $A[c]$; $A[c].cnt\mathrel{+}=1$.
  \item $T_2$: $(b)$, $j{=}b$, còn rỗng $\to$ chỉ $A[b].cnt\mathrel{+}=1$ (Counter-Only).
  \item $T_3$: $(c)$, $j{=}c$, còn rỗng $\to$ chỉ $A[c].cnt\mathrel{+}=1$ (Counter-Only).
\end{itemize}
};
\end{tikzpicture}
 
\medskip
\begin{table}[H]
  \centering
  \caption{Mảng $A$ sau Bước 1 (Ví dụ 2)}
  \label{tab:ex2-after1}
  \begin{tabular}{ccl}
    \toprule
    \textbf{H.mục} & $\mathbf{cnt}$ & \textbf{Danh sách} \\
    \midrule
    $c$ & 3 & $(d,b)$ [từ $T_4$] + $(d,b)$ [reassign $T_1$] \\
    $d$ & 1 & $(b)$ [từ $T_5$, không thay đổi] \\
    $b$ & 1 & \textit{rỗng} (Counter-Only từ reassign $T_2$) \\
    \bottomrule
  \end{tabular}
\end{table}
 
\medskip
% -- EX2 BƯỚC 2 --
\begin{tikzpicture}
\node[ rounded corners=4pt,
      inner xsep=10pt, inner ysep=6pt, text width=\linewidth-28pt] {
\textbf{{Bước 2 - Xử lý hạng mục $c$ \hfill $\mathrm{sup}(c) = 3 \geq \sigma$}}
 
\smallskip
\textbf{Ghi nhận:} $\{c\} \in \mathcal{F}$ với support = 3.
 
\smallskip
$A[c]$ chứa 2 mắt xích với hậu tố $(d,b)$ từ $T_4$ và $(d,b)$ từ reassign $T_1$.

\smallskip
\textbf{Xây dựng $\mathit{DB}_c$}:
 
\medskip
\begin{center}
\begin{tabular}{clll}
  \toprule
  Nguồn & Hậu tố & Ẩn h.mục đầu & Đưa vào $\mathit{DB}_c$ \\
  \midrule
  $T_4$ & $(d,b)$ & $j{=}d$; còn $(b)$ & Mắt xích $\to B_c[d]$; $B_c[d].cnt\mathrel{+}=1$ \\
  $T_1$ & $(d,b)$ & $j{=}d$; còn $(b)$ & Mắt xích $\to B_c[d]$; $B_c[d].cnt\mathrel{+}=1$ \\
  \bottomrule
\end{tabular}
\end{center}
 
\medskip
\textbf{Đếm trong $\mathit{DB}_c$:}
$\mathrm{sup}_{DB_c}(d) = 2$,\quad $\mathrm{sup}_{DB_c}(b) = 0$.
Không có hạng mục nào đạt $\sigma = 3$ $\Rightarrow$ không sinh tập với tiền tố $\{c\}$.
 
Kiểm tra: $\mathrm{sup}(\{c,b\}) = |\{T_1, T_4\}| = 2 < \sigma$ \checkmark{} nhất quán.
 
\medskip
\textbf{Phân bổ lại} 2 mắt xích từ $A[c]$:
\begin{itemize}[topsep=2pt, itemsep=1pt]
  \item $T_4$: $(d,b)$, $j{=}d$, còn $(b)$ $\to$ chèn vào $A[d]$; $A[d].cnt\mathrel{+}=1$.
  \item $T_1$: $(d,b)$, $j{=}d$, còn $(b)$ $\to$ chèn vào $A[d]$; $A[d].cnt\mathrel{+}=1$.
\end{itemize}
Sau reassign: $A[d].cnt = 1 + 2 = 3$, danh sách: $(b)[T_5] \to (b)[T_4] \to (b)[T_1]$.
};
\end{tikzpicture}
 
\medskip

\begin{table}[H]
  \centering
  \caption{Mảng $A$ sau Bước 2 (Ví dụ 2)}
  \label{tab:ex2-after2}
  \begin{tabular}{ccl}
    \toprule
    \textbf{H.mục} & $\mathbf{cnt}$ & \textbf{Danh sách} \\
    \midrule
    $a$ & 0 \textit{(đã xử lý)} & \textit{rỗng} \\
    $c$ & 0 \textit{(đã xử lý)} & \textit{rỗng} \\
    $d$ & 3 & $(b)[T_5] \to (b)[T_4] \to (b)[T_1]$ \\
    $b$ & 1 & \textit{rỗng (counter-only từ $T_2$)} \\
    \bottomrule
  \end{tabular}
\end{table}

\medskip
% -- EX2 BƯỚC 3 --
\begin{tikzpicture}
\node[ rounded corners=4pt,
      inner xsep=10pt, inner ysep=6pt, text width=\linewidth-28pt] {
\textbf{{Bước 3 - Xử lý hạng mục $d$ \hfill $\mathrm{sup}(d) = 3 \geq \sigma$}}
 
\smallskip
\textbf{Ghi nhận:} $\{d\} \in \mathcal{F}$ với support = 3.
 
\smallskip
$A[d]$ chứa 3 mắt xích, mỗi mắt xích có hậu tố $(b)$:
 
\medskip
\textbf{Xây dựng $\mathit{DB}_d$}:
 
\medskip
\begin{center}
\begin{tabular}{clll}
  \toprule
  Nguồn & Hậu tố & Ẩn h.mục đầu & Đưa vào $\mathit{DB}_d$ \\
  \midrule
  $T_5$ & $(b)$ & $j{=}b$; còn rỗng & \textbf{Counter-Only}: $B_d[b].cnt\mathrel{+}=1$ \\
  $T_4$ & $(b)$ & $j{=}b$; còn rỗng & \textbf{Counter-Only}: $B_d[b].cnt\mathrel{+}=1$ \\
  $T_1$ & $(b)$ & $j{=}b$; còn rỗng & \textbf{Counter-Only}: $B_d[b].cnt\mathrel{+}=1$ \\
  \bottomrule
\end{tabular}
\end{center}
 
\medskip
\textbf{Đếm trong $\mathit{DB}_d$:}
$\mathrm{sup}_{DB_d}(b) = 3 \geq \sigma$
$\Rightarrow$ \textbf{Ghi nhận} $\{d, b\} \in \mathcal{F}$ với support = 3.
 
\smallskip
Đệ quy $\mathit{DB}_{db}$: rỗng (toàn bộ Counter-Only, không có mắt xích nào) $\Rightarrow$ dừng.
 
\medskip
\textit{\textbf{Đây là điểm mấu chốt:} tập $\{d,b\}$ được phát hiện nhờ tích lũy
\textbf{ba lần Counter-Only liên tiếp}. Không có mắt xích nào được cấp phát trong
$\mathit{DB}_d$ -- toàn bộ thông tin support được truyền qua biến đếm $B_d[b].cnt$.
Nếu không có tối ưu hóa Counter-Only, ba mắt xích thừa sẽ bị cấp phát rồi ngay lập
tức bị bỏ đi vì hậu tố rỗng, gây lãng phí bộ nhớ.}
 
\medskip
\textbf{Phân bổ lại:} 3 mắt xích, mỗi mắt xích $j{=}b$, còn rỗng $\to$ $A[b].cnt\mathrel{+}=3$.
};
\end{tikzpicture}
 
\medskip

\begin{table}[H]
  \centering
  \caption{Mảng $A$ sau Bước 3 (Ví dụ 2)}
  \label{tab:ex2-after3}
  \begin{tabular}{ccl}
    \toprule
    \textbf{H.mục} & $\mathbf{cnt}$ & \textbf{Danh sách} \\
    \midrule
    $a$ & 0 \textit{(đã xử lý)} & \textit{rỗng} \\
    $c$ & 0 \textit{(đã xử lý)} & \textit{rỗng} \\
    $d$ & 0 \textit{(đã xử lý)} & \textit{rỗng} \\
    $b$ & 4 & \textit{rỗng (counter-only: 1 từ $T_2$ + 3 từ reassign $d$)} \\
    \bottomrule
  \end{tabular}
\end{table}

\medskip
% -- EX2 BƯỚC 4 --
\begin{tikzpicture}
\node[rounded corners=4pt,
      inner xsep=10pt, inner ysep=6pt, text width=\linewidth-28pt] {
\textbf{Bước 4 - Xử lý hạng mục $b$ \hfill $\mathrm{sup}(b) = 4 \geq \sigma$}
 
\smallskip
\textbf{Ghi nhận:} $\{b\} \in \mathcal{F}$ với support = 4.

\smallskip
\textbf{Nguồn gốc $A[b].cnt = 4$:} $0$ (init) $+ 1$ (reassign $T_2$ ở Bước 1)
$+ 3$ (reassign $T_5, T_4, T_1$ ở Bước 3) $= 4$.

\smallskip
Không có hạng mục nào đứng sau $b$ trong thứ tự. \textbf{Thuật toán kết thúc.}
};
\end{tikzpicture}
 
\subsection{Tập kết quả}
 
\begin{center}
\begin{tikzpicture}
\node[ rounded corners=5pt,
      inner xsep=12pt, inner ysep=10pt, font=\small] {
$\mathcal{F} = \bigl\{
  \{a\}{:}3,\;\; \{c\}{:}3,\;\; \{d\}{:}3,\;\; \{b\}{:}4,\;\; \{d,b\}{:}3
\bigr\}$
};
\end{tikzpicture}
\end{center}
 
\subsection{Kiểm tra chéo}
 
\begin{table}[H]
  \centering
  \caption{Kiểm tra chéo toàn bộ tập hạng mục (Ví dụ 2, $\sigma = 3$)}
  \label{tab:ex2-crosscheck}
  \begin{tabular}{lcc@{\hskip 1.5em}lcc}
    \toprule
    \textbf{Tập} & $\mathbf{sup}$ & \textbf{Kết luận}
      & \textbf{Tập} & $\mathbf{sup}$ & \textbf{Kết luận} \\
    \midrule
    \multicolumn{6}{l}{\textit{\textbf{1-tập (4 tập):}}} \\
    $\{a\}$   & 3 & {\textbf{Phổ biến}} & $\{d\}$   & 3 & {\textbf{Phổ biến}} \\
    $\{c\}$   & 3 & {\textbf{Phổ biến}} & $\{b\}$   & 4 & {\textbf{Phổ biến}} \\
    \midrule
    \multicolumn{6}{l}{\textit{\textbf{2-tập ($\binom{4}{2} = 6$ tập):}}} \\
    $\{a,c\}$  & 2 & {Loại} & $\{c,d\}$  & 2 & {Loại} \\
    $\{a,d\}$  & 1 & {Loại} & $\{c,b\}$  & 2 & {Loại} \\
    $\{a,b\}$  & 2 & {Loại} & $\{d,b\}$  & 3 & {\textbf{Phổ biến}} \\
    \midrule
    \multicolumn{6}{l}{\textit{\textbf{3-tập ($\binom{4}{3} = 4$ tập):}}} \\
    $\{a,c,d\}$  & 1 & {Loại}   & $\{a,c,b\}$  & 1 & {Loại} \\
    $\{a,d,b\}$  & 1 & {Loại}   & $\{c,d,b\}$  & 2 & {Loại} \\
    \midrule
    \multicolumn{6}{l}{\textit{\textbf{4-tập ($\binom{4}{4} = 1$ tập):}}} \\
    $\{a,c,d,b\}$ & 1 & {Loại} &               &   &  \\
    \bottomrule
  \end{tabular}
\end{table}
 
Kết quả kiểm tra chéo xác nhận tính đúng đắn. Đặc biệt, $\{d,b\}$ với support 3
chỉ được phát hiện nhờ tích lũy ba lần Counter-Only tại $\mathit{DB}_d[b]$:
mỗi mắt xích trong $A[d]$ có hậu tố chỉ còn 1 item $(b)$, nên khi xây dựng
$\mathit{DB}_d$, thuật toán chỉ tăng biến đếm mà không cấp phát mắt xích nào.
Nếu không có tối ưu hóa Counter-Only, ba mắt xích sẽ bị cấp phát rồi ngay lập tức
bị bỏ đi, gây lãng phí bộ nhớ -- đặc biệt nghiêm trọng ở các tầng đệ quy sâu
nơi giao dịch đã bị thu ngắn đáng kể.
 

% ============================================================
% CHƯƠNG 3: CÀI ĐẶT
% Người thực hiện: Mỹ Linh
% Ngôn ngữ: Julia (nhận +5 điểm thưởng theo quy định đề bài)
% ============================================================
\chapter{CÀI ĐẶT}

Chương này trình bày quá trình cài đặt thuật toán Relim từ đầu bằng ngôn ngữ
Julia, bám sát mã giả đã phân tích ở Chương~1. Phần cài đặt được tổ chức thành
bốn cấp độ (level) theo yêu cầu của đề bài: (1)~thuật toán cơ bản, (2)~kiểm thử
tự động, (3)~tối ưu hóa có đo lường, và (4)~xử lý vào/ra chuẩn SPMF kèm giao
diện dòng lệnh.

% ------------------------------------------------------------
\section{Môi trường và công cụ}
% ------------------------------------------------------------

Toàn bộ mã nguồn được phát triển và kiểm thử trên cấu hình ở
Bảng~\ref{tab:env}. Nhóm chọn Julia vì ngôn ngữ này vừa cho phép viết mã ở mức
trừu tượng cao, vừa biên dịch JIT (Just-In-Time) sang mã máy nên đạt hiệu năng
gần với C trong các vòng lặp nóng (hot loop) - phù hợp với một thuật toán khai
phá dữ liệu thiên về thao tác con trỏ và đệ quy như Relim.

\begin{table}[H]
    \centering
    \caption{Môi trường thực nghiệm}
    \label{tab:env}
    \begin{tabularx}{\textwidth}{|l|X|}
        \hline
        \textbf{Thành phần} & \textbf{Chi tiết} \\
        \hline
        Ngôn ngữ        & Julia 1.12.5 (yêu cầu đề bài: $\geq 1.9$) \\
        \hline
        Hệ điều hành    & Linux Mint 22.3 (nhân Linux 6.17) \\
        \hline
        CPU             & Intel Core i5-11320H @ 3.20\,GHz \\
        \hline
        RAM             & 16\,GB \\
        \hline
        Thư viện        & \texttt{BenchmarkTools.jl} (đo thời gian/bộ nhớ),
                          \texttt{Test}, \texttt{Random} (thư viện chuẩn) \\
        \hline
        Quản lý gói     & \texttt{Project.toml} (tương đương
                          \texttt{requirements.txt} của Python) \\
        \hline
    \end{tabularx}
\end{table}

Theo đúng quy định của đề bài, nhóm \textbf{không sử dụng bất kỳ thư viện khai
phá tập phổ biến có sẵn nào}; mọi cấu trúc dữ liệu và logic của thuật toán đều
được tự cài đặt dựa trên bài báo gốc của Borgelt~\cite{borgelt2005}.

% ------------------------------------------------------------
\section{Kiến trúc mã nguồn}
% ------------------------------------------------------------

Mã nguồn được tổ chức theo cấu trúc thư mục gợi ý trong đề bài, tách bạch giữa
định nghĩa cấu trúc dữ liệu, thuật toán, tiện ích vào/ra và bộ kiểm thử:

\begin{lstlisting}[style=julia, caption={Cấu trúc thư mục mã nguồn},
                   label={lst:tree}]
relim_datamining_julia/
|-- Project.toml             # Dac ta dependencies (Julia)
|-- README.md                # Huong dan cai dat va su dung
|-- src/
|   |-- structures.jl        # Cau truc du lieu: mang A + linked list
|   |-- utils.jl             # I/O dinh dang SPMF + parse tham so dong lenh
|   |-- main.jl              # Entry point (CLI)
|   `-- algorithm/
|       |-- relim.jl         # Thuat toan Relim (ham chinh + ham con)
|       `-- visualize.jl     # Minh hoa tung buoc (Chuong 2)
|-- tests/
|   |-- test_unit.jl         # Unit test (Chuong 3): doi chung brute-force
|   |-- bench_opt.jl         # Do L1 vs L3 (thoi gian + bo nho)
|   |-- test_correctness.jl  # Correctness vs SPMF (Chuong 4)
|   `-- test_benchmark.jl    # Benchmark + so sanh SPMF (Chuong 4)
`-- data/
    |-- toy/                 # CSDL nho cho vi du tay
    |-- benchmark/           # CSDL benchmark (chess, mushroom, retail, ...)
    `-- real_world/          # CSDL thuc te (Market Basket)
\end{lstlisting}

\subsection{Mô tả các module chính}

\begin{itemize}
    \item \texttt{structures.jl}: định nghĩa ba kiểu dữ liệu cốt lõi của Relim
    là \texttt{ListNode} (mắt xích danh sách liên kết), \texttt{RelimEntry}
    (phần tử $A[i]$ gồm bộ đếm và con trỏ đầu danh sách) và \texttt{RelimArray}
    (mảng $A$). Các kiểu này hiện thực hóa trực tiếp cấu trúc ``mảng danh sách
    liên kết đơn'' đã mô tả ở Mục~\ref{subsec:ds} và minh họa ở
    Hình~\ref{fig:relim-structure}.

    \item \texttt{relim.jl}: chứa hàm public \texttt{relim} cùng ba hàm con
    \texttt{recursive\_relim!}, \texttt{build\_conditional} và
    \texttt{reassign!}, ánh xạ một-một với ba thủ tục mã giả
    \textsc{Relim}, \textsc{Recursive-Relim} và \textsc{Build-Conditional-DB}
    (Thuật toán~\ref{alg:relim-main}, \ref{alg:relim-recursive},
    \ref{alg:relim-build-cond}).

    \item \texttt{utils.jl}: tiện ích đọc/ghi file định dạng SPMF
    (\texttt{load\_spmf}, \texttt{save\_spmf}), hàm đếm support phục vụ kiểm thử
    (\texttt{count\_support}) và bộ phân tích tham số dòng lệnh
    (\texttt{parse\_cli}).

    \item \texttt{main.jl}: điểm vào của chương trình, kết nối các module trên
    thành một công cụ dòng lệnh hoàn chỉnh.
\end{itemize}

\section{Cài đặt Level 1: Thuật toán cơ bản}


\subsection{Cấu trúc dữ liệu}

Cấu trúc dữ liệu được định nghĩa trong \texttt{structures.jl}. Mỗi giao dịch
được lưu một lần dưới dạng \texttt{Vector\{Int\}} và được \textbf{chia sẻ} (share)
giữa các mắt xích thông qua trường \texttt{trans} cùng con trỏ vị trí
\texttt{trans\_ptr}; nhờ vậy việc ``ẩn'' một hạng mục dẫn đầu chỉ là phép tăng
con trỏ chứ không phải sao chép mảng.

\begin{lstlisting}[style=julia, caption={Định nghĩa cấu trúc dữ liệu (structures.jl)},
                   label={lst:structures}]
mutable struct ListNode
    trans::Vector{Int}             # tham chieu (share) toi giao dich goc
    trans_ptr::Int                 # offset 1-based cua item ke tiep
    next::Union{ListNode,Nothing}  # con tro toi mat xich ke tiep
end

mutable struct RelimEntry
    cnt::Int                       # so giao dich co item nay lam leading item
    head::Union{ListNode,Nothing}  # con tro dau danh sach lien ket
end

RelimEntry() = RelimEntry(0, nothing)

const RelimArray = Vector{RelimEntry}
new_relim_array(k::Int)::RelimArray = [RelimEntry() for _ in 1:k]
\end{lstlisting}

\subsection{Hàm chính và tiền xử lý}

Hàm \texttt{relim} (Listing~\ref{lst:relim-main}) hiện thực hóa thủ tục
\textsc{Relim} (Thuật toán~\ref{alg:relim-main}). Năm bước trong mã nguồn được
chú thích tương ứng với năm bước trong mã giả: (1)~đếm support một lần quét,
(2)~lọc hạng mục phổ biến, (3)~ánh xạ lại (remap) hạng mục về chỉ số $1..k$ theo
thứ tự support \emph{tăng dần} và sắp các hạng mục trong từng giao dịch,
(4)~dựng mảng $A$ ban đầu, và (5)~gọi đệ quy.

Việc ánh xạ hạng mục gốc sang chỉ số liên tục $1..k$ cho phép dùng trực tiếp
mảng (\texttt{Vector}) thay cho từ điển (\texttt{Dict}) trong vòng đệ quy, giúp
truy cập $A[i]$ với chi phí $O(1)$ và thân thiện với bộ nhớ đệm (cache).

\begin{lstlisting}[style=julia, caption={Hàm chính \texttt{relim} (relim.jl)},
                   label={lst:relim-main}]
function relim(transactions::Vector{Vector{Int}}, minsup::Real;
               relative::Bool=true,
               counter_only::Bool=true)::Vector{Tuple{Vector{Int},Int}}

    n = length(transactions)
    minsup_abs = relative ? max(1, ceil(Int, minsup * n)) : Int(minsup)

    # Buoc 1: dem support cua tung item (1 lan quet)
    counts = Dict{Int,Int}()
    for t in transactions, x in t
        counts[x] = get(counts, x, 0) + 1
    end

    # Buoc 2: loc item pho bien
    freq_items = [it for (it, c) in counts if c >= minsup_abs]
    isempty(freq_items) && return Tuple{Vector{Int},Int}[]

    # Remap: item pho bien -> chi so 1..k, sap theo support TANG DAN
    sort!(freq_items; by = it -> (counts[it], it))
    k = length(freq_items)
    item2idx = Dict{Int,Int}(it => i for (i, it) in enumerate(freq_items))
    idx2item = freq_items

    # Buoc 3: voi moi giao dich, giu lai item pho bien va sort theo index
    sorted_txns = Vector{Vector{Int}}()
    sizehint!(sorted_txns, n)
    @inbounds for t in transactions
        filtered = Int[]
        sizehint!(filtered, length(t))
        for x in t
            idx = get(item2idx, x, 0)
            idx == 0 || push!(filtered, idx)
        end
        if !isempty(filtered)
            sort!(filtered)
            push!(sorted_txns, filtered)
        end
    end

    # Buoc 4: dung mang A ban dau
    A = new_relim_array(k)
    @inbounds for t in sorted_txns
        i = t[1]
        A[i].cnt += 1
        A[i].head = ListNode(t, 2, A[i].head)  # "an" item dau = tang con tro
    end

    # Buoc 5: de quy
    F = Vector{Tuple{Vector{Int},Int}}()
    prefix = Int[]
    recursive_relim!(A, prefix, minsup_abs, F, idx2item, counter_only)
    return F
end
\end{lstlisting}

\subsection{Hàm đệ quy chính}

Hàm \texttt{recursive\_relim!} (Listing~\ref{lst:relim-rec}) là phần lõi của
thuật toán, tương ứng Thuật toán~\ref{alg:relim-recursive}. Với mỗi hạng mục $i$
(duyệt từ hạng mục ít phổ biến nhất): nếu đủ support thì xuất tập phổ biến
$P \cup \{i\}$, dựng cơ sở dữ liệu điều kiện rồi đệ quy; sau đó thực hiện bước
\emph{loại bỏ} (elimination) - chuyển toàn bộ mắt xích của $A[i]$ sang các danh
sách tương ứng với hạng mục dẫn đầu mới. Dấu \texttt{!} ở cuối tên hàm là quy
ước của Julia cho hàm có biến đổi (mutate) đối số.

\begin{lstlisting}[style=julia, caption={Hàm đệ quy \texttt{recursive\_relim!} (relim.jl)},
                   label={lst:relim-rec}]
function recursive_relim!(A::RelimArray, prefix::Vector{Int}, smin::Int,
                          F::Vector{Tuple{Vector{Int},Int}},
                          idx2item::Vector{Int}, counter_only::Bool)
    k = length(A)
    @inbounds for i in 1:k
        c = A[i].cnt
        if c >= smin
            push!(prefix, idx2item[i])
            push!(F, (sort(copy(prefix)), c))           # OUTPUT P u {i}
            A_cond = build_conditional(A[i], k, counter_only)
            recursive_relim!(A_cond, prefix, smin, F, idx2item, counter_only)
            pop!(prefix)
        end
        reassign!(A[i], A)   # ELIMINATION: loai item i khoi CSDL hien tai
    end
end
\end{lstlisting}

Hai hàm con còn lại đảm nhiệm việc dựng cơ sở dữ liệu điều kiện và bước loại bỏ
(Listing~\ref{lst:relim-helpers}). Trong \texttt{build\_conditional}, mỗi giao
dịch đã ẩn hạng mục dẫn đầu được phân loại theo hạng mục kế tiếp $j$. Khối lệnh
\texttt{if p + 1 <= L} chính là vị trí áp dụng \textbf{Counter-Only
Optimization} sẽ trình bày ở Level~3: khi giao dịch chỉ còn đúng một hạng mục,
ta chỉ tăng bộ đếm mà không cấp phát mắt xích mới.

\begin{lstlisting}[style=julia, caption={Hàm \texttt{build\_conditional} và \texttt{reassign!} (relim.jl)},
                   label={lst:relim-helpers}]
function build_conditional(entry::RelimEntry, k_cond::Int,
                           counter_only::Bool)::RelimArray
    A_cond = new_relim_array(k_cond)
    node = entry.head
    @inbounds while node !== nothing
        t = node.trans; p = node.trans_ptr; L = length(t)
        if p <= L
            j = t[p]                       # item dan dau moi (j > i)
            if j <= k_cond
                A_cond[j].cnt += 1
                if p + 1 <= L
                    A_cond[j].head = ListNode(t, p + 1, A_cond[j].head)
                elseif !counter_only       # ban L1: van cap node "rong"
                    A_cond[j].head = ListNode(t, p + 1, A_cond[j].head)
                end
            end
        end
        node = node.next
    end
    return A_cond
end
 
\begin{remark}
\textbf{Sự khác biệt giữa Giả mã và Cài đặt:} Trong Thuật toán 2 (Chương 1),
giả mã sử dụng thao tác $\mathrm{Copy}(e)$ để tạo bản sao hoàn toàn của mắt xích.
Tuy nhiên, trong cài đặt Julia, do đặc thù quản lý bộ nhớ, việc sao chép mảng liên tục rất
tốn kém. Do đó, hàm \texttt{build\_conditional} khởi tạo một \texttt{ListNode} mới nhưng
vẫn \textbf{tham chiếu (share)} lại mảng \texttt{t} của giao dịch gốc và chỉ tịnh tiến
con trỏ lên \texttt{p + 1}. Cách tiếp cận này đảm bảo tính tương đương logic với giả mã
nhưng loại bỏ hoàn toàn chi phí sao chép bộ nhớ thừa.
\end{remark}
function reassign!(entry::RelimEntry, A::RelimArray)
    node = entry.head
    @inbounds while node !== nothing
        nxt = node.next
        t = node.trans; p = node.trans_ptr; L = length(t)
        if p <= L
            j = t[p]
            node.trans_ptr = p + 1         # tai su dung node, khong cap phat moi
            node.next = A[j].head
            A[j].head = node
            A[j].cnt += 1
        end
        node = nxt
    end
    entry.head = nothing
    entry.cnt = 0
    return entry
end
\end{lstlisting}

\subsection{Kiểm chứng nhanh trên ví dụ}

Chạy thuật toán trên cơ sở dữ liệu đồ chơi
$\mathcal{D} = \{\{1,3,4\}, \{2,3,5\}, \{1,2,3,5\}, \{2,5\}\}$ với
$minsup = 0{,}5$ (tương đương support tuyệt đối $\geq 2$) cho ra đúng 9 tập phổ
biến, khớp hoàn toàn với kết quả liệt kê thủ công:
$\{1\}, \{2\}, \{3\}, \{5\}, \{1,3\}, \{2,3\}, \{2,5\}, \{3,5\}, \{2,3,5\}$.


\section{Cài đặt Level 2: Kiểm thử tự động}

\subsection{Thiết kế test case}

Bộ kiểm thử \texttt{tests/test\_unit.jl} dùng thư viện \texttt{Test} của
Julia, đối chứng kết quả của Relim với một hàm \emph{brute-force} (liệt kê toàn
bộ tập con và đếm support trực tiếp) đóng vai trò chuẩn vàng (ground truth). Cách
kiểm thử này độc lập hoàn toàn với dữ liệu và công cụ bên ngoài. Các
nhóm test bao gồm:

\begin{enumerate}[label=(\arabic*)]
    \item \textbf{Toy database từ paper}: cơ sở dữ liệu 4 giao dịch ở trên,
    kiểm tra số lượng và đối chứng brute-force.
    \item \textbf{Trường hợp biên}: cơ sở dữ liệu rỗng; ngưỡng minsup lớn hơn
    số giao dịch (không có hạng mục phổ biến).
    \item \textbf{Giao dịch đơn hạng mục}: kiểm tra việc đếm support cơ bản.
    \item \textbf{Cơ sở dữ liệu dày (dense)}: mọi hạng mục đều phổ biến, sinh
    đầy đủ $2^k - 1$ tập phổ biến.
    \item \textbf{Bất biến của tối ưu hóa}: khẳng định hai phiên bản
    \texttt{counter\_only = true/false} cho kết quả y hệt nhau.
    \item \textbf{Đối chứng ngẫu nhiên}: sinh 10 cơ sở dữ liệu ngẫu nhiên (với
    \texttt{Random.seed!(42)} để tái lập được) và so sánh từng tập phổ biến với
    brute-force.
\end{enumerate}

\subsection{Kết quả kiểm thử}

Toàn bộ \textbf{18/18} khẳng định (assertion) đều đạt
(Bảng~\ref{tab:unittest}). Việc đối chứng với brute-force trên cả ví dụ cố định
lẫn 10 trường hợp ngẫu nhiên cho phép khẳng định tính đúng đắn của cài đặt một
cách độc lập với SPMF.

\begin{table}[H]
    \centering
    \caption{Kết quả chạy bộ kiểm thử \texttt{test\_unit.jl}}
    \label{tab:unittest}
    \begin{tabularx}{\textwidth}{|X|c|c|}
        \hline
        \textbf{Nhóm kiểm thử} & \textbf{Số assertion} & \textbf{Kết quả} \\
        \hline
        Toy database (số lượng + brute-force)   & 2  & Đạt \\
        \hline
        Trường hợp biên (rỗng / minsup lớn)     & 2  & Đạt \\
        \hline
        Giao dịch đơn hạng mục                  & 1  & Đạt \\
        \hline
        Cơ sở dữ liệu dày                       & 2  & Đạt \\
        \hline
        Bất biến Counter-Only ON/OFF            & 1  & Đạt \\
        \hline
        Đối chứng ngẫu nhiên (10 trial)         & 10 & Đạt \\
        \hline
        \textbf{Tổng}                           & \textbf{18} & \textbf{18/18 Đạt} \\
        \hline
    \end{tabularx}
\end{table}

\begin{remark}
Việc \textbf{so sánh số lượng và support của từng itemset với SPMF} trên các bộ
dữ liệu benchmark (Chess, Mushroom, Retail, T10I4D100K) được trình bày ở
Chương~4 (Thực nghiệm) qua bộ \texttt{test\_correctness.jl} và
\texttt{test\_benchmark.jl}, do phần đó cần chạy SPMF làm chuẩn so sánh và thuộc
nhiệm vụ thực nghiệm của nhóm. Bộ kiểm thử ở chương này
(\texttt{test\_unit.jl}) tập trung kiểm chứng tính đúng đắn của \emph{cài đặt}
một cách độc lập, không phụ thuộc dữ liệu ngoài.
\end{remark}


\section{Cài đặt Level 3: Tối ưu hóa}

\subsection{Mô tả kỹ thuật tối ưu hóa}

Cài đặt áp dụng hai kỹ thuật tối ưu được mô tả trong chính bài báo gốc
của Borgelt~\cite{borgelt2005}:

\paragraph{Tối ưu 1: Pointer increment (tránh sao chép giao dịch).}
Thay vì cắt mảng (slice) để loại hạng mục dẫn đầu khi đi sâu vào đệ quy, mỗi mắt
xích chỉ lưu một con trỏ \texttt{trans\_ptr} vào giao dịch gốc dùng chung. Việc
``ẩn'' một hạng mục trở thành phép tăng con trỏ $O(1)$, không sinh cấp phát bộ
nhớ. Hơn nữa, bước \texttt{reassign!} \emph{tái sử dụng} chính các mắt xích cũ
(chỉ đổi con trỏ liên kết) thay vì tạo mắt xích mới.

\paragraph{Tối ưu 2: Counter-Only Optimization.}
Khi một giao dịch trong cơ sở dữ liệu điều kiện chỉ còn lại đúng một hạng mục
$j$ sau khi ẩn hạng mục dẫn đầu, việc tạo một mắt xích danh sách cho nó là vô
ích (nó sẽ không bao giờ sinh thêm cơ sở dữ liệu điều kiện con nào nữa). Khi đó
ta chỉ tăng \texttt{A\_cond[j].cnt} mà không cấp phát mắt xích
(Listing~\ref{lst:relim-helpers}). Tối ưu này loại bỏ phần lớn cấp phát bộ nhớ
động ở các tầng đệ quy sâu, nơi giao dịch đã bị thu ngắn đáng kể.

Cài đặt giữ lại tham số \texttt{counter\_only} để có thể bật/tắt Tối ưu~2, nhờ
đó đo lường trực tiếp được mức cải thiện giữa \emph{bản cơ bản} (Level~1,
\texttt{counter\_only = false}) và \emph{bản tối ưu} (Level~3,
\texttt{counter\_only = true}) trên cùng một mã nguồn. Bộ kiểm thử ở Level~2 đã
khẳng định hai phiên bản cho kết quả giống hệt nhau, đảm bảo tối ưu không làm
thay đổi tính đúng đắn.

\subsection{Đo lường hiệu quả tối ưu hóa}

Phương pháp đo: dùng \texttt{BenchmarkTools.@belapsed} (lấy thời gian nhỏ nhất
qua nhiều mẫu nên ổn định, đơn vị ms) và \texttt{@allocated} (tổng bộ nhớ cấp
phát, đơn vị MB) cho cả hai phiên bản trên cùng một bộ dữ liệu và cùng ngưỡng
minsup, sau bước khởi động (warm-up) để loại bỏ chi phí biên dịch JIT. Toàn bộ
quy trình đo được đóng gói trong \texttt{tests/bench\_opt.jl} để tái lập. Phép
đo thực hiện trên hai bộ dữ liệu \emph{dày} (dense) -- nơi Counter-Only phát huy
rõ rệt nhất -- với kết quả ở Bảng~\ref{tab:opt}.

\begin{table}[H]
    \centering
    \caption{So sánh hiệu năng bản cơ bản (Level~1, \texttt{counter\_only=false})
             và bản tối ưu (Level~3, \texttt{counter\_only=true})}
    \label{tab:opt}
    \begin{tabularx}{\textwidth}{|X|r|r|r|r|r|r|}
        \hline
        \multirow{2}{*}{\textbf{Bộ dữ liệu}} & \multirow{2}{*}{\textbf{$|F|$}}
            & \multicolumn{2}{c|}{\textbf{Thời gian (ms)}}
            & \multicolumn{2}{c|}{\textbf{Bộ nhớ (MB)}}
            & \multirow{2}{*}{\textbf{Giảm RAM}} \\
        \cline{3-6}
            & & Cơ bản & Tối ưu & Cơ bản & Tối ưu & \\
        \hline
        Mushroom (minsup $=0{,}30$) & 2.587 & 98{,}0  & 72{,}3  & 127{,}8 & 68{,}0  & 46{,}8\% \\
        \hline
        Chess (minsup $=0{,}80$)    & 8.227 & 285{,}9 & 194{,}4 & 347{,}5 & 181{,}5 & 47{,}8\% \\
        \hline
    \end{tabularx}
\end{table}

\paragraph{Nhận xét.}
Counter-Only Optimization mang lại cải thiện \emph{nhất quán} trên cả hai bộ dữ
liệu: bộ nhớ cấp phát giảm gần một nửa (46,8\% với Mushroom và 47,8\% với Chess),
kéo theo thời gian chạy nhanh hơn $1{,}36$-$1{,}47$ lần. Mức cải thiện về thời
gian thấp hơn mức cải thiện về bộ nhớ là hợp lý: tối ưu này chủ yếu cắt giảm số
lần \emph{cấp phát} mắt xích (allocation) và áp lực lên bộ thu gom rác (garbage
collector), trong khi tổng số phép duyệt để đếm support gần như không đổi. Kết
quả khớp với phân tích lý thuyết ở Mục~\ref{subsec:ds}: phần lớn giao dịch ở các
tầng đệ quy sâu đã bị thu ngắn còn một hạng mục - đúng trường hợp mà Counter-Only
tránh được việc cấp phát.

\begin{remark}
Hai bộ dữ liệu thưa (sparse) như Retail và T10I4D100K có lợi ích từ Counter-Only
nhỏ hơn, vì giao dịch ngắn khiến số mắt xích sinh ra ở các tầng sâu vốn đã ít.
Các thực nghiệm về thời gian và bộ nhớ theo nhiều mức minsup trên toàn bộ năm bộ
dữ liệu benchmark được trình bày chi tiết ở Chương~4.
\end{remark}


\section{Cài đặt Level 4: Xử lý I/O chuẩn SPMF}

\subsection{Hàm đọc/ghi file (utils.jl)}

Định dạng SPMF quy ước mỗi dòng của tệp đầu vào là một giao dịch, các hạng mục
là số nguyên cách nhau bởi khoảng trắng. Tệp đầu ra liệt kê mỗi tập phổ biến
trên một dòng theo dạng \texttt{i1 i2 ... ik \#SUP: count}. Hai hàm
\texttt{load\_spmf} và \texttt{save\_spmf} (Listing~\ref{lst:io}) hiện thực hóa
đúng quy ước này, giúp kết quả có thể đối chiếu (diff) trực tiếp với đầu ra của
SPMF.

\begin{lstlisting}[style=julia, caption={Đọc/ghi định dạng SPMF (utils.jl)},
                   label={lst:io}]
function load_spmf(path::AbstractString)::Vector{Vector{Int}}
    txns = Vector{Vector{Int}}()
    open(path) do io
        for line in eachline(io)
            s = strip(line)
            isempty(s) && continue
            push!(txns, parse.(Int, split(s)))
        end
    end
    return txns
end

function save_spmf(results::Vector{Tuple{Vector{Int},Int}}, path::AbstractString)
    open(path, "w") do io
        for (iset, sup) in results
            print(io, join(sort(iset), ' '))
            println(io, " #SUP: ", sup)
        end
    end
    return path
end
\end{lstlisting}

\subsection{Giao diện dòng lệnh (main.jl)}

Chương trình hỗ trợ các tham số dòng lệnh theo đúng yêu cầu đề bài: đường dẫn
tệp vào/ra và ngưỡng minsup. Đáng chú ý, minsup được hỗ trợ \textbf{cả hai cách
hiểu}: mặc định là tỉ lệ tương đối trong $[0,1]$ (giống giao diện SPMF), hoặc số
tuyệt đối khi bật cờ \texttt{--absolute} (giống mã giả trong bài báo). Bộ phân
tích tham số được viết thủ công, không phụ thuộc thư viện ngoài.

\begin{lstlisting}[style=julia, caption={Điểm vào chương trình (main.jl)},
                   label={lst:main}]
include(joinpath(@__DIR__, "algorithm", "relim.jl"))

function main()
    opts = parse_cli(copy(ARGS))
    txns = load_spmf(opts.input)
    t = @elapsed result = relim(txns, opts.minsup; relative = !opts.absolute)
    save_spmf(result, opts.output)
    println("# frequent itemsets: ", length(result))
    println("Runtime (s)  : ", round(t; digits = 4))
end

if abspath(PROGRAM_FILE) == @__FILE__
    main()
end
\end{lstlisting}

Ví dụ lời gọi (đầy đủ hướng dẫn trong \texttt{README.md}):

\begin{lstlisting}[style=julia, caption={Ví dụ chạy từ dòng lệnh},
                   label={lst:cli-example}]
julia --project=. src/main.jl --input data/toy/example_basic.txt \
                              --minsup 0.5 --output result.txt
\end{lstlisting}

Lệnh trên đọc cơ sở dữ liệu đồ chơi, khai phá với $minsup = 50\%$ và ghi 9 tập
phổ biến ra \texttt{result.txt} theo đúng định dạng SPMF.


\chapter{THỰC NGHIỆM VÀ ĐÁNH GIÁ}
 
%% Người thực hiện phần thực nghiệm: (xem Plan.docx - phần 3.4.1)
%% Người phân tích kết quả: (xem Plan.docx - phần 3.4.2)
%% Lấy code từ Level 1-3 ở Chương 3 để chạy
 
% ------------------------------------------------------------
\section{Tập dữ liệu benchmark}
% ------------------------------------------------------------
 
%% Mô tả các tập dữ liệu (tham khảo bảng trong đề bài, Section 3.4.1)
%% Nguồn: https://www.philippe-fournier-viger.com/spmf/index.php?link=datasets.php
 
\begin{table}[H]
    \centering
    \caption{Tóm tắt tập dữ liệu benchmark}
    \label{tab:datasets}
    \begin{tabularx}{\textwidth}{|l|r|r|r|X|}
    \hline
    \textbf{Tập dữ liệu} & \textbf{\#Trans.} & \textbf{\#Items} & \textbf{AvgLen} & \textbf{Đặc điểm} \\
    \hline
    Chess       & 3.196   & 75    & 37.0 & Dày đặc (dense), itemset dài \\
    \hline
    Mushroom    & 8.124   & 119   & 23.0 & Dày đặc, nhiều item \\
    \hline
    Retail      & 88.162  & 16.47 & 10.3 & Thưa (sparse), thực tế \\
    \hline
    Accidents   & 340.183 & 468   & 33.8 & Rất lớn, dày đặc \\
    \hline
    T10I4D100K  & 100.000 & 870   & 10.1 & Tổng hợp, thưa \\
    \hline
    \end{tabularx}
\end{table}
 
%% Bổ sung nhận xét về đặc điểm từng dataset liên quan đến RELIM:
%% Chess, Accidents → dense → dự đoán chậm hơn
%% Retail, T10I4D100K → sparse → dự đoán nhanh hơn
 
% ------------------------------------------------------------
\section{Thực nghiệm bắt buộc}
% ------------------------------------------------------------
 
\subsection{Kiểm tra tính đúng đắn (Correctness)}

Mục đích của thực nghiệm này là đảm bảo tính đúng đắn của thuật toán RELIM do nhóm
tự cài đặt bằng ngôn ngữ Julia thông qua việc so sánh đối chiếu chéo
(cross-validation) kết quả khai thác tập phổ biến với thư viện chuẩn SPMF --
một công cụ Data Mining mã nguồn mở bằng Java được công nhận rộng rãi trong giới
học thuật.

\paragraph{Thiết lập thực nghiệm.}
Thực nghiệm sử dụng cả 5 bộ dữ liệu chuẩn với hai ngưỡng \texttt{minsup} khác nhau
cho mỗi bộ. Hai chỉ số được đánh giá:
\begin{enumerate}[label=(\roman*)]
    \item \textit{Match Rate (Tỷ lệ khớp):} Tỷ lệ phần trăm số lượng tập phổ biến
          mà thuật toán của nhóm tìm được trùng khớp với SPMF.
    \item \textit{Support Correct:} Kiểm tra xem giá trị độ hỗ trợ đếm được của
          từng tập phổ biến có hoàn toàn chính xác so với kết quả của SPMF hay không.
\end{enumerate}

\paragraph{Kết quả thực nghiệm.}

\begin{table}[H]
    \centering
    \caption{Kết quả kiểm tra tính đúng đắn trên 5 bộ dữ liệu benchmark}
    \label{tab:correctness}
    \begin{tabularx}{\textwidth}{|l|c|c|c|c|>{\centering\arraybackslash}X|}
    \hline
    \textbf{Dataset} & \textbf{Loại} & \textbf{Minsup Test 1} & \textbf{Minsup Test 2} & \textbf{Match Rate} & \textbf{Support Correct} \\
    \hline
    Chess       & Dense  & 50\% (1\,598)   & 80\% (2\,557)   & 100{,}0\% & \checkmark \\
    \hline
    Mushroom    & Dense  & 10\% (812)      & 50\% (4\,062)   & 100{,}0\% & \checkmark \\
    \hline
    Accidents   & Dense  & 50\% (170\,091) & 75\% (255\,137) & 100{,}0\% & \checkmark \\
    \hline
    Retail      & Sparse & 1\% (882)       & 5\% (4\,408)    & 100{,}0\% & \checkmark \\
    \hline
    T10I4D100K  & Sparse & 1\% (1\,000)    & 5\% (5\,000)    & 100{,}0\% & \checkmark \\
    \hline
    \end{tabularx}
    \par\smallskip
    \footnotesize\textit{Ghi chú: Giá trị minsup tuyệt đối trong ngoặc được tính bằng
    cách làm tròn theo tỉ lệ phần trăm so với tổng số giao dịch của từng bộ dữ liệu.
    Dữ liệu đầu vào được làm sạch, loại bỏ nhiễu trước khi chạy.}
\end{table}

\paragraph{Kết luận.}
Kết quả trong Bảng~\ref{tab:correctness} cho thấy thuật toán RELIM Julia do nhóm
cài đặt đạt Match Rate \textbf{100\%} và Support Correct \textbf{tuyệt đối} trên
toàn bộ 5 bộ dữ liệu ở các ngưỡng \texttt{minsup} khác nhau. Điều này khẳng định:
\begin{enumerate}[label=(\arabic*)]
    \item Thuật toán \textbf{không bỏ sót} bất kỳ tập phổ biến nào so với chuẩn SPMF.
    \item Cấu trúc danh sách liên kết đơn và phép chiếu cơ sở dữ liệu đệ quy
          hoạt động hoàn hảo, không tạo ra sai số đếm (ghost items / duplicates).
    \item Độ tin cậy của cài đặt đủ điều kiện làm nền tảng cho việc tính toán
          Confidence, Lift và rút trích luật kết hợp ở Chương~5.
\end{enumerate}
 
\subsection{Thời gian chạy theo minsup}

Thực nghiệm đo thời gian thực thi của RELIM trên 5 bộ dữ liệu với các mức
\texttt{minsup} giảm dần. Kết quả được tổng hợp trong Bảng~\ref{tab:runtime}.

\begin{table}[H]
    \centering
    \caption{Thời gian chạy của RELIM theo ngưỡng \texttt{minsup}}
    \label{tab:runtime}
    \begin{tabular}{|l|l|r|r|}
    \hline
    \textbf{Dataset} & \textbf{Đặc tính} & \textbf{Minsup} & \textbf{Thời gian (ms)} \\
    \hline
    \multirow{5}{*}{\textbf{Chess}}
        & \multirow{5}{*}{Rất đặc, ngắn}
        & 70\% & 2.264{,}86 \\
        \cline{3-4}
        & & 60\% & 7.003{,}00 \\
        \cline{3-4}
        & & 55\% & 14.415{,}55 \\
        \cline{3-4}
        & & 50\% & 26.836{,}07 \\
        \cline{3-4}
        & & 45\% & 54.725{,}52 \\
    \hline
    \multirow{4}{*}{\textbf{Mushroom}}
        & \multirow{4}{*}{Đặc, ngắn}
        & 80\% & 4{,}97 \\
        \cline{3-4}
        & & 50\% & 22{,}32 \\
        \cline{3-4}
        & & 30\% & 100{,}80 \\
        \cline{3-4}
        & & 10\% & 7.870{,}64 \\
    \hline
    \multirow{3}{*}{\textbf{Accidents}}
        & \multirow{3}{*}{Thưa trung bình, lớn}
        & 75\% & 9.836{,}78 \\
        \cline{3-4}
        & & 60\% & 56.619{,}14 \\
        \cline{3-4}
        & & 50\% & 174.954{,}10 \\
    \hline
    \multirow{4}{*}{\textbf{Retail}}
        & \multirow{4}{*}{Thưa, trung bình}
        & 5\%  & 1.202{,}71 \\
        \cline{3-4}
        & & 2\%  & 1.075{,}07 \\
        \cline{3-4}
        & & 1\%  & 1.438{,}59 \\
        \cline{3-4}
        & & 0{,}5\% & 2.227{,}95 \\
    \hline
    \multirow{3}{*}{\textbf{T10I4D100K}}
        & \multirow{3}{*}{Thưa, cực lớn}
        & 5\%  & 18{,}91 \\
        \cline{3-4}
        & & 2\%  & 34{,}04 \\
        \cline{3-4}
        & & 1\%  & 58{,}04 \\
    \hline
    \end{tabular}
\end{table}

\paragraph{Nhận xét.}
Thời gian thực thi của RELIM tăng tỷ lệ nghịch với ngưỡng \texttt{minsup},
đặc biệt rõ rệt trên các tập dữ liệu dày đặc:
\begin{itemize}
    \item \textbf{Tập Dense (Chess, Mushroom, Accidents):} Khi giảm \texttt{minsup},
          số lượng itemset tăng bùng nổ kéo theo runtime tăng vọt theo hàm mũ.
          Điển hình tại \texttt{chess}, khi \texttt{minsup} giảm từ 60\% xuống 50\%,
          số itemset tăng gấp 5 lần (từ 254.944 lên 1.272.932), làm runtime nhảy
          từ 7.003 ms lên 26.836 ms. Ở mức 45\%, runtime đạt tới 54.725 ms
          ($\approx$55 giây). Tập \texttt{accidents} ở \texttt{minsup} = 50\% tốn
          đến $\approx$175 giây do kết hợp giữa kích thước lớn (340.183 giao dịch)
          và độ dài trung bình cao.
    \item \textbf{Tập Sparse (Retail, T10I4D100K):} Thời gian chạy rất nhỏ và ổn
          định. Tập \texttt{T10I4D100K} với 100.000 giao dịch chỉ mất 58{,}04 ms tại
          \texttt{minsup} 1\%. Đặc biệt, tập \texttt{Retail} cho thấy hiện tượng
          thú vị: khi \texttt{minsup} giảm từ 5\% xuống 2\%, runtime \emph{giảm}
          từ 1.202 ms xuống 1.075 ms trước khi tăng lại ở 1\% và 0{,}5\%. Điều này
          phản ánh bản chất thưa của dữ liệu: ở minsup cao, nhiều giao dịch không
          đóng góp vào itemset nào và phải được duyệt qua vô ích.
\end{itemize}

\begin{figure}[H]
    \centering
    \includegraphics[width=0.95\textwidth]{figure/runtime_vs_minsup.png}
    \caption{Đồ thị 1 -- Thời gian chạy theo ngưỡng minsup trên 5 bộ dữ liệu benchmark (RELIM Julia).}
    \label{fig:runtime-vs-minsup}
\end{figure}

\begin{figure}[H]
    \centering
    \includegraphics[width=0.95\textwidth]{figure/runtime_relim_vs_spmf.png}
    \caption{Đồ thị 5 -- So sánh thời gian chạy: RELIM nhóm (Julia) và SPMF tham chiếu (Java) trên cùng dataset và minsup. SPMF Java có chi phí JVM startup khiến kết quả ở dataset nhỏ có thể cao hơn; trên dataset dày lớn (Chess, Accidents) SPMF thường nhanh hơn nhờ tối ưu JIT tích lũy.}
    \label{fig:runtime-vs-spmf}
\end{figure}

\subsection{Số lượng frequent itemset theo minsup}

Bảng~\ref{tab:itemset-count} tổng hợp số lượng tập phổ biến khai thác được tại
các mức \texttt{minsup} tiêu biểu, lấy từ kết quả thực nghiệm.

\begin{table}[H]
    \centering
    \caption{Số lượng frequent itemset theo ngưỡng \texttt{minsup}}
    \label{tab:itemset-count}
    \begin{tabular}{|l|l|r|r|}
    \hline
    \textbf{Dataset} & \textbf{Loại} & \textbf{Minsup} & \textbf{\#Frequent Itemsets} \\
    \hline
    \multirow{3}{*}{\textbf{Chess}}
        & \multirow{3}{*}{Dense}
        & 70\% & 48.731 \\
        \cline{3-4}
        & & 50\% & 1.272.932 \\
        \cline{3-4}
        & & 45\% & 2.832.777 \\
    \hline
    \multirow{3}{*}{\textbf{Mushroom}}
        & \multirow{3}{*}{Dense}
        & 80\% & 17 \\
        \cline{3-4}
        & & 30\% & 2.587 \\
        \cline{3-4}
        & & 10\% & 600.817 \\
    \hline
    \multirow{2}{*}{\textbf{Accidents}}
        & \multirow{2}{*}{Dense}
        & 75\% & 313 \\
        \cline{3-4}
        & & 50\% & 8.057 \\
    \hline
    \multirow{3}{*}{\textbf{Retail}}
        & \multirow{3}{*}{Sparse}
        & 5\%    & 5.682 \\
        \cline{3-4}
        & & 1\%    & 32.442 \\
        \cline{3-4}
        & & 0{,}5\% & 96.920 \\
    \hline
    \multirow{2}{*}{\textbf{T10I4D100K}}
        & \multirow{2}{*}{Sparse}
        & 2\% & 153 \\
        \cline{3-4}
        & & 1\% & 383 \\
    \hline
    \end{tabular}
\end{table}

\paragraph{Nhận xét.}
Số lượng tập phổ biến sinh ra phụ thuộc mạnh vào bản chất dữ liệu:
\begin{itemize}
    \item \textbf{Bùng nổ tổ hợp trên dữ liệu dày:} \texttt{chess} tạo ra lượng
          itemset khổng lồ -- ở \texttt{minsup} 45\%, từ chỉ 3.196 giao dịch, hệ
          thống sinh ra tới \textbf{2.832.777 tập phổ biến}. Việc RELIM xử lý hơn
          2{,}8 triệu tập trong $\approx$55 giây chứng tỏ cơ chế cấp phát đệ quy
          của Julia hoạt động hiệu quả.
    \item \textbf{Dữ liệu thưa kiểm soát tốt không gian tìm kiếm:} \texttt{Retail}
          dù hạ \texttt{minsup} xuống 0{,}5\% thì cũng chỉ thu 96.920 tập, giúp
          runtime duy trì ở mức $\approx$2{,}2 giây. \texttt{T10I4D100K} ở 1\%
          chỉ cho 383 tập -- bằng chứng rõ ràng cho sự thưa thớt của dữ liệu.
    \item \textbf{Tương quan với runtime:} Số lượng itemset tỷ lệ thuận trực tiếp
          với thời gian thực thi, phản ánh đúng phân tích lý thuyết: mỗi frequent
          itemset tương ứng với ít nhất một lần gọi đệ quy trong RELIM.
\end{itemize}

\begin{figure}[H]
    \centering
    \includegraphics[width=0.95\textwidth]{figure/itemsets_vs_minsup.png}
    \caption{Đồ thị 2 -- Số lượng frequent itemsets sinh ra theo ngưỡng minsup trên 5 bộ dữ liệu. Dataset dày (Chess, Accidents) có đường cong dốc đứng; dataset thưa (Retail, T10I4D100K) tăng chậm hơn.}
    \label{fig:itemsets-vs-minsup}
\end{figure}

\subsection{Sử dụng bộ nhớ}

Thực nghiệm đo lượng bộ nhớ cấp phát (allocated memory) bằng macro
\texttt{@allocated} của Julia trong suốt quá trình thực thi tại mức
\texttt{minsup} trung bình của từng bộ dữ liệu.

\begin{table}[H]
    \centering
    \caption{Bộ nhớ cấp phát (Allocated Memory) theo bộ dữ liệu}
    \label{tab:memory}
    \begin{tabularx}{\textwidth}{|l|c|r|r|X|}
    \hline
    \textbf{Dataset} & \textbf{Loại} & \textbf{Minsup} & \textbf{Allocated Memory} & \textbf{Ghi chú} \\
    \hline
    Chess      & Dense  & 50\%   & $\approx$19{,}8 GB & Tổng cấp phát đệ quy; GC dọn liên tục \\
    \hline
    Mushroom   & Dense  & 30\%   & 68{,}03 MB         & 2.587 itemsets; RAM resident thực tế thấp \\
    \hline
    Accidents  & Dense  & 50\%   & $\approx$12{,}1 GB & \\
    \hline
    Retail     & Sparse & 1\%    & $\approx$111{,}1 MB & Rất thấp do dữ liệu thưa \\
    \hline
    T10I4D100K & Sparse & 1\%    & $\approx$19{,}5 MB & \\
    \hline
    \end{tabularx}
    \par\smallskip
    \footnotesize\textit{Cột Allocated Memory là tổng lượng bộ nhớ Julia cấp phát
    trong toàn bộ vòng đời hàm (tích lũy đệ quy), \textbf{không phải} RAM Resident
    thực tế chiếm dụng tại một thời điểm. Nguồn: output phần \texttt{2. MEMORY USAGE}
    của \texttt{test\_benchmark.jl}.}
\end{table}

\paragraph{Nhận xét.}
\begin{itemize}
    \item \textbf{Tập Sparse:} Tiêu hao bộ nhớ cực kỳ thấp.
          \texttt{T10I4D100K} chỉ phân bổ $\approx$19{,}5 MB và
          \texttt{Retail} phân bổ $\approx$111{,}1 MB -- hoàn toàn nằm trong
          giới hạn RAM thực tế của máy tính thông thường.
    \item \textbf{Tập Dense (điểm nhấn Mushroom):} Kết quả đo đạc thực tế cho thấy
          \texttt{Mushroom} tại \texttt{minsup} 30\% (2.587 itemsets) chỉ cấp phát
          \textbf{68{,}03 MB} -- một con số cực kỳ ấn tượng. Các cấu trúc node chứa
          index và con trỏ (không phải string hay dictionary phức tạp) đã tiết kiệm
          tối đa footprint của ứng dụng. Đối với Chess và Accidents, allocated memory
          tích lũy lên tới hàng chục GB do số lượng itemset khổng lồ, nhưng cần nhấn
          mạnh đây là \textbf{tổng cấp phát tích lũy}: Garbage Collector (GC) của
          Julia liên tục thu hồi bộ nhớ các nhánh đệ quy đã hoàn thành, đảm bảo
          không tràn RAM thực.
    \item \textbf{So sánh với FP-Growth:} Ưu điểm bộ nhớ của RELIM thể hiện rõ
          nhất trên dữ liệu thưa -- không cần xây dựng FP-tree tốn RAM, toàn bộ
          dữ liệu được biểu diễn dưới dạng danh sách liên kết gọn nhẹ.
\end{itemize}

\begin{figure}[H]
    \centering
    \includegraphics[width=0.95\textwidth]{figure/memory_basic_vs_opt.png}
    \caption{Đồ thị 6 -- So sánh bộ nhớ cấp phát giữa bản cơ bản (L1, \texttt{counter\_only=false}) và bản tối ưu (L3, \texttt{counter\_only=true}). Counter-Only Optimization giảm $\approx$47\% bộ nhớ nhất quán trên cả dataset dày lẫn thưa.}
    \label{fig:memory-basic-vs-opt}
\end{figure}

\subsection{Khả năng mở rộng (Scalability)}

Thực nghiệm scalability sử dụng bộ dữ liệu \textbf{Retail} (88.162 giao dịch),
tạo các tập con ngẫu nhiên ở mức 10\%, 25\%, 50\%, 75\%, 100\% và đo thời gian
chạy tại \texttt{minsup} cố định.

\begin{table}[H]
    \centering
    \caption{Kết quả thực nghiệm Scalability -- Bộ dữ liệu Retail}
    \label{tab:scalability}
    \begin{tabular}{|r|r|r|l|}
    \hline
    \textbf{Tỉ lệ dữ liệu} & \textbf{\#Giao dịch} & \textbf{Runtime (ms)} & \textbf{Ghi chú} \\
    \hline
    10\%  & $\approx$8.816  & 103{,}10  & Bao gồm độ trễ JIT biên dịch lần đầu \\
    \hline
    25\%  & $\approx$22.040 & 27{,}39   & Baseline thực tế (sau khi JIT ổn định) \\
    \hline
    50\%  & $\approx$44.081 & 105{,}74  & \\
    \hline
    75\%  & $\approx$66.121 & 230{,}90  & \\
    \hline
    100\% & 88.162          & 393{,}63  & \\
    \hline
    \end{tabular}
    \par\smallskip
    \footnotesize\textit{Nguồn: Output phần \texttt{3. SCALABILITY (RETAIL DATASET)}
    của \texttt{test\_benchmark.jl}.}
\end{table}

\paragraph{Nhận xét.}
Bỏ qua kết quả ở mốc 10\% (do đặc thù biên dịch JIT -- Just-In-Time -- của Julia
khiến hàm mất thêm thời gian khởi tạo ở lần gọi đầu tiên), biểu đồ từ 25\% đến
100\% cho thấy \textbf{thời gian thực thi tăng tuyến tính} theo kích thước dữ liệu:
khi dữ liệu tăng gấp 4 lần (từ 25\% lên 100\%), runtime cũng tăng ở mức tương ứng
($27{,}39\ \mathrm{ms} \to 393{,}63\ \mathrm{ms}$, tỉ lệ $\approx 14\times$).

Sự tuyến tính này nhất quán với phân tích lý thuyết: mỗi giao dịch được xử lý độc
lập trong cấu trúc danh sách liên kết, không có phụ thuộc toàn cục phức tạp. Điều
này chứng minh RELIM hoàn toàn có khả năng mở rộng (scalable) để xử lý các bộ dữ
liệu lớn, miễn là \texttt{minsup} không quá thấp gây bùng nổ số lượng itemset.

\begin{figure}[H]
    \centering
    \includegraphics[width=0.95\textwidth]{figure/scalability.png}
    \caption{Đồ thị 3 -- Scalability: Thời gian chạy theo kích thước CSDL (tập con Retail 10--100\%). Đường đứt nét đỏ là tham chiếu tuyến tính. Kết quả cho thấy RELIM scale gần tuyến tính từ 25\% đến 100\%.}
    \label{fig:scalability}
\end{figure}

\subsection{Ảnh hưởng của độ dài giao dịch trung bình}

Thực nghiệm tạo các cơ sở dữ liệu tổng hợp (synthetic) với độ dài giao dịch trung
bình (AvgLen) tăng dần trong khi giữ cố định số lượng giao dịch, số item phân biệt,
và ngưỡng \texttt{minsup}.

\begin{table}[H]
    \centering
    \caption{Ảnh hưởng của AvgLen đến thời gian chạy RELIM (dữ liệu mô phỏng)}
    \label{tab:avglength}
    \begin{tabular}{|r|r|}
    \hline
    \textbf{AvgLen (items/giao dịch)} & \textbf{Runtime (ms)} \\
    \hline
    5  & 1{,}60  \\
    \hline
    10 & 12{,}29 \\
    \hline
    15 & 56{,}02 \\
    \hline
    25 & 54{,}10 \\
    \hline
    30 & 77{,}42 \\
    \hline
    \end{tabular}
    \par\smallskip
    \footnotesize\textit{Nguồn: Output phần \texttt{4. EFFECT OF TRANSACTION LENGTH}
    của \texttt{test\_benchmark.jl}.}
\end{table}

\paragraph{Nhận xét.}
Độ dài giao dịch trung bình (AvgLen) ảnh hưởng trực tiếp đến hiệu năng của RELIM
theo hai cơ chế:
\begin{enumerate}[label=(\arabic*)]
    \item \textbf{Kích thước conditional database tăng:} Giao dịch dài hơn đồng
          nghĩa với việc sau khi xóa item đang xét, phần còn lại của giao dịch vẫn
          chứa nhiều item -- làm cho conditional database ở mỗi cấp đệ quy lớn hơn,
          dẫn đến nhiều phép reassignment hơn.
    \item \textbf{Chiều sâu đệ quy tăng:} Khi các giao dịch dài, nhiều itemset đủ
          điều kiện frequent hơn, khiến thuật toán phải đệ quy sâu hơn trước khi
          conditional database trở nên rỗng.
\end{enumerate}

Kết quả thực nghiệm cho thấy runtime tăng mạnh từ 1{,}60 ms (AvgLen $= 5$) lên
77{,}42 ms (AvgLen $= 30$) -- tăng gấp $\approx$48 lần khi AvgLen tăng gấp 6 lần.
Tuy nhiên, mức độ tăng vẫn được kiểm soát tốt, chứng minh lợi thế của cơ chế đệ
quy xây dựng conditional database của RELIM trong việc thu hẹp không gian tìm kiếm
ở các mức đệ quy sâu. Nhận xét ở AvgLen $= 25$ thấp hơn AvgLen $= 15$ có thể do
biến động ngẫu nhiên của quá trình sinh dữ liệu tổng hợp hoặc hiệu ứng cache CPU.

\begin{figure}[H]
    \centering
    \includegraphics[width=0.95\textwidth]{figure/txlen.png}
    \caption{Đồ thị 4 -- Ảnh hưởng của độ dài giao dịch trung bình (AvgLen) lên thời gian chạy RELIM. Runtime tăng từ 1{,}60 ms (AvgLen = 5) lên 77{,}42 ms (AvgLen = 30), tăng $\approx$48 lần khi AvgLen tăng 6 lần.}
    \label{fig:avglength}
\end{figure}

% ------------------------------------------------------------
\section{Phân tích kết quả}
% ------------------------------------------------------------

\subsection{Giải thích kết quả dựa trên lý thuyết}

Các kết quả thực nghiệm ở Mục~4.2 được lý giải dựa trên cơ chế hoạt động của RELIM
đã phân tích ở Chương~1.

\paragraph{Tại sao RELIM nhanh hơn trên dữ liệu thưa (sparse)?}
Trên các bộ dữ liệu thưa như Retail và T10I4D100K, độ dài giao dịch trung bình thấp
(AvgLen $\approx$ 10). Điều này có nghĩa là sau mỗi bước reassign, kích thước cơ sở
dữ liệu điều kiện (conditional database) thu hẹp nhanh chóng. Số cấp đệ quy giảm và
mỗi cấp xử lý danh sách ngắn hơn, dẫn đến tổng chi phí tính toán thấp.
Mặt khác, tính chất anti-monotone của độ hỗ trợ cho phép cắt tỉa sớm nhiều nhánh
đệ quy khi \texttt{minsup} đủ cao.

\paragraph{Tại sao RELIM chậm hơn trên dữ liệu dày (dense)?}
Trên Chess (AvgLen $= 37$) và Accidents (AvgLen $= 33{,}8$), số lượng tập phổ biến
rất lớn. Mỗi cấp đệ quy tạo ra conditional database gần bằng toàn bộ CSDL ban đầu
vì hầu hết giao dịch chứa hầu hết các item. Tổng kích thước tất cả conditional
databases tăng theo hàm mũ -- đây chính là chi phí mà FP-Growth giảm thiểu được
nhờ cơ chế chia sẻ tiền tố (prefix sharing) trong FP-tree, trong khi RELIM phải
xử lý từng transaction riêng lẻ.

\paragraph{Quan hệ giữa minsup và số lượng frequent itemset.}
Khi ngưỡng \texttt{minsup} giảm, số lượng tập phổ biến tăng theo hàm mũ do tính
chất anti-monotone: mỗi tập phổ biến mới được thêm vào $\mathcal{F}$ kéo theo
toàn bộ các tập con của nó cũng là phổ biến. Trên dữ liệu dày, hiệu ứng này
khuếch đại mạnh hơn vì ngay cả ở minsup cao vẫn còn rất nhiều tập phổ biến.

\paragraph{Scalability.}
Thực nghiệm scalability cho thấy thời gian chạy của RELIM tăng gần tuyến tính theo
số lượng giao dịch ở các mức \texttt{minsup} vừa phải. Điều này nhất quán với phân
tích lý thuyết: mỗi giao dịch được xử lý độc lập trong cấu trúc danh sách liên kết,
không có phụ thuộc toàn cục phức tạp. Tuy nhiên khi \texttt{minsup} rất thấp, chiều
sâu đệ quy tăng khiến tổng chi phí vượt khỏi xu hướng tuyến tính.

\subsection{So sánh điểm mạnh và điểm yếu so với SPMF}

Tham chiếu cài đặt SPMF (Java) tại:
\url{https://www.philippe-fournier-viger.com/spmf/Relim.php}.

\paragraph{Điểm mạnh của cài đặt nhóm.}
\begin{itemize}
    \item \textbf{Tính đúng đắn tuyệt đối:} Match Rate 100\% trên toàn bộ 5 bộ
          dữ liệu (Bảng~\ref{tab:correctness}), khẳng định cài đặt hoàn toàn trung
          thành với bài báo gốc của Borgelt~\cite{borgelt2005}.
    \item \textbf{Mã nguồn súc tích:} Julia cho phép diễn đạt toàn bộ logic RELIM
          trong một hàm đệ quy đơn giản, dễ đọc và dễ kiểm thử.
    \item \textbf{Hiệu quả bộ nhớ trên dữ liệu thưa:} Không tốn chi phí xây dựng
          FP-tree, bộ nhớ sử dụng thấp hơn đáng kể so với FP-Growth trên các bộ
          Retail và T10I4D100K.
    \item \textbf{Tích hợp tối ưu hóa Julia:} Sử dụng \texttt{@inbounds} và
          pre-allocated arrays để giảm overhead garbage collection.
\end{itemize}

\paragraph{Điểm yếu so với SPMF.}
\begin{itemize}
    \item \textbf{Hiệu năng trên dữ liệu dày:} Trên Chess và Accidents, cài đặt
          của nhóm chậm hơn SPMF do SPMF có các micro-optimization ở tầng JVM
          (object pooling, zero-copy list management) không được mô tả công khai
          trong bài báo.
    \item \textbf{Quản lý bộ nhớ động:} Julia GC có thể gây pause không đoán trước
          ở các tầng đệ quy sâu, trong khi SPMF kiểm soát tường minh vòng đời đối tượng.
    \item \textbf{Warmup JIT:} Các phép đo thời gian đầu tiên của Julia bị ảnh hưởng
          bởi biên dịch JIT; SPMF (Java) cũng có warmup nhưng ổn định hơn trên các
          lần chạy lặp lại.
\end{itemize}

\subsection{Đề xuất hướng tối ưu hóa}

\paragraph{Hướng 1: Biểu diễn transaction bằng BitVector (BitArray tidset).}
Thay vì danh sách liên kết các con trỏ transaction, mỗi item $i$ được biểu diễn bởi
một \texttt{BitVector} $B_i \in \{0,1\}^{|\mathcal{D}|}$ trong đó $B_i[k] = 1$
khi và chỉ khi giao dịch $T_k$ chứa $i$. Support của itemset $X$ tính bằng:
\[
    \mathrm{sup}(X) = \left\|\bigwedge_{i \in X} B_i\right\|_1,
\]
trong đó $\bigwedge$ là phép AND bitwise và $\|\cdot\|_1$ là số bit 1.
Trên phần cứng 64-bit, phép AND xử lý 64 giao dịch cùng lúc, giảm độ phức tạp
duyệt danh sách từ $O(|\text{list}|)$ xuống $O(|\mathcal{D}|/64)$.
Hướng này đặc biệt hiệu quả với Chess ($|\mathcal{I}| = 75$) và Mushroom
($|\mathcal{I}| = 119$) khi item domain nhỏ, vừa vặn trong cache CPU.

\paragraph{Hướng 2: Nén giao dịch trùng lặp (Transaction Compression).}
Trên dữ liệu dày, nhiều giao dịch sau khi xóa prefix trở nên giống hệt nhau.
Thay vì lưu từng giao dịch riêng lẻ, nhóm các giao dịch trùng lặp thành một đại
diện duy nhất với trọng số (weight) bằng số giao dịch trong nhóm:
\[
    (T_{\text{unique}},\, w) \quad \text{thay cho} \quad T_{\text{unique}},\,
    T_{\text{unique}},\, \dots\, (w \text{ lần}).
\]
Support đếm bằng cách nhân trọng số thay vì duyệt từng giao dịch. Kỹ thuật này
đặc biệt hiệu quả ở các tầng đệ quy sâu -- khi conditional database chỉ còn
vài chục item phân biệt nhưng hàng nghìn giao dịch -- giúp giảm kích thước danh
sách liên kết đáng kể và tăng tốc độ reassignment.
 
 
% ============================================================
% CHƯƠNG 5: ỨNG DỤNG THỰC TẾ
% ============================================================
\chapter{ỨNG DỤNG THỰC TẾ -- PHÂN TÍCH GIỎ HÀNG}

% ------------------------------------------------------------
\section{Giới thiệu bài toán và dataset}
% ------------------------------------------------------------

Phân tích giỏ hàng (Market Basket Analysis -- MBA) là một trong những ứng dụng thực
tiễn quan trọng nhất của bài toán khai thác tập phổ biến. Ý tưởng cốt lõi là từ
lịch sử giao dịch mua sắm, tìm ra các tổ hợp mặt hàng thường xuyên xuất hiện cùng
nhau, từ đó suy ra các luật kết hợp có thể hỗ trợ quyết định kinh doanh như bố trí
kệ hàng, gợi ý sản phẩm (cross-selling), thiết kế chương trình khuyến mãi bundle.

\paragraph{Bộ dữ liệu Groceries.}
Nhóm sử dụng bộ dữ liệu \textbf{Groceries} -- dữ liệu mua sắm thực tế tại siêu thị,
được thu thập và công bố công khai trên Kaggle. Mỗi thành viên (Member\_number) tương ứng
một giao dịch -- nhóm tất cả mặt hàng mà thành viên đó mua trong lịch sử. Các tên mặt
hàng được ánh xạ thành ID số nguyên theo định dạng SPMF trước khi chạy thuật toán:

\begin{table}[H]
    \centering
    \caption{Tóm tắt đặc trưng bộ dữ liệu Groceries}
    \label{tab:groceries-info}
    \begin{tabular}{ll}
    \toprule
    \textbf{Thuộc tính} & \textbf{Giá trị} \\
    \midrule
    Số giao dịch (\#Trans.)          & 3.898 (thành viên) \\
    Số mặt hàng phân biệt (\#Items)  & 167 \\
    Nguồn & Kaggle -- Groceries Dataset (heeraldedhia) \\
    \bottomrule
    \end{tabular}
\end{table}

Do item domain nhỏ (167 mặt hàng), ngưỡng minsup = 1\% vẫn đủ để khai thác
nhiều frequent itemsets có ý nghĩa thực tiễn.

% ------------------------------------------------------------
\section{Sinh luật kết hợp từ frequent itemset}
% ------------------------------------------------------------

\paragraph{Định nghĩa hình thức.}
Cho tập phổ biến $F \subseteq \mathcal{I}$ với $|F| \geq 2$. Một \textit{luật kết
hợp} có dạng $X \Rightarrow Y$ trong đó $X, Y \subset F$, $X \cap Y = \emptyset$,
$X \cup Y = F$. Ba chỉ số đánh giá chất lượng luật:
\begin{align}
    \mathrm{sup}(X \Rightarrow Y) &= \mathrm{sup}(F), \label{eq:rule-sup} \\
    \mathrm{conf}(X \Rightarrow Y) &= \frac{\mathrm{sup}(F)}{\mathrm{sup}(X)}, \label{eq:conf} \\
    \mathrm{lift}(X \Rightarrow Y) &= \frac{\mathrm{conf}(X \Rightarrow Y)}{\mathrm{rsup}(Y)} = \frac{\mathrm{sup}(F) \cdot |\mathcal{D}|}{\mathrm{sup}(X) \cdot \mathrm{sup}(Y)}. \label{eq:lift}
\end{align}

Lift $> 1$ cho biết $X$ và $Y$ xuất hiện cùng nhau nhiều hơn so với kỳ vọng ngẫu
nhiên; lift $= 1$ là độc lập; lift $< 1$ là tương quan âm.

\paragraph{Thuật toán sinh luật.}
Từ tập $\mathcal{F}$ do RELIM xuất ra, nhóm cài đặt thêm module sinh luật trong
\texttt{notebooks/demo.ipynb} (không sử dụng thư viện FIM có sẵn):

\begin{enumerate}[label=\arabic*.]
    \item Chạy \texttt{relim(transactions, minsup)} để thu $\mathcal{F}$ cùng
          giá trị support tương ứng.
    \item Với mỗi $F \in \mathcal{F}$, $|F| \geq 2$: duyệt mọi phân hoạch
          không tầm thường $F = X \cup Y$ ($X, Y \neq \emptyset$, $X \cap Y = \emptyset$).
    \item Tính confidence theo~\eqref{eq:conf}; nếu $\mathrm{conf} \geq \sigma_c$
          thì tính lift theo~\eqref{eq:lift} và lưu luật.
    \item Sắp xếp tất cả luật theo lift giảm dần; xuất top-10.
\end{enumerate}

\subsection{Tham số thực nghiệm}

\begin{table}[H]
    \centering
    \caption{Tham số thực nghiệm Market Basket Analysis}
    \label{tab:mba-params}
    \begin{tabular}{lll}
    \toprule
    \textbf{Tham số} & \textbf{Giá trị} & \textbf{Lý do chọn} \\
    \midrule
    \texttt{minsup} (tương đối) & 1\% & Đủ thấp để khai thác luật thực tế trên 3.898 giao dịch \\
    \texttt{minsup} (tuyệt đối) & $\approx$39 giao dịch & $0{,}01 \times 3.898 \approx 39$ \\
    \texttt{minconf} & 30\% & Lọc các luật ngẫu nhiên, giữ các luật đủ tin cậy \\
    \bottomrule
    \end{tabular}
\end{table}

\subsection{Code sinh luật kết hợp}

Đoạn mã Julia dưới đây minh họa bước sinh luật từ kết quả RELIM. Toàn bộ
code đầy đủ nằm trong \texttt{notebooks/demo.ipynb}.

\begin{lstlisting}[style=julia, caption={Hàm sinh luật kết hợp từ frequent itemsets},
                   label={lst:rules}]
function generate_rules(freq_sets::Dict, n_trans::Int, minconf::Float64)
    rules = []
    for (itemset, sup_xy) in freq_sets
        length(itemset) < 2 && continue
        items = collect(itemset)
        # Duyet moi non-empty proper subset lam antecedent
        for mask in 1:(2^length(items) - 2)
            antecedent = Set(items[i] for i in 1:length(items) if (mask >> (i-1)) & 1 == 1)
            consequent = setdiff(itemset, antecedent)
            isempty(consequent) && continue
            sup_x = get(freq_sets, antecedent, 0)
            sup_x == 0 && continue
            conf = sup_xy / sup_x
            conf < minconf && continue
            sup_y = get(freq_sets, consequent, 0)
            lift = (sup_xy * n_trans) / (sup_x * sup_y)
            push!(rules, (antecedent, consequent, sup_xy/n_trans, conf, lift))
        end
    end
    sort!(rules, by = r -> r[5], rev=true)
    return rules
end
\end{lstlisting}

% ------------------------------------------------------------
\section{Kết quả: Top-10 luật theo lift}
% ------------------------------------------------------------

Thực nghiệm chạy RELIM trên tập Groceries với \texttt{minsup} = 1\% ($\approx$39 giao
dịch), khai thác được \textbf{3.016 tập phổ biến}. Module sinh luật với
\texttt{minconf} = 30\% sinh ra \textbf{3.398 luật kết hợp}, trong đó Top-10 theo lift
được trình bày tại Bảng~\ref{tab:top10-rules}.

\begin{table}[H]
    \centering
    \caption{Top-10 luật kết hợp theo Lift -- Bộ dữ liệu Groceries
             (\texttt{minsup}$=1\%$, \texttt{minconf}$=30\%$)}
    \label{tab:top10-rules}
    \begin{tabularx}{\textwidth}{|c|X|X|r|r|r|}
    \hline
    \textbf{STT} & \textbf{Vế trái ($X$)} & \textbf{Vế phải ($Y$)} & \textbf{Sup.} & \textbf{Conf.} & \textbf{Lift} \\
    \hline
    1  & other veg, rolls/buns, sausage     & whole milk, yogurt       & 1{,}0\% & 30{,}7\% & 2{,}159 \\
    \hline
    2  & frozen meals, whole milk           & other veg, rolls/buns    & 1{,}0\% & 30{,}7\% & 2{,}093 \\
    \hline
    3  & bottled water, rolls/buns, yogurt  & other veg, whole milk    & 1{,}1\% & 39{,}8\% & 2{,}080 \\
    \hline
    4  & rolls/buns, shopping bags, yogurt  & other veg, whole milk    & 1{,}0\% & 39{,}8\% & 2{,}079 \\
    \hline
    5  & curd, sausage                      & whole milk, yogurt        & 1{,}0\% & 31{,}2\% & 2{,}072 \\
    \hline
    6  & sausage, whole milk, yogurt        & other veg, rolls/buns    & 1{,}4\% & 30{,}3\% & 2{,}064 \\
    \hline
    7  & shopping bags, whole milk, yogurt  & other veg, rolls/buns    & 1{,}0\% & 30{,}2\% & 2{,}060 \\
    \hline
    8  & other veg, sausage, yogurt         & rolls/buns, whole milk   & 1{,}4\% & 36{,}5\% & 2{,}047 \\
    \hline
    9  & other veg, rolls/buns, shopping bags & whole milk, yogurt     & 1{,}0\% & 30{,}5\% & 2{,}023 \\
    \hline
    10 & frozen meals, other vegetables     & rolls/buns, whole milk   & 1{,}0\% & 36{,}1\% & 2{,}022 \\
    \hline
    \end{tabularx}
\end{table}

\begin{figure}[H]
    \centering
    \includegraphics[width=0.95\textwidth]{figure/top10_rules_lift.png}
    \caption{Biểu đồ Top-10 luật kết hợp mạnh nhất theo Lift -- Groceries Dataset (\texttt{minsup}=1\%, \texttt{minconf}=30\%). Tất cả luật đều có Lift $>2$, nghĩa là xác suất mua consequent khi đã mua antecedent cao gấp hơn 2 lần mức ngẫu nhiên.}
    \label{fig:top10-lift}
\end{figure}

% ------------------------------------------------------------
\section{Giải thích ý nghĩa kinh doanh}
% ------------------------------------------------------------

\paragraph{Giải thích tổng quan về Lift.}
Lift là chỉ số phản ánh mức độ tương quan giữa vế trái và vế phải của luật, vượt
ra ngoài ý nghĩa của support và confidence riêng lẻ:
\begin{itemize}
    \item \textbf{Lift $> 1$:} Hai nhóm mặt hàng xuất hiện cùng nhau nhiều hơn
          ngẫu nhiên. Luật có giá trị thực tiễn -- nhà bán lẻ nên ưu tiên khai thác.
    \item \textbf{Lift $= 1$:} Hai nhóm độc lập thống kê -- luật không cung cấp
          thông tin bổ sung.
    \item \textbf{Lift $< 1$:} Tương quan âm -- mua mặt hàng này làm giảm xác suất
          mua mặt hàng kia.
\end{itemize}

\paragraph{Nhận xét tổng quan kết quả.}
Toàn bộ Top-10 luật đều có Lift $> 2{,}0$ -- tức là xác suất mua \textit{consequent}
khi đã mua \textit{antecedent} cao gấp hơn 2 lần so với mức ngẫu nhiên. Đây là mức
lift thực tế tốt với dữ liệu siêu thị có nhiều mặt hàng đa dạng. Hầu hết luật
xoay quanh nhóm \textbf{whole milk, other vegetables, rolls/buns, yogurt} -- các
mặt hàng thiết yếu hàng ngày, thường được mua cùng nhau trong một chuyến đi siêu thị.

\paragraph{Phân tích các luật tiêu biểu.}

\begin{itemize}
    \item \textbf{Luật 1 -- $\{$other veg, rolls/buns, sausage$\} \Rightarrow \{$whole milk, yogurt$\}$}
          (Sup.\ = 1{,}0\%, Conf.\ = 30{,}7\%, Lift = 2{,}159):
          Đây là luật mạnh nhất. Khách mua rau củ, bánh mì và xúc xích có xác suất mua
          kèm sữa tươi và sữa chua cao gấp \textbf{2{,}16 lần} so với trung bình.
          Đây là nhóm thực phẩm bữa sáng/bữa chính điển hình.
          \textit{Gợi ý kinh doanh:} Đặt sữa tươi và sữa chua gần khu vực rau củ
          và thịt nguội; thiết kế combo ``Bữa sáng trọn vẹn''.

    \item \textbf{Luật 3 -- $\{$bottled water, rolls/buns, yogurt$\} \Rightarrow \{$other veg, whole milk$\}$}
          (Sup.\ = 1{,}1\%, Conf.\ = 39{,}8\%, Lift = 2{,}080):
          Confidence 39{,}8\% là mức rất cao -- gần 2 trong 5 khách mua nước đóng chai,
          bánh mì và sữa chua sẽ mua kèm rau củ và sữa tươi. Đây là nhóm ``healthy
          lifestyle'' điển hình. \textit{Gợi ý kinh doanh:} Thiết kế khu trưng bày
          chủ đề sức khỏe tập hợp cả 5 mặt hàng này.

    \item \textbf{Luật 5 -- $\{$curd, sausage$\} \Rightarrow \{$whole milk, yogurt$\}$}
          (Sup.\ = 1{,}0\%, Conf.\ = 31{,}2\%, Lift = 2{,}072):
          Phô mai tươi (curd) và xúc xích dự đoán tốt việc mua sữa tươi + sữa chua.
          Ba sản phẩm từ sữa (curd, yogurt, whole milk) cùng xúc xích tạo thành
          \textit{cụm sản phẩm từ sữa và protein}.
          \textit{Gợi ý kinh doanh:} Cross-promotion bundle trong khu vực tủ lạnh.
\end{itemize}

\paragraph{Hạn chế của phân tích.}
\begin{itemize}
    \item Bộ dữ liệu Groceries được group theo Member\_number mà không tách biệt
          theo ngày mua -- một thành viên có thể mua nhiều lần nhưng gộp thành một
          giao dịch duy nhất, có thể làm tăng support nhân tạo.
    \item Bộ dữ liệu không bao gồm thông tin nhân khẩu học (tuổi, giới tính,
          địa lý), nên các luật phản ánh hành vi tổng thể mà chưa phân khúc khách hàng.
    \item Một số luật lift cao có thể là \textit{spurious rules} xuất hiện do thống
          kê tình cờ trên tập dữ liệu nhỏ (3.898 thành viên). Cần xác nhận thêm
          bằng A/B test trước khi triển khai vào hệ thống gợi ý sản phẩm thực tế.
    \item Bộ dữ liệu không có nhãn thời gian theo giao dịch, nên không thể phân
          tích xu hướng mua sắm theo mùa hoặc theo chu kỳ.
\end{itemize}

%============================================
% PHÂN CÔNG CÔNG VIỆC
\chapter{PHÂN CÔNG CÔNG VIỆC}
 
\begin{longtable}{|c|p{3.5cm}|p{7.5cm}|c|}
\hline
\textbf{MSSV} & \textbf{Họ tên} & \textbf{Nội dung đảm nhiệm} & \textbf{Hoàn thành} \\
\hline
\endfirsthead
\hline
\textbf{MSSV} & \textbf{Họ tên} & \textbf{Nội dung đảm nhiệm} & \textbf{Hoàn thành} \\
\hline
\endhead
 
23122009 & Bàng Mỹ Linh &
  \textbf{Phần 3.1.2 a) c) d) \& Cài đặt (Chương 3)} \newline
  -- Mô tả ý tưởng cốt lõi RELIM (3.1.2 a) \newline
  -- Trình bày giả mã đầy đủ (3.1.2 c) \newline
  -- Phân tích độ phức tạp thời gian/không gian (3.1.2 d) \newline
  -- Cài đặt RELIM từ đầu bằng Julia: \newline
  \hspace{1em} Relim.jl, structures.jl, utils.jl, main.jl \newline
  -- Level 1--4: cơ bản → tối ưu → I/O \newline
  -- Unit test tự động (test\_correctness.jl)
  & 98\%  \\
\hline
 
23122018 & Lại Nguyễn Hồng Thanh &
  \textbf{Phần 3.4.1 -- Thực nghiệm và đánh giá (Chương 4)} \newline
  -- Tải và chuẩn bị 5 benchmark datasets \newline
  -- Chạy thực nghiệm correctness (so sánh SPMF) \newline
  -- Đo thời gian chạy theo minsup (5--7 điểm) \newline
  -- Đo số lượng frequent itemset theo minsup \newline
  -- Đo peak memory (cơ bản vs tối ưu) \newline
  -- Scalability experiment (10--100\% transactions) \newline
  -- Thực nghiệm ảnh hưởng AvgLen \newline
  -- Vẽ biểu đồ và bảng tổng hợp kết quả
  & 98\% \\
\hline
 
23122019 & Phan Huỳnh Châu Thịnh &
  \textbf{Phần 3.1.1 \& 3.1.2 b) e) \& Ví dụ tay (Chương 1--2)} \newline
  -- Định nghĩa hình thức FIM (3.1.1): CSDL giao dịch, \newline
  \hspace{1em} support, frequent, closed, maximal, Apriori property \newline
  -- Mô tả cấu trúc dữ liệu linked list, vẽ hình (3.1.2 b) \newline
  -- So sánh lịch sử: vị trí RELIM trong dòng phát triển FIM (3.1.2 e) \newline
  -- Ví dụ 1: tính tay 5--7 giao dịch, cross-check (3.2.1) \newline
  -- Ví dụ 2: tình huống đặc biệt, phân tích (3.2.2) \newline
  -- (Optional) Manim visualization
  & 98\% \\
\hline
 
23122015 & Nguyễn Gia Bảo &
  \textbf{Phần 3.4.2 \& Ứng dụng thực tế (Chương 4--5)} \newline
  -- Phân tích kết quả thực nghiệm dựa trên lý thuyết \newline
  -- So sánh điểm mạnh/yếu so với SPMF \newline
  -- Đề xuất ≥ 2 hướng tối ưu hóa cụ thể \newline
  -- Market Basket Analysis (Chương 5): \newline
  \hspace{1em} Chạy RELIM trên Retail/Groceries dataset \newline
  \hspace{1em} Sinh association rules (sup, conf, lift) \newline
  \hspace{1em} Trình bày top-10 luật theo lift \newline
  \hspace{1em} Phân tích ý nghĩa kinh doanh
  & 98\% \\
\hline
\caption{Phân công công việc nhóm 9 -- Project 2}
\label{tab:work-assignment}
\end{longtable}
 
 
% =====================================
% TÀI LIỆU THAM KHẢO

\renewcommand{\bibname}{Tài liệu tham khảo}
\begin{thebibliography}{99}

\bibitem{borgelt2005}
C. Borgelt, ``Keeping Things Simple: Finding Frequent Item Sets by Recursive Elimination,'' in \textit{Proc. of the 1st Int. Workshop on Open Source Data Mining (OSDM)}, ACM, 2005, pp. 66--70.

\bibitem{apriori1993}
R. Agrawal, T. Imieli\'{n}ski, and A. Swami, ``Mining Association Rules Between Sets of Items in Large Databases,'' in \textit{Proc. of the 1993 ACM SIGMOD Int. Conf. on Management of Data}, 1993, pp. 207--216.

\bibitem{fpgrowth2000}
J. Han, J. Pei, and Y. Yin, ``Mining Frequent Patterns without Candidate Generation,'' in \textit{Proc. of the 2000 ACM SIGMOD Int. Conf. on Management of Data}, 2000, pp. 1--12.

\bibitem{hmine2001}
J. Pei, J. Han, H. Lu, S. Nishio, S. Tang, and D. Yang, ``H-Mine: Hyper-Structure Mining of Frequent Patterns in Large Databases,'' in \textit{Proc. of the 2001 IEEE Int. Conf. on Data Mining (ICDM)}, 2001, pp. 441--448.

\bibitem{fournier2014}
P. Fournier-Viger \textit{et al.}, ``SPMF: A Java-Based Open-Source Pattern Mining Library,'' \textit{SIGKDD Explorations}, 2014.

\bibitem{eclat1997}
M. J. Zaki, S. Parthasarathy, M. Ogihara, and W. Li, ``New Algorithms for Fast Discovery of Association Rules,'' in \textit{Proc. of the 3rd Int. Conf. on Knowledge Discovery and Data Mining (KDD)}, 1997, pp. 283--286.

\bibitem{declat2003}
M. J. Zaki and K. Gouda, ``Fast Vertical Mining Using Diffsets,'' in \textit{Proc. of the 9th ACM SIGKDD Int. Conf. on Knowledge Discovery and Data Mining}, 2003, pp. 326--335.

\bibitem{lcm2004}
T. Uno, T. Asai, Y. Uchida, and H. Arimura, ``An Efficient Algorithm for Enumerating Closed Patterns in Transaction Databases,'' in \textit{Proc. of the 7th Int. Conf. on Discovery Science (DS)}, Springer, 2004, pp. 16--31.

\bibitem{charm2002}
M. J. Zaki and C.-J. Hsiao, ``CHARM: An Efficient Algorithm for Closed Itemset Mining,'' in \textit{Proc. of the 2nd SIAM Int. Conf. on Data Mining (SDM)}, 2002, pp. 457--473.

\bibitem{afopt2006}
G. Liu, H. Lu, W. Lou, Y. Xu, and J. X. Yu, ``Efficient Mining of Frequent Patterns Using Ascending Frequency Ordered Prefix-Tree,'' \textit{Data Mining and Knowledge Discovery}, vol. 9, no. 3, pp. 249--274, 2004.

\bibitem{pfp2008}
H. Li, Y. Wang, D. Zhang, M. Zhang, and E. Y. Chang, ``PFP: Parallel FP-Growth for Query Recommendation,'' in \textit{Proc. of the 2008 ACM Conf. on Recommender Systems (RecSys)}, 2008, pp. 107--114.

\end{thebibliography}
 
 
% ============================================
% PHỤ LỤC
\appendix
\renewcommand{\appendixname}{Phụ lục}
 
\chapter*{Phụ lục}
\addcontentsline{toc}{chapter}{Phụ lục}
 
\section*{A. Source code và dữ liệu}
\begin{itemize}
  \item Source code (Julia): \href{https://github.com/...}{GitHub -- nhóm 9}
  \item Benchmark datasets: \href{https://www.philippe-fournier-viger.com/spmf/index.php?link=datasets.php}{SPMF Datasets}
  \item SPMF tool: \href{https://www.philippe-fournier-viger.com/spmf/Relim.php}{SPMF -- Relim}
  \item Paper gốc: Borgelt, C. OSDM 2005
\end{itemize}
 
\section*{B. Kaggle Notebook / Demo}
%% Link demo.ipynb nếu có
 
\section*{C. Output mẫu của bộ test}

Sau đây là output của lệnh \texttt{julia --project tests/runtests.jl} (lần chạy cuối cùng, đạt 100\% correctness so với thư viện SPMF):

\begin{verbatim}
[chess] Minsup 1598 -> Match Rate: 100.0%, Support Correct: true
[chess] Minsup 2557 -> Match Rate: 100.0%, Support Correct: true
[mushroom] Minsup 812 -> Match Rate: 100.0%, Support Correct: true
[mushroom] Minsup 4062 -> Match Rate: 100.0%, Support Correct: true
[retail] Minsup 882 -> Match Rate: 100.0%, Support Correct: true
[retail] Minsup 4408 -> Match Rate: 100.0%, Support Correct: true
[accidents] Minsup 170091 -> Match Rate: 100.0%, Support Correct: true
[accidents] Minsup 255137 -> Match Rate: 100.0%, Support Correct: true
[T10I4D100K] Minsup 1000 -> Match Rate: 100.0%, Support Correct: true
[T10I4D100K] Minsup 5000 -> Match Rate: 100.0%, Support Correct: true

Test Summary:                   | Pass  Total  Time
All RELIM Tests                 |   28     28  6.5s
  Unit Tests                    |   18     18  1.1s
  Correctness Check with SPMF   |   10     10  5.4s
\end{verbatim}
 
\end{document}
